\section{Flat spectral geometry on the positive radial cusp}
\label{cusp-spec-section}

The period family used in this section is the family already constructed and
proved nondegenerate in \cref{prop:period-input-audit}.  In particular, this
section introduces no additional period hypothesis.  We retain the proved
ordered lattices
\[
 V=\mathbb Z\langle\gamma,u,w,\delta\rangle,
 \qquad
 \Lambda=V^*
 =\mathbb Z\langle\widehat\gamma,\widehat u,
                         \widehat w,\widehat\delta\rangle
\]
and the proved period matrix
\[
 \Pi=[Z\mid I_2],
 \qquad
 Z=\begin{pmatrix}6\mu&\tau\\ \beta&\mu\end{pmatrix}.
\]
Thus the columns of \(\Pi\), in the displayed ordered basis of \(\Lambda\),
are
\[
 (6\mu,\beta)^{\mathsf t},\quad
 (\tau,\mu)^{\mathsf t},\quad
 (1,0)^{\mathsf t},\quad
 (0,1)^{\mathsf t}.
\]

The canonical cusp calculation has already produced holomorphic functions
\(\mu(t_c),b(t_c),h(t_c)\) on a full disc about \(t_c=0\), where
\[
 b=\beta+\tau,\qquad s=\tau-h,
 \qquad e^{2\pi i s}=t_c,
\]
and has proved the exact identity
\begin{equation}
 Z=s\begin{pmatrix}0&1\\-1&0\end{pmatrix}
 +\begin{pmatrix}6\mu&h\\b-h&\mu\end{pmatrix}.
 \label{cusp-spec-canonical-cusp-input}
\end{equation}
Choose \(T_0\) so that \(e^{-2\pi T}\) lies in the proved cusp disc whenever
\(T\ge T_0\).  On the positive radial ray
\[
 t_c=e^{-2\pi T},\qquad T\ge T_0,
\]
we choose the lift \(s=iT\).  The other lifts differ by the already accounted
for integral deck transformation.  Put
\begin{equation}
 \begin{aligned}
  \mu_T&:=\mu(e^{-2\pi T})=\alpha_T+i m_T,\\
  h_T&:=h(e^{-2\pi T})=\eta_T+i H_T,\\
  b_T&:=b(e^{-2\pi T})=\xi_T+i B_T,\\
  c_T&:=\xi_T-\eta_T,\qquad
  L_T:=T+H_T,\qquad q_T:=L_T-B_T,\\
  D_T&:=L_Tq_T+6m_T^2.
 \end{aligned}
 \label{cusp-spec-radial-notation}
\end{equation}
All letters on the right-hand side of
\cref{cusp-spec-radial-notation} are real.  We reserve
\begin{equation}
 m:=\operatorname{Im}\mu(0),\qquad
 H:=\operatorname{Im}h(0),\qquad
 B:=\operatorname{Im}b(0)
 \label{cusp-spec-limiting-imaginary-parts}
\end{equation}
for their limiting values.  Holomorphy at the origin gives, with the same
statement for the real parts,
\begin{equation}
 m_T=m+O(e^{-2\pi T}),\qquad
 H_T=H+O(e^{-2\pi T}),\qquad
 B_T=B+O(e^{-2\pi T}).
 \label{cusp-spec-holomorphic-errors}
\end{equation}

Identify \(\mathbb C^2\) with \(\mathbb R^4\) by the ordered coordinates
\begin{equation}
 (z_1,z_2)=(x_1+iy_1,x_2+iy_2)
 \longmapsto (x_1,y_1,x_2,y_2).
 \label{cusp-spec-real-coordinate-order}
\end{equation}
The metric in this section is, by definition, the quotient of
\(dx_1^2+dy_1^2+dx_2^2+dy_2^2\).  This makes the metric choice explicit; the
complex period construction itself neither requires nor singles out this
Euclidean metric.

\begin{proposition}[Exact radial realification and dual lattice]
\label{cusp-spec-exact-dual-lattice}
In the domain basis
\((\widehat\gamma,\widehat u,\widehat w,\widehat\delta)\) and the real
coordinate order \cref{cusp-spec-real-coordinate-order}, the positive radial
period matrix is
\begin{equation}
 P_T=
 \begin{pmatrix}
  6\alpha_T&\eta_T&1&0\\
  6m_T&L_T&0&0\\
  c_T&\alpha_T&0&1\\
  -q_T&m_T&0&0
 \end{pmatrix}.
 \label{cusp-spec-real-period-matrix}
\end{equation}
Its oriented determinant, positive covolume, and relation with the determinant
function in \cref{prop:period-input-audit} are
\begin{equation}
 \det P_T=-D_T,
 \qquad \operatorname{covol}(P_T\mathbb Z^4)=D_T>0,
 \qquad
 D_{\mathrm{can},T}=-\frac{D_T}{L_T}<0.
 \label{cusp-spec-determinant}
\end{equation}
Here
\(D_{\mathrm{can},T}=\operatorname{Im}\beta_T
-6(\operatorname{Im}\mu_T)^2/\operatorname{Im}\tau_T\) is the canonical
determinant function; in particular, \(D_T\) is not a new nondegeneracy
assumption.

The complete inverse transpose is
\begin{equation}
 P_T^{-\mathsf T}=
 \begin{pmatrix}
 0&0&1&0\\[2pt]
 \dfrac{m_T}{D_T}&\dfrac{q_T}{D_T}&
 \dfrac{-6m_T\alpha_T-q_T\eta_T}{D_T}&
 \dfrac{-m_Tc_T-q_T\alpha_T}{D_T}\\[7pt]
 0&0&0&1\\[2pt]
 -\dfrac{L_T}{D_T}&\dfrac{6m_T}{D_T}&
 \dfrac{6(L_T\alpha_T-m_T\eta_T)}{D_T}&
 \dfrac{L_Tc_T-6m_T\alpha_T}{D_T}
 \end{pmatrix}.
 \label{cusp-spec-full-inverse-transpose}
\end{equation}
Equivalently, for
\(k=(r,s,n_1,n_2)^{\mathsf t}\in\mathbb Z^4\), define
\begin{equation}
 R_T(k):=r-6\alpha_Tn_1-c_Tn_2,
 \qquad
 S_T(k):=s-\eta_Tn_1-\alpha_Tn_2.
 \label{cusp-spec-shifted-indices}
\end{equation}
Then the Euclidean dual lattice is exactly
\begin{equation}
 \begin{split}
 (P_T\mathbb Z^4)^*
 =\biggl\{\biggl(&n_1,
 \frac{m_TR_T(k)+q_TS_T(k)}{D_T},n_2,
 \frac{-L_TR_T(k)+6m_TS_T(k)}{D_T}\biggr):\\[-2pt]
 &k=(r,s,n_1,n_2)\in\mathbb Z^4\biggr\}.
 \end{split}
 \label{cusp-spec-exact-dual-lattice-formula}
\end{equation}
Consequently, if
\[
 F_T:=\mathbb R^4/P_T\mathbb Z^4
\]
 and \(\Delta_T\) is the nonnegative self-adjoint scalar Laplacian on
 \(L^2(F_T,\mathrm d\operatorname{vol}_{F_T})\), obtained from the closed
 Dirichlet form
 \[
  \mathfrak q_T[f]=\int_{F_T}|\mathrm df|^2\,
       \mathrm d\operatorname{vol}_{F_T},
  \qquad \operatorname{Dom}(\mathfrak q_T)=H^1(F_T),
 \]
 with operator domain \(\operatorname{Dom}(\Delta_T)=H^2(F_T)\), then, with
multiplicity,
\begin{equation}
 \operatorname{Spec}(\Delta_T)
 =\{\lambda_T(k):k\in\mathbb Z^4\},
 \label{cusp-spec-spectrum-set}
\end{equation}
where
\begin{equation}
 \begin{aligned}
 \lambda_T(k)=4\pi^2\biggl(&n_1^2+n_2^2\\
 &+\frac{(m_TR_T(k)+q_TS_T(k))^2
       +(-L_TR_T(k)+6m_TS_T(k))^2}{D_T^2}\biggr).
 \end{aligned}
 \label{cusp-spec-exact-eigenvalue}
\end{equation}
The eigenfunction belonging to \(k\) is the character
\begin{equation}
 \chi_{T,k}([x])
 =\exp\bigl(2\pi i\langle P_T^{-\mathsf T}k,x\rangle\bigr).
 \label{cusp-spec-characters}
\end{equation}
\end{proposition}

\begin{proof}
Substitution of \(s=iT\) into
\cref{cusp-spec-canonical-cusp-input} gives
\[
 Z_T=
 \begin{pmatrix}
  6\mu_T&h_T+iT\\ b_T-h_T-iT&\mu_T
 \end{pmatrix}.
\]
Realifying its four columns in the interlaced coordinate order
\((x_1,y_1,x_2,y_2)\) gives
\cref{cusp-spec-real-period-matrix}; in particular, the two divergent entries
have signs \(+L_T\) and \(-q_T=B_T-L_T\).

Expanding the determinant along the last two columns gives
\[
 \det P_T=-L_Tq_T-6m_T^2=-D_T.
\]
Moreover,
\[
 \operatorname{Im}\tau_T=L_T,
 \qquad \operatorname{Im}\beta_T=B_T-L_T=-q_T,
\]
so the canonical identity
\(\det_{\mathbb R}\Pi_T=\operatorname{Im}\tau_TD_{\mathrm{can},T}\)
becomes
\[
 L_TD_{\mathrm{can},T}
 =L_T\left(-q_T-\frac{6m_T^2}{L_T}\right)=-D_T.
\]
The already-proved inequalities \(L_T>0\) and
\(D_{\mathrm{can},T}<0\) imply \(D_T>0\).

To compute the inverse transpose without suppressing the real parts, solve
\(P_T^{\mathsf T}x=k\).  Its third and fourth equations give
\[
 x_1=n_1,\qquad x_3=n_2.
\]
The remaining equations are
\begin{equation}
 \begin{pmatrix}6m_T&-q_T\\L_T&m_T\end{pmatrix}
 \begin{pmatrix}x_2\\x_4\end{pmatrix}
 =\begin{pmatrix}R_T(k)\\S_T(k)\end{pmatrix}.
 \label{cusp-spec-two-by-two-system}
\end{equation}
The determinant of the matrix in
\cref{cusp-spec-two-by-two-system} is \(D_T\), and its inverse is
\[
 \frac1{D_T}
 \begin{pmatrix}m_T&q_T\\-L_T&6m_T\end{pmatrix}.
\]
This proves both \cref{cusp-spec-full-inverse-transpose} and
\cref{cusp-spec-exact-dual-lattice-formula}.

A character \(x\mapsto e^{2\pi i\langle\ell,x\rangle}\) descends to
\(F_T\) precisely when
\(\langle\ell,P_Tk\rangle\in\mathbb Z\) for every
\(k\in\mathbb Z^4\), which is precisely
\(\ell\in P_T^{-\mathsf T}\mathbb Z^4\).  These characters are the complete
Fourier basis of the compact abelian group \(F_T\), and applying the positive
Euclidean Laplacian gives \(4\pi^2\lVert\ell\rVert^2\).  Inserting
\cref{cusp-spec-exact-dual-lattice-formula} proves
\cref{cusp-spec-spectrum-set,cusp-spec-exact-eigenvalue,cusp-spec-characters}.
\end{proof}

\begin{theorem}[Shortest cusp modes and their first corrections]
\label{cusp-spec-shortest-modes}
Let
\begin{equation}
 \ell_{\gamma,T}:=P_T^{-\mathsf T}(1,0,0,0)^{\mathsf t},
 \qquad
 \ell_{u,T}:=P_T^{-\mathsf T}(0,1,0,0)^{\mathsf t}.
 \label{cusp-spec-unit-dual-modes}
\end{equation}
For all sufficiently large \(T\), every shortest nonzero vector of
\((P_T\mathbb Z^4)^*\) belongs to
\[
 \{\pm\ell_{\gamma,T},\ \pm\ell_{u,T}\}.
\]
Their exact squared lengths are
\begin{equation}
 \lVert\ell_{\gamma,T}\rVert^2
 =\frac{L_T^2+m_T^2}{D_T^2},
 \qquad
 \lVert\ell_{u,T}\rVert^2
 =\frac{q_T^2+36m_T^2}{D_T^2}.
 \label{cusp-spec-unit-mode-exact-lengths}
\end{equation}
Their exact comparison is
\begin{equation}
 D_T^2\bigl(\lVert\ell_{\gamma,T}\rVert^2
              -\lVert\ell_{u,T}\rVert^2\bigr)
 =B_T(2L_T-B_T)-35m_T^2.
 \label{cusp-spec-unit-mode-comparison}
\end{equation}
Thus the shortest Fourier indices are \(\pm(0,1,0,0)\) if the right-hand
side of \cref{cusp-spec-unit-mode-comparison} is positive,
\(\pm(1,0,0,0)\) if it is negative, and both pairs if it is zero.

In the standard basis \(e_1,e_2,e_3,e_4\) of the real coordinate space,
the vector expansions are
\begin{align}
 T\ell_{\gamma,T}
 &=-e_4+\frac{m e_2+(H-B)e_4}{T}+O(T^{-2}),
 \label{cusp-spec-gamma-vector-correction}\\
 T\ell_{u,T}
 &=e_2+\frac{-He_2+6me_4}{T}+O(T^{-2}).
 \label{cusp-spec-u-vector-correction}
\end{align}
Accordingly,
\begin{align}
 4\pi^2\lVert\ell_{\gamma,T}\rVert^2
 &=\frac{4\pi^2}{T^2}
   \left(1+\frac{2(B-H)}{T}+O(T^{-2})\right),
 \label{cusp-spec-gamma-eigenvalue-correction}\\
 4\pi^2\lVert\ell_{u,T}\rVert^2
 &=\frac{4\pi^2}{T^2}
   \left(1-\frac{2H}{T}+O(T^{-2})\right).
 \label{cusp-spec-u-eigenvalue-correction}
\end{align}
More generally, every fixed plane index has the first correction
\begin{equation}
 \frac{T^2\lambda_T(r,s,0,0)}{4\pi^2}
 =r^2+s^2+
 \frac{2(B-H)r^2-10mrs-2Hs^2}{T}+O_{r,s}(T^{-2}).
 \label{cusp-spec-general-plane-mode-correction}
\end{equation}
In particular, the first positive eigenvalue, including the coefficient after
its leading term, is
\begin{equation}
 \boxed{
 \lambda_1(F_T)
 =\frac{4\pi^2}{T^2}
  \left(1+
   \frac{-2H+2\min\{B,0\}}{T}+O(T^{-2})\right).}
 \label{cusp-spec-first-eigenvalue-correction}
\end{equation}
If \(B>0\), the bottom indices are eventually
\(\pm(0,1,0,0)\); if \(B<0\), they are eventually
\(\pm(1,0,0,0)\).  If \(B=0\), both pairs have the same displayed
first correction, and the exact decision is made by
\cref{cusp-spec-unit-mode-comparison}.  In particular, if \(B=0\) but
\(m\ne0\), the bottom indices are eventually \(\pm(1,0,0,0)\).
The first positive eigenvalue has complex multiplicity two when exactly one
pair minimizes and multiplicity four when both pairs minimize.
\end{theorem}

\begin{proof}
For \(k=(r,s,0,0)\), define
\[
 K_T:=\frac{T}{D_T}
 \begin{pmatrix}m_T&q_T\\-L_T&6m_T\end{pmatrix}.
\]
Equations \cref{cusp-spec-holomorphic-errors,cusp-spec-radial-notation}
give
\begin{equation}
 K_T=
 \begin{pmatrix}0&1\\-1&0\end{pmatrix}+O_{\mathrm{op}}(T^{-1}).
 \label{cusp-spec-scaled-dual-plane}
\end{equation}
Hence there is a constant \(C\), independent of
\((r,s)\in\mathbb R^2\), such that
\begin{equation}
 (1-C/T)\sqrt{r^2+s^2}
 \le T\lVert P_T^{-\mathsf T}(r,s,0,0)\rVert
 \le(1+C/T)\sqrt{r^2+s^2}
 \label{cusp-spec-uniform-singular-bound}
\end{equation}
for all sufficiently large \(T\).  Thus among the vectors with
\((n_1,n_2)=(0,0)\), every vector with \(r^2+s^2\ge2\) is eventually
longer than all four vectors for which \(r^2+s^2=1\).

If instead \((n_1,n_2)\ne(0,0)\),
\cref{cusp-spec-exact-dual-lattice-formula} gives
\[
 \lVert P_T^{-\mathsf T}k\rVert^2\ge n_1^2+n_2^2\ge1,
\]
whereas each of the four unit vectors just identified has length
\(T^{-1}(1+O(T^{-1}))\).  This proves the asserted exhaustion of the
shortest vectors.

Taking the first two columns of
\cref{cusp-spec-full-inverse-transpose} proves
\cref{cusp-spec-unit-mode-exact-lengths}.  Since \(q_T=L_T-B_T\),
their numerator difference is
\[
 L_T^2+m_T^2-q_T^2-36m_T^2
 =B_T(2L_T-B_T)-35m_T^2,
\]
which proves \cref{cusp-spec-unit-mode-comparison} and the exact mode
classification.

Finally,
\begin{equation}
 D_T
 =T^2+(2H-B)T+H(H-B)+6m^2+O(Te^{-2\pi T}).
 \label{cusp-spec-D-expansion}
\end{equation}
Expanding the first two columns of
\cref{cusp-spec-full-inverse-transpose} with
\cref{cusp-spec-D-expansion} gives
\cref{cusp-spec-gamma-vector-correction,cusp-spec-u-vector-correction}.
Taking their squared norms gives
\cref{cusp-spec-gamma-eigenvalue-correction,cusp-spec-u-eigenvalue-correction}.
The same multiplication for a fixed \((r,s)\) gives
\cref{cusp-spec-general-plane-mode-correction}; explicitly, the scaled
two-dimensional matrix is
\[
 \begin{pmatrix}0&1\\-1&0\end{pmatrix}
 +\frac1T\begin{pmatrix}m&-H\\H-B&6m\end{pmatrix}
 +O(T^{-2}),
\]
and twice the Euclidean inner product of its leading and first-order terms is
\(2(B-H)r^2-10mrs-2Hs^2\).
The smaller relative correction is
\[
 \min\{2(B-H),-2H\}=-2H+2\min\{B,0\},
\]
which proves \cref{cusp-spec-first-eigenvalue-correction}.
\end{proof}

\begin{corollary}[Volume and unit-volume scaling]
\label{cusp-spec-unit-volume}
The Euclidean volume is exactly
\begin{equation}
 V_T:=\operatorname{Vol}(F_T)=D_T
 =(T+H_T)(T+H_T-B_T)+6m_T^2,
 \label{cusp-spec-exact-volume}
\end{equation}
and it has the expansion \cref{cusp-spec-D-expansion}.  Define the
unit-volume metric
\begin{equation}
 \widehat g_T:=D_T^{-1/2}g_T.
 \label{cusp-spec-unit-volume-metric}
\end{equation}
Let \(\widehat\Delta_T\) denote the nonnegative self-adjoint scalar
Laplacian on
\(L^2(F_T,\mathrm d\operatorname{vol}_{\widehat g_T})\) obtained from the
closed Dirichlet form
\[
 \widehat{\mathfrak q}_T[f]
 =\int_{F_T}|\mathrm df|_{\widehat g_T}^2\,
       \mathrm d\operatorname{vol}_{\widehat g_T},
 \qquad
 \operatorname{Dom}(\widehat{\mathfrak q}_T)=H^1(F_T),
\]
with operator domain
\(\operatorname{Dom}(\widehat\Delta_T)=H^2(F_T)\).  Then the first
positive eigenvalue of \(\widehat\Delta_T\) is
\begin{equation}
 \boxed{
 \lambda_1(F_T,\widehat g_T)
 =\frac{4\pi^2}{T}
 \left(1+
  \frac{-H-B/2+2\min\{B,0\}}{T}+O(T^{-2})\right).}
 \label{cusp-spec-unit-volume-gap}
\end{equation}
Thus unit-volume normalization does not prevent the first eigenvalue from
tending to zero.  In the normalized metric the periods \(e_1,e_3\) have
length \(D_T^{-1/4}=T^{-1/2}(1+O(T^{-1}))\), while the two periods with
indices \(\widehat\gamma,\widehat u\) have length
\(T^{1/2}(1+O(T^{-1}))\).
\end{corollary}

\begin{proof}
The volume of a Euclidean lattice quotient is its covolume, so
\cref{cusp-spec-determinant} proves
\cref{cusp-spec-exact-volume}.  In real dimension four, replacing \(g\) by
\(ag\) multiplies volume by \(a^2\) and every Laplace eigenvalue by
\(a^{-1}\).  Therefore \cref{cusp-spec-unit-volume-metric} has volume one
and
\[
 \lambda_1(F_T,\widehat g_T)=D_T^{1/2}\lambda_1(F_T,g_T).
\]
Equation \cref{cusp-spec-D-expansion} gives
\[
 D_T^{1/2}
 =T\left(1+\frac{H-B/2}{T}+O(T^{-2})\right).
\]
Multiplication by \cref{cusp-spec-first-eigenvalue-correction} proves
\cref{cusp-spec-unit-volume-gap}.  The stated period-length estimates follow
directly from \cref{cusp-spec-real-period-matrix}: \(e_1,e_3\) have original
length one, whereas the \(\widehat\gamma\)- and \(\widehat u\)-columns have
original length \(T(1+O(T^{-1}))\); lengths are multiplied by
\(D_T^{-1/4}\).
\end{proof}

\begin{proposition}[Exact counting formula and its two cusp regimes]
\label{cusp-spec-counting}
For \(E\ge0\), count eigenvalues with their Fourier multiplicities by
\begin{equation}
 N_T(E):=\#\{k\in\mathbb Z^4:\lambda_T(k)\le E\}.
 \label{cusp-spec-counting-definition}
\end{equation}
Put
\[
 R_E^2:=\frac{E}{4\pi^2},\qquad
 (x)_+:=\max\{x,0\}.
\]
For every fixed \(E\),
\begin{equation}
 N_T(E)
 =\pi D_T\sum_{n\in\mathbb Z^2}
        (R_E^2-\lVert n\rVert^2)_+ +O_E(T+1).
 \label{cusp-spec-fixed-energy-count}
\end{equation}
The sum is finite.  Equivalently,
\begin{equation}
 \lim_{T\to\infty}\frac{N_T(E)}{D_T}
 =\pi\sum_{n\in\mathbb Z^2}
        (R_E^2-\lVert n\rVert^2)_+.
 \label{cusp-spec-counting-density-limit}
\end{equation}
In the first band, \(0<E<4\pi^2\), this reduces to
\begin{equation}
 N_T(E)=\frac{E}{4\pi}D_T+O_E(T+1).
 \label{cusp-spec-first-band-count}
\end{equation}

For the unit-volume metric \(\widehat g_T\), put
\begin{equation}
 \widehat N_T(E):=
 \#\{k\in\mathbb Z^4:\lambda_T(k;\widehat g_T)\le E\}.
 \label{cusp-spec-unit-volume-count-definition}
\end{equation}
For every fixed \(E>0\),
\begin{equation}
 \widehat N_T(E)
 =\frac{E}{4\pi}D_T^{1/2}+O_E(D_T^{1/4}+1),
 \qquad
 \lim_{T\to\infty}\frac{\widehat N_T(E)}{D_T^{1/2}}
 =\frac{E}{4\pi}.
 \label{cusp-spec-unit-volume-count}
\end{equation}

There is also a finite counting law at the collapsing spectral scale.  Define
\begin{equation}
 N_T^{\mathrm{coll}}(E)
 :=\#\left\{k\in\mathbb Z^4:
       \lambda_T(k)\le\frac{4\pi^2E}{T^2}\right\}.
 \label{cusp-spec-collapsing-count-definition}
\end{equation}
If \(E\) is not a sum of two integer squares, then, for all sufficiently
large \(T\),
\begin{equation}
 N_T^{\mathrm{coll}}(E)
 =\#\{(r,s)\in\mathbb Z^2:r^2+s^2\le E\}.
 \label{cusp-spec-collapsing-count}
\end{equation}
At a threshold \(E\in\{r^2+s^2:r,s\in\mathbb Z\}\), the sharp statement
available from the zeroth-order limit is
\begin{equation}
 \begin{split}
 \#\{(r,s):r^2+s^2<E\}
 &\le\liminf_{T\to\infty}N_T^{\mathrm{coll}}(E)\\
 &\le\limsup_{T\to\infty}N_T^{\mathrm{coll}}(E)
 \le\#\{(r,s):r^2+s^2\le E\}.
 \end{split}
 \label{cusp-spec-threshold-count}
\end{equation}
For a threshold vector \((r,s)\), the coefficient in
\cref{cusp-spec-general-plane-mode-correction} decides its inclusion whenever
that coefficient is nonzero.  If it vanishes, the exact eigenvalue formula
and, if necessary, higher Taylor coefficients decide the inclusion.
\end{proposition}

\begin{proof}
Set
\begin{equation}
 \mathcal A_T:=\begin{pmatrix}6m_T&-q_T\\L_T&m_T\end{pmatrix},
 \qquad
 \theta_T(n_1,n_2):=
 \begin{pmatrix}6\alpha_Tn_1+c_Tn_2\\
                 \eta_Tn_1+\alpha_Tn_2\end{pmatrix}.
 \label{cusp-spec-counting-matrices}
\end{equation}
Then \(\det\mathcal A_T=D_T\), and the exact eigenvalue formula can be
rewritten as
\begin{equation}
 \frac{\lambda_T(r,s,n_1,n_2)}{4\pi^2}
 =\lVert(n_1,n_2)\rVert^2+
  \left\lVert\mathcal A_T^{-1}
  \left(\begin{pmatrix}r\\s\end{pmatrix}
       -\theta_T(n_1,n_2)\right)\right\rVert^2.
 \label{cusp-spec-counting-eigenvalue-form}
\end{equation}
Fix \(n=(n_1,n_2)\).  If \(\lVert n\rVert^2>R_E^2\), this index contributes
nothing.  Otherwise put
\(\rho_n=(R_E^2-\lVert n\rVert^2)^{1/2}\).  The allowed pairs
\((r,s)\) are the integer points in the translated ellipse
\begin{equation}
 \theta_T(n)+\mathcal A_T\overline B_{\mathbb R^2}(0,\rho_n).
 \label{cusp-spec-counting-ellipse}
\end{equation}
Its area is \(\pi D_T\rho_n^2\).  Its boundary length is \(O_E(T+1)\),
because \(\lVert\mathcal A_T\rVert_{\mathrm{op}}=O(T+1)\).
If \(\rho_n=0\), the set in \cref{cusp-spec-counting-ellipse} is one point
and contributes either zero or one integer point, which is absorbed in the
\(O_E(1)\) remainder below.  We apply the planar area--perimeter estimate only
when \(\rho_n>0\).

For completeness, the elementary planar lattice estimate used here is
\[
 \#(\mathbb Z^2\cap K)=\operatorname{area}(K)
 +O(\operatorname{length}(\partial K)+1)
\]
for every translated compact convex set \(K\).  Indeed, attach to each
integer point its centered unit square and let \(U\) be the union of the
squares centered at points of \(K\cap\mathbb Z^2\).  A point of
\(U\mathbin\triangle K\) lies within distance \(\sqrt2\) of
\(\partial K\): this is immediate for \(U\setminus K\), while a point of
\(K\setminus U\) whose distance from the boundary exceeded \(\sqrt2\) would
have its nearest integer point inside \(K\) and would belong to the
corresponding square.  If \(P(K)\) is the perimeter, the outer parallel-set
formula for a convex polygon gives
\(\operatorname{area}(K+B_\rho)-\operatorname{area}(K)
=\rho P(K)+\pi\rho^2\): its added pieces are the edge rectangles and vertex
sectors, whose angles total \(2\pi\).  Polygonal approximation gives the same
bound for every compact convex \(K\).  Applying it to the outer and inner
boundary strips with \(\rho=\sqrt2\) gives
\(\operatorname{area}(U\mathbin\triangle K)=O(P(K)+1)\).  Since the unit
squares have disjoint interiors and total area
\(\#(K\cap\mathbb Z^2)\), the lattice estimate follows.

Only finitely many \(n\)'s occur for fixed \(E\).  Applying the estimate to
\cref{cusp-spec-counting-ellipse} and summing proves
\cref{cusp-spec-fixed-energy-count}.  Since \(D_T\sim T^2\), division by
\(D_T\) proves \cref{cusp-spec-counting-density-limit}; if
\(0<E<4\pi^2\), only \(n=0\) occurs and
\cref{cusp-spec-first-band-count} follows.

For \(\widehat g_T\), eigenvalues are multiplied by \(D_T^{1/2}\), so
\(\widehat N_T(E)=N_T(E/D_T^{1/2})\).  For large \(T\), only \(n=0\)
can occur in \cref{cusp-spec-counting-eigenvalue-form}.  Its ellipse has
radius parameter
\(\rho=(E/(4\pi^2D_T^{1/2}))^{1/2}\), hence area
\(ED_T^{1/2}/(4\pi)\) and perimeter
\(O_E(TD_T^{-1/4}+1)=O_E(D_T^{1/4}+1)\).  The same lattice estimate proves
\cref{cusp-spec-unit-volume-count}.

For the collapsing scale, any \((n_1,n_2)\ne(0,0)\) contributes at least
\(4\pi^2\), so it contributes nothing for all sufficiently large \(T\).
For \(n_1=n_2=0\),
\cref{cusp-spec-scaled-dual-plane} says uniformly that
\[
 \frac{T^2\lambda_T(r,s,0,0)}{4\pi^2}
 =\lVert K_T(r,s)^{\mathsf t}\rVert^2
 \longrightarrow r^2+s^2.
\]
The lower bound in \cref{cusp-spec-uniform-singular-bound} also shows that
only a fixed finite set of \((r,s)\)'s can enter a fixed rescaled window.
Termwise comparison on that finite set proves
\cref{cusp-spec-collapsing-count} away from thresholds and the two-sided
bound \cref{cusp-spec-threshold-count} at thresholds.
\end{proof}

\begin{theorem}[Pointed limits and expanding local covers]
\label{cusp-spec-pointed-limits}
Let \(o_T\) be the origin of \(F_T\), and set
\begin{equation}
 \mathbb T^2_{\mathrm{re}}
 :=(\mathbb Re_1\oplus\mathbb Re_3)/
   (\mathbb Ze_1\oplus\mathbb Ze_3).
 \label{cusp-spec-real-two-torus}
\end{equation}
With the product Euclidean metric,
\begin{equation}
 (F_T,o_T)\longrightarrow
 (\mathbb T^2_{\mathrm{re}}\times
  (\mathbb Re_2\oplus\mathbb Re_4),(0,0))
 \cong(\mathbb T^2\times\mathbb R^2,(0,0))
 \label{cusp-spec-pointed-torus-plane-limit}
\end{equation}
in pointed smooth, and hence pointed Gromov--Hausdorff, topology.  More
precisely, for every \(R<\infty\), the radius-\(R\) balls in the two sides
of \cref{cusp-spec-pointed-torus-plane-limit} are isometric for all
sufficiently large \(T\).

For \(M\in\mathbb Z_{\ge1}\), define the marked local sublattice
\begin{equation}
 \Lambda_{T,M}:=
 \mathbb Zv_{\gamma,T}\oplus\mathbb Zv_{u,T}
 \oplus\mathbb Z(Me_1)\oplus\mathbb Z(Me_3),
 \label{cusp-spec-cover-sublattice}
\end{equation}
where
\begin{equation}
 \begin{aligned}
 v_{\gamma,T}&=(6\alpha_T,6m_T,c_T,-q_T),\\
 v_{u,T}&=(\eta_T,L_T,\alpha_T,m_T).
 \end{aligned}
 \label{cusp-spec-long-period-vectors}
\end{equation}
Then
\begin{equation}
 \widetilde F_{T,M}:=\mathbb R^4/\Lambda_{T,M},
 \qquad
 c_{T,M}:\widetilde F_{T,M}\longrightarrow F_T
 \label{cusp-spec-degree-cover}
\end{equation}
 is a regular holomorphic local-isometric cover of degree \(M^2\), with deck
 group \((\mathbb Z/M\mathbb Z)^2\).  For every integer-valued
function \(M(T)\to\infty\),
\begin{equation}
 (\widetilde F_{T,M(T)},\widetilde o_{T,M(T)})
 \longrightarrow(\mathbb R^4,0)
 \label{cusp-spec-cover-Euclidean-limit}
\end{equation}
in pointed smooth and pointed Gromov--Hausdorff topology.  Again, every fixed
radius ball is eventually exactly isometric to its Euclidean counterpart.
If the cover is instead equipped with the pullback of the unit-volume metric
\(\widehat g_T\), then it converges in pointed smooth topology, and hence in
pointed Gromov--Hausdorff topology, to \(\mathbb R^4\) if and only if
\begin{equation}
 \frac{M(T)}{D_T^{1/4}}\longrightarrow\infty,
 \quad\text{equivalently}\quad
 \frac{M(T)}{\sqrt T}\longrightarrow\infty.
 \label{cusp-spec-unit-volume-cover-condition}
\end{equation}
By contrast, if every covering torus is itself rescaled to volume one, no
choice of \(M(T)\) gives a pointed smooth limit \(\mathbb R^4\).
\end{theorem}

\begin{proof}
The period lattice is
\begin{equation}
 P_T\mathbb Z^4
 =\mathbb Zv_{\gamma,T}\oplus\mathbb Zv_{u,T}
  \oplus\mathbb Ze_1\oplus\mathbb Ze_3.
 \label{cusp-spec-primal-lattice-decomposition}
\end{equation}
First quotient \(\mathbb R^4\) by
\(\mathbb Ze_1\oplus\mathbb Ze_3\).  This gives the proposed limit
\[
 X_\infty:=\mathbb T^2_{\mathrm{re}}
 \times(\mathbb Re_2\oplus\mathbb Re_4).
\]
The remaining quotient translations are generated by the images of
\(v_{\gamma,T},v_{u,T}\).  The projection of their integer combination
\(av_{\gamma,T}+bv_{u,T}\) to the noncompact
\((e_2,e_4)\)-plane is
\begin{equation}
 Y_T\begin{pmatrix}a\\b\end{pmatrix},
 \qquad
 Y_T:=\begin{pmatrix}6m_T&L_T\\-q_T&m_T\end{pmatrix}.
 \label{cusp-spec-long-plane-matrix}
\end{equation}
More precisely,
\begin{equation}
 Y_T
 =T\begin{pmatrix}0&1\\-1&0\end{pmatrix}
 +\begin{pmatrix}6m_T&H_T\\B_T-H_T&m_T\end{pmatrix}.
 \label{cusp-spec-long-plane-limit}
\end{equation}
Holomorphy makes the second matrix uniformly bounded for \(T\ge T_0\).
Choose \(C\) above its operator norm.  The first matrix is \(T\) times an
orthogonal matrix, so the reverse triangle inequality gives
\begin{equation}
 \left\lVert Y_T\begin{pmatrix}a\\b\end{pmatrix}\right\rVert
 \ge(T-C)\sqrt{a^2+b^2}.
 \label{cusp-spec-long-plane-lower-bound}
\end{equation}
The displacement in \(X_\infty\) of a translation is at least the norm of
its noncompact-plane projection.  Hence every nonidentity element of the
remaining deck group has displacement at least \(T-C\).
Together with the fixed lattice vectors \(e_1,e_3\), this also proves
\begin{equation}
 \operatorname{sys}(P_T\mathbb Z^4)=1,
 \qquad \operatorname{inj}(F_T,g_T)=\frac12
 \label{cusp-spec-original-injectivity-radius}
\end{equation}
for all sufficiently large \(T\).

Fix \(R\).  Once \(T-C>4R\), the quotient map
\(X_\infty\to F_T\) is injective on the radius-\(R\) ball about the origin.
It is a local isometry, and every minimizing geodesic of length at most
\(R\) in the quotient lifts from the origin, so the closed radius-\(R\)
ball maps onto the closed radius-\(R\) ball in \(F_T\).  Moreover, a
nontrivial translate cannot
shorten the distance between two points of the ball: their original distance
is at most \(2R\), while every translated distance is at least
\((T-C)-2R>2R\).  The balls are therefore isometric.  This
proves \cref{cusp-spec-pointed-torus-plane-limit}, including its pointed
smooth assertion.

Because the four vectors in
\cref{cusp-spec-primal-lattice-decomposition} are a lattice basis, the
sublattice \cref{cusp-spec-cover-sublattice} has index \(M^2\).  Thus the
identity of \(\mathbb R^4\) induces the regular holomorphic degree-\(M^2\)
local-isometric cover \cref{cusp-spec-degree-cover}.  The quotient of the two
lattices is \((\mathbb Z/M\mathbb Z)^2\), which is its deck group.  Consider
a nonzero vector of
\(\Lambda_{T,M}\):
\begin{equation}
 av_{\gamma,T}+bv_{u,T}+M(pe_1+qe_3),
 \qquad a,b,p,q\in\mathbb Z.
 \label{cusp-spec-general-cover-vector}
\end{equation}
If \((a,b)\ne(0,0)\), its \((e_2,e_4)\)-projection and
\cref{cusp-spec-long-plane-lower-bound} show that its norm is at least
\(T-C\).  If \((a,b)=(0,0)\), its norm is at least \(M\).  Therefore
\begin{equation}
 \operatorname{sys}(\Lambda_{T,M})\ge\min\{T-C,M\}.
 \label{cusp-spec-cover-systole}
\end{equation}
Conversely, \(Me_1\) has length \(M\), while either long generator has
length at most \(T+C'\) for a fixed \(C'\).  Hence
\begin{equation}
 \min\{T-C,M\}
 \le\operatorname{sys}(\Lambda_{T,M})
 \le\min\{T+C',M\},
 \qquad
 \operatorname{inj}(\widetilde F_{T,M})
 =\tfrac12\operatorname{sys}(\Lambda_{T,M}).
 \label{cusp-spec-cover-systole-two-sided}
\end{equation}
For \(M=M(T)\to\infty\), this lower bound tends to infinity.  Once it is
greater than \(4R\), the universal covering map
\(\mathbb R^4\to\widetilde F_{T,M(T)}\) identifies the radius-\(R\) balls
isometrically.  This proves \cref{cusp-spec-cover-Euclidean-limit}.

For the pullback of \(\widehat g_T\), all lengths in
\cref{cusp-spec-cover-systole} are multiplied by \(D_T^{-1/4}\).  The long
lower bound becomes
\((T-C)D_T^{-1/4}\sim\sqrt T\), which diverges automatically.  The explicit
lattice vector \(Me_1\) has normalized length \(MD_T^{-1/4}\).  Hence the
normalized systole tends to infinity exactly under
\cref{cusp-spec-unit-volume-cover-condition}.  Divergence gives Euclidean
balls by the preceding covering argument.  We now prove necessity from the
particular flat lattices rather than from any general injectivity-radius
principle.  Put
\[
 \ell_T:=M(T)D_T^{-1/4},
 \qquad A_T:=(T-C)D_T^{-1/4}.
\]
Then \(A_T\sim\sqrt T\to\infty\).  Quotienting the normalized Euclidean
cover first by
\(\ell_T\mathbb Ze_1\oplus\ell_T\mathbb Ze_3\) gives
\[
 X_{\ell_T}:=
 (\mathbb R/\ell_T\mathbb Z)^2\times\mathbb R^2.
\]
Every remaining nonidentity deck transformation has displacement at least
\(A_T\).  Consequently, for each fixed \(R\), once \(A_T>4R\), the
closed radius-\(R\) ball in the normalized covering torus is exactly the
closed radius-\(R\) ball in \(X_{\ell_T}\).  If
\(\ell_T\not\to\infty\), pass to a subsequence with
\(\ell_T\to\ell<\infty\).  When \(\ell>0\), these pointed product spaces
converge to
\((\mathbb R/\ell\mathbb Z)^2\times\mathbb R^2\); when \(\ell=0\),
the two circle factors have diameter tending to zero and the pointed
Gromov--Hausdorff limit is \(\mathbb R^2\).  Neither limit is
\(\mathbb R^4\).  Thus \(\ell_T\to\infty\), which is precisely
\cref{cusp-spec-unit-volume-cover-condition}, is necessary for both the
pointed smooth and the pointed Gromov--Hausdorff conclusion.

Finally, the volume of the covering torus with its unscaled pulled-back metric
is \(M^2D_T\).  If one instead puts
\begin{equation}
 g^{\mathrm{cov},1}_{T,M}:=(M^2D_T)^{-1/2}g_T
 \quad\text{on }\widetilde F_{T,M},
 \label{cusp-spec-cover-own-unit-volume-metric}
\end{equation}
then its volume is one and its length scale is multiplied by
\((M^2D_T)^{-1/4}\).  The explicit periods \(Me_1\) and
\(v_{\gamma,T}\) consequently have respective lengths
\begin{equation}
 \sqrt{\frac{M}{T}}\,(1+o(1)),
 \qquad
 \sqrt{\frac{T}{M}}\,(1+o(1)).
 \label{cusp-spec-cover-own-unit-volume-lengths}
\end{equation}
Their leading factors have product one, so at least one stays bounded.  The
injectivity radius therefore cannot tend to infinity, which excludes a
pointed smooth \(\mathbb R^4\) limit for the individually unit-volume covers.
\end{proof}

\begin{proposition}[Exact compact quotient of the native cusp fibre]
\label{cusp-spec-native-quotient}
Let
\begin{equation}
 \rho:\mathbb R^4\longrightarrow\mathbb R^2,
 \qquad
 \rho(x_1,x_2,x_3,x_4)=(x_2,x_4),
 \label{cusp-spec-native-projection}
\end{equation}
and let \(Y_T\) be the matrix in
\cref{cusp-spec-long-plane-matrix}.  Then
\begin{equation}
 \det Y_T=D_T>0.
 \label{cusp-spec-long-plane-determinant}
\end{equation}
Consequently
\begin{equation}
 Q_T:=\mathbb R^2/Y_T\mathbb Z^2
 \label{cusp-spec-expanding-quotient}
\end{equation}
is a compact real two-torus, and \(\rho\) induces an exact sequence of
compact real Lie groups
\begin{equation}
 0\longrightarrow\mathbb T^2_{\mathrm{re}}
 \longrightarrow F_T
 \xrightarrow{\ \pi_T\ }Q_T
 \longrightarrow0.
 \label{cusp-spec-native-exact-sequence}
\end{equation}
Moreover,
\begin{equation}
 \operatorname{sys}(Q_T)\ge T-C,
 \label{cusp-spec-expanding-quotient-systole}
\end{equation}
where \(C\) is the constant in
\cref{cusp-spec-long-plane-lower-bound}.  Thus the native
\(\mathbb T^2\times\mathbb R^2\) limit separates exactly into a fixed
unit two-torus kernel and an expanding two-torus quotient.
\end{proposition}

\begin{proof}
Direct calculation gives
\[
 \det Y_T=(6m_T)(m_T)-L_T(-q_T)
 =L_Tq_T+6m_T^2=D_T.
\]
The period decomposition
\cref{cusp-spec-primal-lattice-decomposition} gives
\[
 \rho(P_T\mathbb Z^4)=Y_T\mathbb Z^2,
\]
because \(\rho(e_1)=\rho(e_3)=0\) and the projected long generators
are the two columns of \(Y_T\).  Hence
\[
 \pi_T([x]_{P_T\mathbb Z^4})
 :=[\rho(x)]_{Y_T\mathbb Z^2}
\]
is well defined and surjective.

It remains to determine its kernel without suppressing the integral
lattices.  If \(\pi_T([x])=0\), there are \(a,b\in\mathbb Z\) such that
\[
 \rho(x)=Y_T(a,b)^{\mathsf t}
 =\rho(av_{\gamma,T}+bv_{u,T}).
\]
Thus
\[
 x-av_{\gamma,T}-bv_{u,T}
 \in\mathbb Re_1\oplus\mathbb Re_3,
\]
so \([x]\) lies in the displayed copy of
\(\mathbb T^2_{\mathrm{re}}\).  Conversely that copy is plainly killed by
\(\pi_T\).  Finally, if a vector of \(P_T\mathbb Z^4\) lies in
\(\mathbb Re_1\oplus\mathbb Re_3\), its long coefficients \(a,b\) satisfy
\(Y_T(a,b)^{\mathsf t}=0\).  Since \(\det Y_T=D_T>0\), one has
\(a=b=0\); the intersection is therefore exactly
\(\mathbb Ze_1\oplus\mathbb Ze_3\).  This proves injectivity at the left
and exactness.  The systolic bound is
\cref{cusp-spec-long-plane-lower-bound}.
\end{proof}

\begin{theorem}[The maximal regular-locus
\texorpdfstring{\(\mathbb R^4\)}{R4} quotient and its compactification
obstruction]
\label{cusp-spec-globalization-obstruction}
\mbox{\Verified}\ 
Let
\[
 \Lambda=\mathbb Z\langle\widehat\gamma,\widehat u,
                   \widehat w,\widehat\delta\rangle,\qquad
 K_\Lambda=\mathbb Z\langle\widehat w,\widehat\delta\rangle,
\]
and define
\[
 \varepsilon_\gamma
 (a\widehat\gamma+b\widehat u+c\widehat w+d\widehat\delta)=a,\qquad
 L_\Lambda=\ker\varepsilon_\gamma
 =\mathbb Z\langle\widehat u,\widehat w,\widehat\delta\rangle.
\tag{G.1}
\]
Then the following statements hold.

\smallskip
\noindent
\emph{(i) Regular group quotient.}
Over \(B^\circ\), the period family \(J=X^\circ=f^{-1}(B^\circ)\) is a
holomorphic family of complex Lie groups, not merely an affine torus
bundle.  The formula
\begin{equation}
 \chi_\gamma\bigl([z,\Pi(z)x]\bigr)
 =\bigl([z],\varepsilon_\gamma(x)\bmod\mathbb Z\bigr)
 \label{cusp-spec-global-circle-map}
\end{equation}
defines a global smooth real-analytic fibrewise group epimorphism
\begin{equation}
 \chi_\gamma:J\longrightarrow B^\circ\times(\mathbb R/\mathbb Z).
 \label{cusp-spec-global-circle-quotient}
\end{equation}
Write
\begin{equation}
 \theta_\gamma:=\operatorname{pr}_{\mathbb R/\mathbb Z}\circ\chi_\gamma:
 J\longrightarrow\mathbb R/\mathbb Z
 \label{cusp-spec-global-circle-coordinate}
\end{equation}
for its circle coordinate.  Thus
\(\chi_\gamma=(f,\theta_\gamma)\); the two symbols have different
codomains and will not be interchanged below.
Its kernel is the connected real three-torus family
\begin{equation}
 \mathcal H:=
 \frac{\mathfrak h_z^\circ\times
       ((L_\Lambda\otimes\mathbb R)/L_\Lambda)}{\Delta},
 \label{cusp-spec-global-kernel-family}
\end{equation}
where \(\Delta=\pi_1(B^\circ)\) acts diagonally through the deck action
on \(\mathfrak h_z^\circ\) and the restricted monodromy on
\(L_\Lambda\).  Thus every fibre of \(\mathcal H\) is
\((L_\Lambda\otimes\mathbb R)/L_\Lambda\); the family need not be a
product.
The quotient is real and is not a holomorphic quotient of complex tori.

\smallskip
\noindent
\emph{(ii) Exact radial quotient metric.}
On the positive radial cusp, the restriction of
\cref{cusp-spec-global-circle-map} to
\(F_T=\mathbb R^4/P_T\mathbb Z^4\) is
\begin{equation}
 \chi_{\gamma,T}([y])
 =\left\langle P_T^{-\mathsf T}e_\gamma,y\right\rangle
   \bmod\mathbb Z,\qquad e_\gamma=(1,0,0,0)^{\mathsf T}.
 \label{cusp-spec-radial-circle-character}
\end{equation}
For the Riemannian quotient of the Euclidean metric, the target circle has
circumference
\begin{equation}
 \ell_T
 =\frac1{\lVert P_T^{-\mathsf T}e_\gamma\rVert}
 =\frac{D_T}{\sqrt{L_T^2+m_T^2}}
 =T+H-B+O(T^{-1}).
 \label{cusp-spec-global-circle-circumference}
\end{equation}
In particular, \(\ell_T\to\infty\).

\smallskip
\noindent
\emph{(iii) Four copies and minimality.}
The source
\begin{equation}
 J^{[4]}=J\times_{B^\circ}J\times_{B^\circ}J\times_{B^\circ}J
 \label{cusp-spec-fourfold-fibre-product}
\end{equation}
is the actual fourfold fibre product over all of \(B^\circ\).  The product
map
\begin{equation}
 \chi_\gamma^{[4]}=(\chi_\gamma,\chi_\gamma,\chi_\gamma,\chi_\gamma):
 J^{[4]}\longrightarrow B^\circ\times(\mathbb R/\mathbb Z)^4
 \label{cusp-spec-fourfold-circle-map}
\end{equation}
has kernel
\(\mathcal H^{[4]}=\mathcal H\times_{B^\circ}\cdots
\times_{B^\circ}\mathcal H\), whose fibres are
\(\mathbb T^{12}\).  Along the synchronized positive
radial cusp, its quotient fibre, with product quotient metric, is
\begin{equation}
 G_T=(\mathbb R/\ell_T\mathbb Z)^4,\qquad
 \operatorname{sys}(G_T)=\ell_T.
 \label{cusp-spec-four-circle-fibre}
\end{equation}
Consequently
\begin{equation}
 (G_T,0)\longrightarrow(\mathbb R^4,0)
 \label{cusp-spec-four-circle-limit}
\end{equation}
in pointed local-isometric, pointed smooth, and pointed
Gromov--Hausdorff topology.  If \(\ell_T>4R\), the radius-\(R\) ball
about zero is exactly Euclidean.

Equivalently, the external Cartesian product gives a quotient
\[
 J^4\longrightarrow(B^\circ)^4\times(\mathbb R/\mathbb Z)^4,
\]
and \(J^{[4]}\) is its pullback along the diagonal
\(B^\circ\to(B^\circ)^4\).  At four independent cusp parameters its target
lattice is
\(\operatorname{diag}(\ell_{T_1},\ldots,\ell_{T_4})\mathbb Z^4\), so the
same pointed limit holds whenever
\(\min_i\ell_{T_i}\to\infty\).

Among connected torus quotients obtained from \(r\) copies by a primitive
full-monodromy-invariant kernel and whose quotient systole tends to infinity
on the radial cusp, the quotient rank is at most \(r\).  Thus four copies
are minimal for a four-dimensional all-expanding quotient, and
\cref{cusp-spec-fourfold-circle-map} attains the bound.

The cusp-local two-copy map
\begin{equation}
 \pi_T\times\pi_T:F_T\times F_T\longrightarrow Q_T\times Q_T
 \label{cusp-spec-doubled-local-quotient}
\end{equation}
from \cref{cusp-spec-native-quotient} also has pointed limit
\(\mathbb R^4\).  It is not a quotient over the full \(B^\circ\), however,
and no finite base cover makes it one.

\smallskip
\noindent
\emph{(iv) Finite fillings.}
At \(p_j\), use the notation of
\cref{prop:candidate-finite-fillings} and put
\[
 (m_1,v_1,\kappa_1)
 =\bigl(3,\widehat\gamma+2\widehat u-4\widehat w,1\bigr),
\]
\[
 (m_2,v_2,\kappa_2)
 =\bigl(4,-\widehat\gamma-3\widehat u+3\widehat w,-1\bigr),\qquad
 \kappa_j=\varepsilon_\gamma(v_j).
\tag{G.2}
\]
In the un-conjugated flat coordinates, the finite filling generator is
\begin{equation}
 \Phi_j(s,x)=
 \left(e^{-2\pi i/m_j}s,A_jx+\frac{v_j}{m_j}\right).
 \label{cusp-spec-finite-affine-action}
\end{equation}
The raw character is not invariant:
\begin{equation}
 \varepsilon_\gamma\left(A_jx+\frac{v_j}{m_j}\right)
 =\varepsilon_\gamma(x)+\frac{\kappa_j}{m_j}\pmod{\mathbb Z}.
 \label{cusp-spec-finite-character-shift}
\end{equation}
It instead descends equivariantly to the Seifert circle family
\begin{equation}
 C_j=
 \frac{\Delta_j\times(\mathbb R/\mathbb Z)}
 {(s,\theta)\sim
  (e^{-2\pi i/m_j}s,\theta+\kappa_j/m_j)}.
 \label{cusp-spec-finite-Seifert-target}
\end{equation}
The total space \(C_j\) is a smooth solid torus, the map
\(C_j\to D_j\) has a multiple central circle of order \(m_j\), and the
induced map \(N_j\to C_j\) is a smooth proper \(\mathbb T^3\)-bundle.
The target translation in \cref{cusp-spec-finite-Seifert-target} is
exactly compatible with the logarithmic gluing map \(h_j\).

For four synchronized copies, the smooth simultaneous-root extension is
\begin{equation}
 \widetilde N_j^{[4]}
 =\frac{(J'_j)^{\times_{\Delta_j}4}}
 {\langle(\widetilde g_j^{\log})^{\times4}\rangle},
 \label{cusp-spec-finite-fourfold-normalization}
\end{equation}
with target
\begin{equation}
 C_j^{[4]}=
 \frac{\Delta_j\times(\mathbb R/\mathbb Z)^4}
 {(s,\theta_1,\ldots,\theta_4)\sim
  (e^{-2\pi i/m_j}s,
   \theta_1+\kappa_j/m_j,\ldots,\theta_4+\kappa_j/m_j)}.
 \label{cusp-spec-finite-fourfold-target}
\end{equation}
The map \(\widetilde N_j^{[4]}\to C_j^{[4]}\) is a smooth proper
\(\mathbb T^{12}\)-bundle.  There is a canonical finite map
\begin{equation}
 \nu_j:\widetilde N_j^{[4]}
 \longrightarrow N_j\times_{D_j}N_j\times_{D_j}N_j\times_{D_j}N_j
 \label{cusp-spec-finite-normalization-map}
\end{equation}
which is an isomorphism over \(D_j^*\), is the normalization of the
reduced ordinary filled fibre product, and has degree \(m_j^3\) on the
reduced central fibre.  These two degrees are \(27\) and \(64\).  Gluing
these two local normalizations to \(J^{[4]}\) by
\((h_j)^{\times4}\), and gluing their targets to the product
\(B^\circ\times(\mathbb R/\mathbb Z)^4\) by
\((h_j^{S^1})^{\times4}\), gives the corresponding normalization and
Seifert target over \(B\setminus\{p_0\}\).

\smallskip
\noindent
\emph{(v) Cusp nonextension.}
Neither the literal map \(\chi_\gamma\) nor
\(\chi_\gamma^{[4]}\) extends continuously through the toric cusp fibre
\(W\).  Thus the four-copy quotient is global over \(B^\circ\), and its
finite-filling normalization is global over \(B\setminus\{p_0\}\), but
the displayed map \(\chi_\gamma^{[4]}\) has no continuous extension to
the compact fourfold fibre product or to \(X^4\) with the stated
regular-fibre restriction, and it does not factor through
\(X^{\wedge4}\cong S^{24}\).  The cohomology class of
\(\chi_\gamma\) does extend over
\(N_0\); this weaker homotopy statement does not extend the displayed
quotient map or its metric.

\smallskip
\noindent
\emph{(vi) Products, smashes, and sphere constructions.}
If the recognition \(X\cong_{\mathrm{diff}}S^6\) is used, then
\begin{equation}
 X^2\cong_{\mathrm{diff}}S^6\times S^6\not\simeq S^{12},\qquad
 H_6(X^2;\mathbb Z)\cong\mathbb Z^2,
 \label{cusp-spec-product-not-sphere}
\end{equation}
and
\begin{equation}
 \begin{aligned}
 X^4&\cong_{\mathrm{diff}}(S^6)^4\not\simeq S^{24},\\
 H_6(X^4;\mathbb Z)&\cong\mathbb Z^4,&
 H_{12}(X^4;\mathbb Z)&\cong\mathbb Z^6,&
 H_{18}(X^4;\mathbb Z)&\cong\mathbb Z^4.
 \end{aligned}
 \label{cusp-spec-four-product-not-sphere}
\end{equation}
The synchronized regular fourfold \(J^{[4]}\) has real dimension \(18\),
not \(24\).  Choosing a base point gives based homeomorphisms
\begin{equation}
X\wedge X\cong S^{12},\qquad X^{\wedge4}\cong S^{24},
\label{cusp-spec-smash-spheres}
\end{equation}
but these smash spheres do not carry either \emph{literal fixed-target}
quotient above.  The later cone constructions
\cref{cusp-cone-S12-route,cusp-cone-S24-local-route,cusp-cone-S24-route}
change the codomain by retaining the radial parameter and collapsing the
obstructing boundary cross-section; they therefore do not contradict this
nonextension theorem.  Moreover,
neither \(S^{12}\) nor \(S^{24}\), with any differentiable structure,
admits an almost-complex structure
\cite[\S15, Proposition~15.1 and the following remark,
p.~434]{BorelSerre1953}; hence neither can be promoted to a compact
complex manifold.  This is an obstruction to a proposed complex
\(S^{12}\)- or \(S^{24}\)-globalization, not to the real quotient or its
pointed metric limit.  Precisely,
\(\pi_T\times\pi_T\) is a map on \(F_T^2\), and it fails to factor
through \(F_T\wedge F_T\); \(\chi_{\gamma,T}^{[4]}\) is a map on
\(F_T^4\), and it fails to factor through \(F_T^{\wedge4}\).
Neither map is a map on \(X^2\) or \(X^4\), and no extension with the
stated fibre restriction exists.

There is one exact but different sphere construction.  If
\(C_T=\mathbb R/\ell_T\mathbb Z\) is based at zero, postcomposition with
the target fat-wedge collapse gives
\begin{equation}
 F_T^4\xrightarrow{\chi_{\gamma,T}^{[4]}}C_T^4
 \longrightarrow C_T^{\wedge4}\cong S^4.
 \label{cusp-spec-target-smash-map}
\end{equation}
This map does factor through \(F_T^{\wedge4}\).  With the quotient length
metric and base point
\([\ell_T/2,\ell_T/2,\ell_T/2,\ell_T/2]\), its pointed limit is
\(\mathbb R^4\), with every radius-\(R\) ball exactly Euclidean once
\(\ell_T>4R\).  The metric is singular at the collapsed fat wedge, and
the map still does not extend to \(X^4\) or \(X^{\wedge4}\).
\end{theorem}

\begin{proof}
\medskip\noindent\emph{Invariant lattices and the regular group law.}\par
The matrices in \cref{prop:monodromy-certificate} give
\[
 (M_0-I)(a,b,c,d)^{\mathsf T}=(0,0,b,-a)^{\mathsf T},
\]
so
\begin{equation}
 K_\Lambda=\ker(M_0-I)=\operatorname{im}(M_0-I).
 \label{cusp-spec-cusp-kernel}
\end{equation}
They also give
\begin{equation}
 \begin{aligned}
 A_1\widehat u&=-\widehat w+\widehat\delta,&
 A_1\widehat w&=\widehat u-\widehat w,&
 A_1\widehat\delta&=\widehat\delta,\\
 A_2\widehat u&=\widehat w,&
 A_2\widehat w&=-\widehat u+\widehat\delta,&
 A_2\widehat\delta&=\widehat\delta.
 \end{aligned}
 \label{cusp-spec-invariant-flag-calculation}
\end{equation}
Thus \(L_\Lambda\) is invariant.  It is the smallest invariant lattice
containing \(K_\Lambda\), because
\(\widehat u=A_1\widehat w+\widehat w\).  Both lattices are primitive.

For a row covector \((a,b,c,d)\), invariance under \(A_1\) gives
\(b=2c\) and \(d=3c\), whereas invariance under \(A_2\) gives
\(b=c\) and \(d=2c\).  Hence \(b=c=d=0\), and
\begin{equation}
 \operatorname{Hom}(\Lambda,\mathbb Z)^\Delta
 =\mathbb Z\varepsilon_\gamma.
 \label{cusp-spec-invariant-character-lattice}
\end{equation}

The family over the cover is
\((\mathfrak h_z^\circ\times\mathbb C^2)/\Lambda\), and the action used
to form \(J\) is
\[
 [z,\zeta]\longmapsto[gz,R_g(z)\zeta].
\]
It is complex linear in every fibre and fixes the zero section.
Fibrewise addition, inversion, and zero are therefore equivariant and
descend to \(J\).  In the flat real-analytic marking
\[
 \Xi(z,x+\Lambda)=[z,\Pi(z)x],
\]
equation (2.3d) makes the action
\[
 (z,x+\Lambda)\longmapsto(gz,A_gx+\Lambda).
\]
Equation \cref{cusp-spec-invariant-character-lattice} proves that
\cref{cusp-spec-global-circle-map} is well defined.  Primitivity of
\(\varepsilon_\gamma\) proves surjectivity and connectedness of the
kernel.  The induced action on \(\Lambda/L_\Lambda\simeq\mathbb Z\) is
trivial, so the target is the product circle bundle.  Since the kernel has
odd real dimension, this is not a holomorphic quotient.

\medskip\noindent\emph{The quotient metric and the four-copy limit.}\par
The first column of
\cref{cusp-spec-full-inverse-transpose} is
\begin{equation}
 P_T^{-\mathsf T}e_\gamma
 =\left(0,\frac{m_T}{D_T},0,-\frac{L_T}{D_T}\right)^{\mathsf T}.
 \label{cusp-spec-circle-dual-vector}
\end{equation}
The quotient metric on \(\mathbb R/\mathbb Z\) is therefore
\(\ell_T^2d\theta^2\).  Equivalently,
\[
 P_T(L_\Lambda\otimes\mathbb R)
 =\operatorname{span}_{\mathbb R}\{e_1,e_3,v_{u,T}\}
\]
has unit normal
\[
 \nu_T=\frac{(0,m_T,0,-L_T)}{\sqrt{L_T^2+m_T^2}},
\]
and
\[
 \operatorname{dist}
 \bigl(v_{\gamma,T},P_T(L_\Lambda\otimes\mathbb R)\bigr)
 =\left|\langle v_{\gamma,T},\nu_T\rangle\right|
 =\frac{L_Tq_T+6m_T^2}{\sqrt{L_T^2+m_T^2}}
 =\ell_T.
\]
The expansions in
\cref{cusp-spec-radial-notation,cusp-spec-holomorphic-errors} give
\[
 D_T=T^2+(2H-B)T+O(1),\qquad
 \sqrt{L_T^2+m_T^2}=T+H+O(T^{-1}),
\]
which proves \cref{cusp-spec-global-circle-circumference}.

After rescaling the circle coordinate by \(\ell_T\), the product target is
\(\mathbb R^4/\ell_T\mathbb Z^4\).  If
\(x,y\in B_R(0)\) and \(0\ne\lambda\in\ell_T\mathbb Z^4\), then
\[
 \lVert x-y+\lambda\rVert
 \ge\lVert\lambda\rVert-\lVert x-y\rVert
 >\ell_T-2R>2R\ge\lVert x-y\rVert
\]
when \(\ell_T>4R\).  Hence the covering projection is an exact isometry
on \(B_R(0)\), proving all three convergences in
\cref{cusp-spec-four-circle-limit}.  The quotient map is global over
\(B^\circ\); only the metric limit is evaluated along the radial cusp.
The separately chosen Euclidean fibre metric is not claimed to be a
canonical global metric away from that cusp.

\medskip\noindent\emph{Minimality and the finite-base-cover obstruction.}\par
Let \(N\subset\Lambda^{\oplus r}\) be a primitive invariant kernel whose
connected quotient has radial systole tending to infinity.  Every element
of \(K_\Lambda^{\oplus r}\) has uniformly bounded physical length, since
\[
 P_T\widehat w=e_1,\qquad P_T\widehat\delta=e_3.
\]
If such an element survived in the quotient lattice, its quotient norm
would be no greater than its source norm and would give a nonzero uniformly
bounded period.  Therefore \(K_\Lambda^{\oplus r}\subset N\).
Invariance and
\cref{cusp-spec-invariant-flag-calculation} then force
\(L_\Lambda^{\oplus r}\subset N\), and
\[
 \operatorname{rank}(\Lambda^{\oplus r}/N)\le4r-3r=r.
\]
This proves the lower bound \(r\ge4\), while
\(N=L_\Lambda^{\oplus4}\) proves sharpness.

For the finite-cover assertion, put
\begin{equation}
 H_*=A_2A_1
 =\begin{pmatrix}
 1&0&0&0\\
 6&1&1&0\\
 0&0&1&0\\
 -5&0&-1&1
 \end{pmatrix}=I+N_*.
 \label{cusp-spec-no-finite-cover-matrix}
\end{equation}
Direct multiplication gives \(N_*^2=0\), and hence
\begin{equation}
 H_*^n\widehat w
 =\widehat w+n(\widehat u-\widehat\delta),\qquad
 H_*^n\widehat\delta=\widehat\delta.
 \label{cusp-spec-no-finite-cover-orbit}
\end{equation}
Thus \(H_*^nK_\Lambda\ne K_\Lambda\) for every \(n\ne0\).
Every finite-index subgroup of the base fundamental group contains a
positive power of the loop with monodromy \(H_*\).  Such a subgroup cannot
preserve \(K_\Lambda\), proving that no finite base cover globalizes the
native two-torus quotient.

\medskip\noindent\emph{The affine finite fillings and the actual overlap
maps.}\par
Equation \cref{cusp-spec-invariant-character-lattice} gives
\(\varepsilon_\gamma A_j=\varepsilon_\gamma\), so
\cref{cusp-spec-finite-character-shift} follows.  Define
\[
 \overline\Phi_j(s,\theta)
 =\left(e^{-2\pi i/m_j}s,\theta+\frac{\kappa_j}{m_j}\right).
\]
Then \((s,x)\mapsto(s,\varepsilon_\gamma(x))\) is
\((\Phi_j,\overline\Phi_j)\)-equivariant; denote the descended map by
\(\overline\chi_j:N_j\to C_j\).  A nonidentity power which fixes
the base must lie over \(s=0\), where it translates the circle by
\(k\kappa_j/m_j\).  Since \(\gcd(\kappa_j,m_j)=1\), this is nonintegral
for \(0<k<m_j\).  The target action is free; the source action is free by
\cref{prop:candidate-finite-fillings}.

For an explicit solid-torus chart, put
\(\lambda_j=e^{2\pi i/m_j}\), \(u=e^{2\pi i\theta}\), and choose
\(r_j\in\mathbb Z\) with \(r_j\kappa_j\equiv1\pmod {m_j}\).  Then
\[
 w=su^{r_j},\qquad U=u^{m_j}
\]
are invariant and identify \(C_j\) real-analytically with a solid torus.
In these coordinates
\[
 t_j=s^{m_j}=w^{m_j}U^{-r_j},
\]
which proves the central multiplicity.  Equivariance and the primitive
kernel give the proper \(\mathbb T^3\)-bundle \(N_j\to C_j\).

It remains to check the actual logarithmic gluing.  Let
\(\Pi_{\mathbb R}(z):\Lambda\otimes\mathbb R\to\mathbb C^2\) be the
underlying real isomorphism.  On a logarithm branch \(\ell\), define
\(b_j^\ell(z)\) uniquely by
\begin{equation}
 \Pi_{\mathbb R}(z)b_j^\ell(z)
 =\frac{\ell(z)}{2\pi i}\Pi(z)v_j,\qquad
 \beta_j(z)=\varepsilon_\gamma(b_j^\ell(z))\bmod\mathbb Z.
 \label{cusp-spec-finite-target-conjugacy}
\end{equation}
Changing \(\ell\) by \(2\pi in\) adds \(nv_j\) to \(b_j^\ell\), hence
adds \(n\kappa_j\in\mathbb Z\) to its
\(\varepsilon_\gamma\)-value.  Thus \(\beta_j\) is branch independent.
In flat coordinates the source conjugacy (2.9e) is translation by
\(b_j^\ell\).  Applying \(\varepsilon_\gamma\) to that conjugacy gives
\begin{equation}
 \beta_j(g_jz)=\beta_j(z)-\frac{\kappa_j}{m_j}\pmod{\mathbb Z}.
 \label{cusp-spec-finite-target-cocycle}
\end{equation}
Therefore
\[
 h_j^{S^1}(z,\theta)=(z,\theta+\beta_j(z))
\]
conjugates \(\overline\Phi_j\) to the linear target action
\((z,\theta)\mapsto(g_jz,\theta)\), and the overlap identity is
\begin{equation}
 \chi_\gamma\circ h_j=h_j^{S^1}\circ\overline\chi_j.
 \label{cusp-spec-finite-overlap-identity}
\end{equation}

The diagonal action on
\((J'_j)^{\times_{\Delta_j}4}\) is free, and the same coprimality argument
proves that \(C_j^{[4]}\) is smooth.  Every fibre of the resulting map is
a \(\mathbb T^{12}\)-torsor; over \(D_j^*\) its structure torus is the
fourfold fibre product of \(\mathcal H\).  No group zero is asserted at
the affine filling.
In local etale coordinates along the central fibre, the ordinary filled
fibre product has equations
\begin{equation}
 s_i^{m_j}=s_1^{m_j},\qquad i=2,3,4.
 \label{cusp-spec-finite-fibre-product-branches}
\end{equation}
It has \(m_j^3\) smooth branches along the reduced central fibre.  The
simultaneous-root model separates them, is smooth and hence normal, and is
an isomorphism away from the central fibre.  The map \(\nu_j\) is proper:
both spaces are proper over \(D_j\), and the inverse image of a compact
subset of the target is closed in the inverse image of its compact image
in \(D_j\).  It is quasi-finite, since its fibres are singletons away from
the central fibre and have at most \(m_j^3\) points over the central
fibre.  Hence \(\nu_j\) is finite.  The universal property of
normalization now proves that it is the normalization of the reduced
ordinary fibre product.  On the central fibre the relative root choices form
\((\mathbb Z/m_j)^4/\operatorname{diag}(\mathbb Z/m_j)\), of order
\(m_j^3\), which proves the asserted degrees.

On \(D_j^*\), the source identification is \((h_j)^{\times4}\), and
\cref{cusp-spec-finite-overlap-identity} identifies the target overlap
with \((h_j^{S^1})^{\times4}\).  These are precisely the fourfolds of
the one-copy overlap maps and therefore satisfy the cocycle identities.
They glue the two local normalizations and Seifert targets to
\(J^{[4]}\) and to
\(B^\circ\times(\mathbb R/\mathbb Z)^4\), respectively.  Equivalently,
the source gluing follows from uniqueness of normalization.

The fixed product-target trivialization does not extend over \(X\), even if
base winding is allowed.  Indeed, in the presentation of
\cref{prop:geometric-pi1}, let a homomorphism to \(\mathbb Z\) restrict to
\(n\varepsilon_\gamma\) on \(\Lambda\) and take values \(a,b\) on
\(\widehat\rho_1,\widehat\rho_2\).  The three attachment relations give
\[
 3a=n,\qquad4b=-n,\qquad a+b=0.
\]
The first two give \(a=n/3\), \(b=-n/4\), while the third gives
\(n/12=0\).  Hence \(n=0\).  This is the exact global obstruction behind
the fact that the two Seifert targets do not combine into a fixed-circle
map on \(X\).

\medskip\noindent\emph{The explicit toric cusp obstruction.}\par
On the positive radial ray, write
\[
 Z_T=iTB_0+C(t_c),\qquad
 B_0=\begin{pmatrix}0&1\\-1&0\end{pmatrix},\qquad
 t_c=e^{-2\pi T}.
\]
For \(a=(a_1,a_2)^{\mathsf T}\in\mathbb R^2\), take the point of the
regular fibre represented by \(Z_Ta\), with zero short-period coordinates,
and exponentiate it to the dense torus of the cusp model:
\[
 x_T(a)=\exp(2\pi iZ_Ta)\in(\mathbb C^*)^2.
\]
Its logarithmic position is
\begin{equation}
 y_T(a)=
 \frac{(\log|x_{T,1}(a)|,\log|x_{T,2}(a)|)}{\log|t_c|}
 =B_0a+\frac{\operatorname{Im}C(t_c)a}{T}.
 \label{cusp-spec-cusp-nonextension-position}
\end{equation}
The circle character has the constant value
\begin{equation}
 \chi_{\gamma,T}(x_T(a))=a_1\bmod\mathbb Z.
 \label{cusp-spec-cusp-nonextension-character}
\end{equation}

Use the affine toric chart of
\(\sigma=\operatorname{conv}(0,e_1,e_2)\).  Its coordinates are
\begin{equation}
 z_0=\frac{t_c}{x_1x_2},\qquad z_1=x_1,\qquad z_2=x_2.
 \label{cusp-spec-cusp-nonextension-chart}
\end{equation}
If \(B_0a=(y_1,y_2)\) is in the interior of \(\sigma\), then
\(y_1>0\), \(y_2>0\), and \(1-y_1-y_2>0\).  For all sufficiently large
\(T\), the three absolute values in
\cref{cusp-spec-cusp-nonextension-chart} are
\[
 |t_c|^{1-y_{T,1}-y_{T,2}},\qquad
 |t_c|^{y_{T,1}},\qquad |t_c|^{y_{T,2}},
\]
and all tend to zero.  Every such path converges to the same toric fixed
point and hence to the same triple point of \(W\).

Take
\[
 a=\left(-\frac14,\frac14\right),\qquad
 a'=\left(-\frac12,\frac14\right).
\]
Then
\[
 B_0a=\left(\frac14,\frac14\right),\qquad
 B_0a'=\left(\frac14,\frac12\right)
\]
are both in the interior of \(\sigma\).  The paths converge to the same
point of \(W\), whereas
\[
 \chi_{\gamma,T}(x_T(a))=\frac34,\qquad
 \chi_{\gamma,T}(x_T(a'))=\frac12.
\]
No continuous extension of the literal map \(\chi_\gamma\) exists.
Choosing identical marked radial paths in the other three factors and
using these two paths in one factor proves the same assertion for
\(\chi_\gamma^{[4]}\).

The radial family of translations by \(\widehat u/2\) in \(\mathcal H\)
fails to extend as well.  Take two paths
whose limiting logarithmic positions are
\[
 y=\left(\frac18,\frac18\right),\qquad
 y'=\left(\frac38,\frac38\right).
\]
Both approach the fixed point of
\(\operatorname{conv}(0,e_1,e_2)\).  Translation by
\(\widehat u/2\in\mathcal H_T\) adds the exact displacement
\[
 \frac12B_0e_2+
 \frac1{2T}\operatorname{Im}C(t_c)e_2
 \longrightarrow\left(\frac12,0\right)
\]
to logarithmic position.  The translated limits are
\[
 \left(\frac58,\frac18\right)
 \quad\hbox{and}\quad
 \left(\frac78,\frac38\right).
\]
The first lies in the interior of
\(\operatorname{conv}(0,e_1,e_2)\), while the second lies in the interior
of \(\operatorname{conv}(e_1,e_2,e_1+e_2)\).  These two triangle types are
the two translation orbits in the \(A_2\) triangulation and give the two
distinct triple points of \(W\), as proved in
\cref{thm:candidate-cusp-geometry}.  Thus one regular-fibre translation
sends two paths with the same limit to paths with different limits; it has
no continuous extension to \(N_0\).

This pointwise obstruction must not be confused with the cohomology class.
By \cref{prop:geometric-pi1},
\[
 \pi_1(N_0)=\Lambda/K_\Lambda,
\]
and \(\varepsilon_\gamma\) kills \(K_\Lambda\).  It therefore defines an
integral \(H^1\)-class on \(N_0\), hence a map to the Eilenberg--Mac Lane
space \(S^1\) whose restriction is homotopic to \(\chi_\gamma\).  The
explicit paths prove that no such representative agrees with
\(\chi_\gamma\) on the whole punctured cusp.  The phase torus
\((K_\Lambda\otimes\mathbb R)/K_\Lambda\) extends as the compact real torus
action of the toric model, generally with boundary stabilizers; the full
radial \(\widehat u/2\)-translation and the full quotient do not.

\medskip\noindent\emph{Products and smash quotients.}\par
The integral Kunneth theorem gives
\[
 H_k(X^n;\mathbb Z)\cong
 \begin{cases}
  \mathbb Z^{\binom nj},&k=6j,\quad0\le j\le n,\\
  0,&\text{otherwise},
 \end{cases}
 \qquad \chi(X^n)=2^n,
\]
which proves
\cref{cusp-spec-product-not-sphere,cusp-spec-four-product-not-sphere}.
The synchronized regular \(n\)-fold fibre product has dimension
\[
 6n-2(n-1)=4n+2,
\]
which gives \(18\) for four copies.

A product map factors through a smash quotient if and only if it is
constant on the corresponding fat wedge.  On
\(F_T\times\{0\}\), the local two-copy quotient is
\[
 x\longmapsto(\pi_T(x),0),
\]
which is nonconstant.  Thus the literal conclusion is that
\(\pi_T\times\pi_T\) does not factor through \(F_T\wedge F_T\); it would
be a type error to call it a map on \(X\wedge X\).  The restriction of
\(\chi_{\gamma,T}^{[4]}\) to \(F_T\times\{0\}^3\) is similarly
\[
 x\longmapsto(\chi_{\gamma,T}(x),0,0,0),
\]
so it does not factor through \(F_T^{\wedge4}\).

There is no global extension which repairs the domain mismatch:
\(H^1(X^n;\mathbb Z)=0\), whereas both local quotient maps pull nonzero
target \(H^1\)-classes back to nonzero classes on the cusp fibre.  Hence no
map on \(X^2\) or \(X^4\) has the required restriction, and precomposition
with the smash collapses in \cref{cusp-spec-smash-spheres} cannot produce
the local quotient maps.

Identify the metric quotient circle with
\(C_T=\mathbb R/\ell_T\mathbb Z\) by
\(\theta\mapsto\ell_T\theta\).  For
\cref{cusp-spec-target-smash-map}, target-smash collapse is constant
whenever one source coordinate is the identity, so the modified map does
factor through \(F_T^{\wedge4}\).  Its chosen base point has distance
\(\ell_T/2\) from the collapsed target fat wedge.  If two points lie
within distance \(R\) of it, a route through the collapsed point has
length at least \(\ell_T-2R\), while their direct Euclidean distance is at
most \(2R\).  When \(\ell_T>4R\), neither the collapse nor circle
wrap-around shortens distances in that ball.

Let \(u_i\in H^1(F_T^4;\mathbb Z)\) be the pullback of the generator from
the \(i\)-th quotient circle.  The orientation class of
\(S^4=C_T^{\wedge4}\) pulls back under
\cref{cusp-spec-target-smash-map} to
\[
 \pm u_1\smile u_2\smile u_3\smile u_4\ne0.
\]
But \(H^4(X^4;\mathbb Z)=0\).  Thus even this modified \(S^4\)-target map
does not extend to \(X^4\), and a fortiori does not arise from a map on
\(X^{\wedge4}\cong S^{24}\).

For completeness, Borel--Serre prove that \(S^{2n}\) admits no
almost-complex structure for \(n\ge4\).  Their argument applies to every
differentiable structure on the sphere: on \(S^{2n}\), all lower Chern
classes vanish, while their reduced-power calculation forces the top
Chern class to vanish modulo a suitable odd prime \(p<n\); on the other
hand that class is the Euler class, namely \(2h\) for a generator
\(h\in H^{2n}(S^{2n};\mathbb Z)\), a contradiction.  Taking \(n=6\) and
\(n=12\) gives the asserted obstructions for \(S^{12}\) and \(S^{24}\).
This topological obstruction does not itself prohibit the real
four-circle target \(G_T\), the target-smash \(S^4\), or their metric
limits.

Finally,
\(X*X\cong S^{13}\) and \(X^{*4}\cong S^{27}\), so joins do not produce
the proposed \(S^{12}\) or \(S^{24}\) dimensions.  Suspensions can produce
spheres of those dimensions only by discarding the torus-fibration maps.
\end{proof}


\subsection{The Kummer cover extends across the toric cusp}
\label{cusp-cover-section}

This section determines exactly what is true in the assertion that the
degree-$N^2$ cover of the regular fibres extends by a toric power map.  The
answer is positive for the completed \emph{local} cusp family, but the source
action has to be changed, and two different degrees must not be conflated.

Retain the cusp data
\[
 B_0=\begin{pmatrix}0&1\\-1&0\end{pmatrix},\qquad
 Z(s)=sB_0+C(t),\qquad t=e^{2\pi i s},
\]
where $C$ is holomorphic on a disc $\Delta_t$ about zero.  Put
\[
 c_{\bar\lambda}(t)
 :=\exp\!\bigl(2\pi iC(t)\bar\lambda\bigr)
 \in(\mathbb C^*)^2,
 \qquad \bar\lambda\in\bar\Lambda\simeq\mathbb Z^2.
\]
Let $\Sigma$ be the fan in
$N'=\mathbb Z^2\oplus\mathbb Z$ consisting of the cones over the cells of
the translation-invariant $A_2$ triangulation at height one, and write
$Y=Y_\Sigma$.  On the dense torus, with coordinates $(x_1,x_2,t)$, the
target action is
\begin{equation}
 \Psi_{\bar\lambda}(x,t)
 =\bigl(c_{\bar\lambda}(t)t^{B_0\bar\lambda}x,t\bigr).
 \label{cusp-cover-target-action}
\end{equation}
Its quotient is the completed cusp family
$f_0\colon N_0=Y/\bar\Lambda\to\Delta_t$.

Fix an integer $N\geq1$.  Let $v$ be a coordinate on a second disc and
make the base change
\begin{equation}
 \rho_N\colon\Delta_v\longrightarrow\Delta_t,\qquad
 \rho_N(v)=v^N,\qquad t=v^N.
 \label{cusp-cover-base-change}
\end{equation}
On a copy $Y^{(N)}$ of the same toric variety, with dense-torus coordinates
$(\xi_1,\xi_2,v)$, define
\begin{equation}
 C_N(v):=\frac1N C(v^N),\qquad
 c^{(N)}_{\bar\lambda}(v)
 :=\exp\!\bigl(2\pi iC_N(v)\bar\lambda\bigr)
 \label{cusp-cover-source-factor}
\end{equation}
and
\begin{equation}
 \Psi^{(N)}_{\bar\lambda}(\xi,v)
 :=\bigl(c^{(N)}_{\bar\lambda}(v)
          v^{B_0\bar\lambda}\xi,v\bigr).
 \label{cusp-cover-source-action}
\end{equation}

\begin{theorem}[Exact Kummer extension through the cusp]
\label{cusp-cover-extension-theorem}
\mbox{\Verified}\ 
After shrinking the $v$-disc, the following statements hold.

\begin{enumerate}
\item The maps in \eqref{cusp-cover-source-action} extend to a free,
properly discontinuous action on $Y^{(N)}$.  The quotient
\[
 f_N\colon N_0^{\langle N\rangle}
 :=Y^{(N)}/\bar\Lambda\longrightarrow\Delta_v
\]
is a proper holomorphic map from a smooth threefold.  Its central fibre
$W_N$ is reduced, irreducible, and nonnormal, and its normalization is a
toric del Pezzo surface of degree six.

\item Multiplication by $N$ on the cocharacter lattice defines a finite
flat toric morphism
\begin{equation}
 P_N\colon Y^{(N)}\longrightarrow Y,\qquad
 P_N(\xi_1,\xi_2,v)=(\xi_1^N,\xi_2^N,v^N),
 \label{cusp-cover-toric-power-map}
\end{equation}
of degree $N^3$.  It is equivariant for
\eqref{cusp-cover-source-action} and
\eqref{cusp-cover-target-action}, hence descends to a finite flat map
\begin{equation}
 \kappa_N\colon N_0^{\langle N\rangle}\longrightarrow N_0,
 \qquad f_0\circ\kappa_N=\rho_N\circ f_N,
 \label{cusp-cover-completed-map}
\end{equation}
of degree $N^3$.

\item The map $\kappa_N$ is etale over $v\ne0$, and
\begin{equation}
 \kappa_N^*W=NW_N,
 \qquad
 \operatorname{Ram}(\kappa_N)=(N-1)W_N.
 \label{cusp-cover-direct-ramification}
\end{equation}
In particular, the branch support of the map to the original cusp family
is the whole central fibre.

\item Let
\begin{equation}
 N_0[N]:=N_0\times_{\Delta_t}\Delta_v,
 \qquad t=v^N.
 \label{cusp-cover-base-changed-target}
\end{equation}
Then $N_0[N]$ is normal and $\kappa_N$ factors uniquely as
\begin{equation}
 N_0^{\langle N\rangle}
 \xrightarrow{\ Q_N\ }N_0[N]
 \longrightarrow N_0.
 \label{cusp-cover-relative-factorization}
\end{equation}
The map $Q_N$ has generic degree $N^2$.  For $N>1$ it is quasi-etale but
not flat: it is etale over the smooth locus of $N_0[N]$, and its non-etale
and nonflat loci are exactly the inverse images of the double and triple
strata of the base-changed central fibre.  Thus $Q_N$ is not a
normalization map.

\item On the reduced central fibres, the induced map
$W_N\to W$ has generic degree $N^2$.  It lifts to a commutative diagram
\begin{equation}
\begin{CD}
 E_N\simeq\mathrm{dP}_6 @>{[N]}>> E\simeq\mathrm{dP}_6\\
 @V{\nu_N}VV @VV{\nu}V\\
 W_N @>>> W,
\end{CD}
 \label{cusp-cover-normalization-diagram}
\end{equation}
where $[N]$ is toric multiplication by $N$, of degree $N^2$, ramified
with index $N$ along each of the six toric boundary curves.  The diagram
respects the paired conductor curves and the two triple-point classes.

\item Over $v\ne0$, $Q_N$ is exactly the cover defined by
\[
 \mathbb Z\langle\widehat\gamma,\widehat u,
 N\widehat w,N\widehat\delta\rangle
 \subset
 \mathbb Z\langle\widehat\gamma,\widehat u,
 \widehat w,\widehat\delta\rangle.
\]
Its deck group is $(\mathbb Z/N\mathbb Z)^2$.  The pulled-back cusp
monodromy is $M_0^N$, and in the displayed cover basis it is $M_0$:
\begin{equation}
 S_N^{-1}M_0^NS_N=M_0,\qquad
 S_N=\operatorname{diag}(1,1,N,N).
 \label{cusp-cover-monodromy-conjugacy}
\end{equation}
\end{enumerate}
\end{theorem}

\begin{proof}
We prove the assertions in the order in which the constructions occur.

\medskip\noindent
\emph{The source action.}
For $\bar\lambda\in\bar\Lambda$, the integral linear map
\[
 \varphi_{\bar\lambda}(y,y_3)
 =(y+y_3B_0\bar\lambda,y_3)
\]
translates the height-one slice by the integral vector
$B_0\bar\lambda$.  It preserves the $A_2$ triangulation and hence induces
a toric automorphism $\Phi_{\bar\lambda}$ of $Y^{(N)}$.  On the dense
torus this automorphism is
\[
 \Phi_{\bar\lambda}(\xi,v)
 =(v^{B_0\bar\lambda}\xi,v).
\]
Multiplication by the nowhere-zero holomorphic function
$c^{(N)}_{\bar\lambda}(v)$ extends across every toric chart.  Since
$C_N(v)\bar\lambda$ is additive in $\bar\lambda$, the resulting maps
satisfy
$\Psi^{(N)}_{\bar\lambda}\Psi^{(N)}_{\bar\mu}
=\Psi^{(N)}_{\bar\lambda+\bar\mu}$.

For completeness, freeness and proper discontinuity can be checked
directly.  Put
\[
 R_N(v):=-2\pi\operatorname{Im}C_N(v).
\]
It is bounded at $v=0$.  At a dense-torus point fixed by
$\bar\lambda$, logarithmic absolute values give
\[
 \log|v|B_0\bar\lambda+R_N(v)\bar\lambda=0.
\]
Choose the disc so small that
$|\log|v||>2\|R_N\|_\infty$.  Because $B_0$ is orthogonal, the last
equality is impossible for $\bar\lambda\ne0$.  On the central fibre,
$\Psi^{(N)}_{\bar\lambda}$ translates every bounded cell by the nonzero
vector $B_0\bar\lambda$, so it fixes no toric orbit.

For a dense-torus point set
\[
 y(\xi,v)=\frac{(\log|\xi_1|,\log|\xi_2|)}{\log|v|},
 \qquad
 B_{N,v}=B_0+\frac{R_N(v)}{\log|v|}.
\]
Then
\[
 y(\Psi^{(N)}_{\bar\lambda}(\xi,v))
 =y(\xi,v)+B_{N,v}\bar\lambda,
 \qquad
 |B_{N,v}\bar\lambda|\ge\tfrac12|\bar\lambda|.
\]
The logarithmic-position regions associated with any two fixed toric
charts are bounded.  Hence only finitely many $\bar\lambda$ can translate
one of them to meet the other.  Local finiteness of the fan reduces the
same assertion for compact sets to finitely many chart pairs, proving
proper discontinuity.  Finally, reducing
$B_{N,v}^{-1}y(\xi,v)$ modulo $\mathbb Z^2$ places every orbit meeting a
closed subdisc in finitely many closed unit toric charts.  Their union is
compact.  This proves properness of $f_N$.  The central-fibre and
normalization assertions follow exactly from the height-one fan: locally
$v$ is the product of the coordinates belonging to the rays through a
cell, every component occurs with multiplicity one, the lattice action is
transitive on vertices, and the star of a vertex is the complete smooth
hexagonal fan.

\medskip\noindent
\emph{The toric map and equivariance.}
Multiplication by $N$ on $N'$ sends every cone $\sigma\in\Sigma$ onto
itself as a real cone.  It therefore defines the toric morphism $P_N$.
For a maximal cone, let $z_0,z_1,z_2$ be the target characters dual to its
three primitive height-one ray generators, and let $w_0,w_1,w_2$ be the
corresponding source characters.  The map on convergent local rings is
\begin{equation}
 \mathbb C\{z_0,z_1,z_2\}\longrightarrow
\mathbb C\{w_0,w_1,w_2\},\qquad z_i\longmapsto w_i^N.
 \label{cusp-cover-maximal-chart-map}
\end{equation}
The target-chart preimage is the corresponding source chart, and the
right-hand ring is free of rank $N^3$ over the left-hand ring, with basis
\[
 \{w_0^{a_0}w_1^{a_1}w_2^{a_2}:0\le a_i<N\}.
\]
Thus $P_N$ is finite flat of degree $N^3$ on every maximal chart, hence
globally.

The choice \eqref{cusp-cover-source-factor} is forced by equivariance.
Indeed, writing powers coordinatewise,
\begin{align}
 P_N\Psi^{(N)}_{\bar\lambda}(\xi,v)
 &=\bigl((c^{(N)}_{\bar\lambda}(v))^N
          v^{NB_0\bar\lambda}\xi^N,v^N\bigr)\\
 &=\bigl(c_{\bar\lambda}(v^N)
          (v^N)^{B_0\bar\lambda}\xi^N,v^N\bigr)\\
 &=\Psi_{\bar\lambda}P_N(\xi,v).
 \label{cusp-cover-equivariance-computation}
\end{align}
Consequently $P_N$ descends to \eqref{cusp-cover-completed-map}.  To see
finiteness and flatness after the infinite-lattice quotient without any
implicit descent theorem, take an evenly covered open set in $N_0$ and
one of its lifts $V\subset Y$.  Equivariance gives a unique representative
of every source orbit above that open set whose image lies in $V$; two
such representatives cannot differ by a nontrivial lattice translate.
The descended map is therefore locally identified with
$P_N^{-1}(V)\to V$.  It is finite flat of degree $N^3$ by
\eqref{cusp-cover-maximal-chart-map}.

In the same chart the base characters are
\[
 t=z_0z_1z_2,\qquad v=w_0w_1w_2,\qquad P_N^*t=v^N.
\]
The Jacobian determinant of \eqref{cusp-cover-maximal-chart-map} is
\[
 N^3(w_0w_1w_2)^{N-1}=N^3v^{N-1}.
\]
This proves \eqref{cusp-cover-direct-ramification}; the local quotient
maps are biholomorphisms, so the calculation is unchanged downstairs.

\medskip\noindent
\emph{The relative degree-$N^2$ factor.}
The equality
\((f_0\circ\kappa_N)([\xi,v])=v^N=(\rho_N\circ f_N)([\xi,v])\)
gives the factorization
\eqref{cusp-cover-relative-factorization}.  At a triple point the exact
local-ring map is
\begin{equation}
 \begin{gathered}
 A_N:=
 \frac{\mathbb C\{z_0,z_1,z_2,v\}}
      {(z_0z_1z_2-v^N)},
 \qquad
 B:=\mathbb C\{w_0,w_1,w_2\},\\
 A_N\longrightarrow B,\qquad
 z_i\longmapsto w_i^N,\qquad
 v\longmapsto w_0w_1w_2.
 \end{gathered}
 \label{cusp-cover-basechange-local-ring}
\end{equation}
Let
\[
 G_N=\{(\varepsilon_0,\varepsilon_1,\varepsilon_2)\in\mu_N^3:
       \varepsilon_0\varepsilon_1\varepsilon_2=1\}
 \simeq\mu_N^2
\]
act diagonally on $B$.  A monomial
$w_0^{a_0}w_1^{a_1}w_2^{a_2}$ is $G_N$-invariant if and only if
\[
 a_0\equiv a_1\equiv a_2\pmod N.
\]
Writing their common residue as $r$, with $0\le r<N$, expresses that
monomial uniquely in the form
\[
 (w_0w_1w_2)^r
 (w_0^N)^{b_0}(w_1^N)^{b_1}(w_2^N)^{b_2}.
\]
The same coefficientwise argument applies to convergent power series.
Hence
\begin{equation}
 A_N=B^{G_N}.
 \label{cusp-cover-invariant-ring}
\end{equation}
It follows that $A_N$ is normal, that the map has generic degree
$|G_N|=N^2$, and that it is a finite quotient rather than a normalization.
The corresponding face-chart calculations prove normality along the
double and smooth strata as well, so $N_0[N]$ is normal everywhere.

The stabilizer in $G_N$ is trivial when none or exactly one of the $w_i$
vanishes, is isomorphic to $\mu_N$ when exactly two vanish, and is all of
$G_N$ when all three vanish.  Thus $Q_N$ is etale in codimension one and
has no ramification divisor, but it is not etale above the double and
triple strata.  It is also nonflat at every such point when $N>1$.  At a
triple point the closed fibre of \eqref{cusp-cover-basechange-local-ring}
has length
\begin{align}
 \dim_{\mathbb C}
 \frac{\mathbb C\{w_0,w_1,w_2\}}
      {(w_0^N,w_1^N,w_2^N,w_0w_1w_2)}
 &=N^3-(N-1)^3\\
 &=3N^2-3N+1,
 \label{cusp-cover-triple-fibre-length}
\end{align}
which exceeds the generic rank $N^2$.  At a point of a double stratum,
localizing at the nonzero third coordinate and choosing its holomorphic
$N$th root reduces the ramified factor to
\[
 \frac{\mathbb C\{z_0,z_1,v\}}{(z_0z_1-v^N)}
\longrightarrow\mathbb C\{w_0,w_1\},\qquad
 z_i\longmapsto w_i^N,\quad v\longmapsto w_0w_1.
\]
Its local generic rank is $N$, whereas its closed fibre has length
\begin{equation}
 \dim_{\mathbb C}
 \frac{\mathbb C\{w_0,w_1\}}
      {(w_0^N,w_1^N,w_0w_1)}=2N-1>N.
 \label{cusp-cover-double-fibre-length}
\end{equation}
The remaining transverse nonzero torus coordinate is etale locally.
Equations \eqref{cusp-cover-triple-fibre-length} and
\eqref{cusp-cover-double-fibre-length}, together with etaleness elsewhere,
identify the nonflat locus exactly.

\medskip\noindent
\emph{The regular fibres and monodromy.}
Choose a logarithm $r$ with $v=e^{2\pi ir}$.  The target period matrix at
$t=v^N$ is
\[
 Z(Nr)=NrB_0+C(v^N),
\]
whereas the source residual action is determined by
\[
 Z_N(r)=rB_0+\frac1NC(v^N)=\frac1N Z(Nr).
\]
The source fibre is therefore
\[
 \frac{\mathbb C^2}{\mathbb Z^2+Z_N(r)\mathbb Z^2},
\]
and multiplication by $N$ on $\mathbb C^2$ maps it to the target fibre.
Equivalently, in the unscaled additive coordinate it is the identity map
\begin{equation}
 \frac{\mathbb C^2}
      {N\mathbb Z^2+Z(Nr)\mathbb Z^2}
 \longrightarrow
 \frac{\mathbb C^2}
      {\mathbb Z^2+Z(Nr)\mathbb Z^2}.
 \label{cusp-cover-additive-fibre-map}
\end{equation}
The imaginary part of $Z(Nr)$ is invertible on a regular fibre, so its
period lattice meets $\mathbb Z^2$ trivially in the required real sense;
the index in \eqref{cusp-cover-additive-fibre-map} is $N^2$.  Its deck
group is $(\mathbb Z/N\mathbb Z)^2$.  In the ordered lattice basis
$(\widehat\gamma,\widehat u,\widehat w,\widehat\delta)$, this is exactly
the sublattice obtained by $S_N$.  Since $M_0=I+R$ with $R^2=0$ and
$R(a,b,c,d)=(0,0,b,-a)$, one has
\[
 M_0^N=I+NR,\qquad S_N^{-1}M_0^NS_N=M_0,
\]
proving \eqref{cusp-cover-monodromy-conjugacy}.

\medskip\noindent
\emph{The central normalization and conductor.}
Multiplication by $N$ fixes every ray of $\Sigma$.  It therefore maps a
chosen source component of the toric central fibre to the corresponding
target component.  The star of that ray is the hexagonal fan, so the
restriction is the toric power map
$[N]\colon\mathrm{dP}_6\to\mathrm{dP}_6$.  On its dense two-torus it has
degree $N^2$; in each toric boundary chart it is an $N$th-power map in the
normal coordinate, hence has ramification index $N$ along all six boundary
curves.  Equation \eqref{cusp-cover-equivariance-computation} says that
the source and target lattice identifications commute with this map.
Consequently opposite boundary curves, their conductor identifications,
and the two translation-orbits of triangle vertices descend compatibly.
This proves \eqref{cusp-cover-normalization-diagram} and completes the
proof.
\end{proof}

\begin{remark}[Why the unmodified action does not work]
\label{cusp-cover-naive-action-failure}
If one reuses the target action with $v$ in place of $t$ at the source,
then the two sides of equivariance have multiplicative factors
\[
 c_{\bar\lambda}(v)^N v^{NB_0\bar\lambda}
 \quad\hbox{and}\quad
 c_{\bar\lambda}(v^N)v^{NB_0\bar\lambda},
\]
respectively.  In general
$c_{\bar\lambda}(v)^N\ne c_{\bar\lambda}(v^N)$.  Thus the bare formula
$(\xi_1,\xi_2,v)\mapsto(\xi_1^N,\xi_2^N,v^N)$ extends on the ambient toric
variety but does not descend through the manuscript's quotient unless the
source factor \eqref{cusp-cover-source-factor} is inserted.  The explicit
logarithm $2\pi iC(v^N)\bar\lambda/N$ makes this root holomorphic; no
unjustified choice of an $N$th root is involved.
\end{remark}

\begin{remark}[Exact scope]
\label{cusp-cover-global-scope}
The theorem constructs a finite morphism of completed families over the
local cusp discs.  It is a degree-$N^3$ map to the original family and a
degree-$N^2$ map only after the target is base changed.  It is neither a
degree-$N^2$ self-map over the unchanged disc nor the normalization of the
base change.

Nothing here supplies a cover over the entire base $B$.  In fact the
displayed sublattice is not preserved by the original full monodromy when
$N>1$: for example, the printed matrix $A_1$ sends
$\widehat u$ to a vector whose $\widehat w$- and
$\widehat\delta$-coefficients are $-1$ and $1$.  A further global base
cover and compatible extensions through both finite fillings would have
to be constructed separately.  The local Kummer morphism also supplies
no map to a product sphere, smash sphere, $S^{12}$, or $S^{24}$.
\end{remark}


\subsection{Cone-compactified sphere carriers for the native cusp quotient}
\label{cusp-cone-section}

This subsection answers a topological question which is strictly different
from the finite-cover construction.  No period is enlarged by passing to a
cover.  Instead, the native expanding quotient is retained on the cusp, the
parts on which that quotient is unavailable are collapsed, and the resulting
based map is passed through a smash product.  The outcome is an exact global
map with source \(S^{12}\), and there is an analogous \(S^{24}\) map from the
full-monodromy-invariant circle quotient.  The targets are conically
stratified spaces.  The Euclidean limits below use the intrinsic metrics
inherited from the flat quotient tori on the regular slices; they are not
ambient metric limits in one fixed metric
on the compact targets.

Fix a closed coordinate disc
\begin{equation}
 D_0=\{t\in\mathbb C:|t|\leq\epsilon\}\subset B,\qquad
 p_0=\{t=0\},\qquad N_0=f^{-1}(D_0),\qquad W=f^{-1}(0),
 \label{cusp-cone-disc}
\end{equation}
small enough that it contains no singular value other than \(p_0\).  Put
\begin{equation}
 K_\Lambda=\mathbb Z\langle\widehat w,\widehat\delta\rangle,\qquad
 \overline\Lambda=\Lambda/K_\Lambda,\qquad
 \overline T^2=
 \frac{\overline\Lambda\otimes\mathbb R}{\overline\Lambda}.
 \label{cusp-cone-quotient-lattice}
\end{equation}
The ordered basis of \(\overline\Lambda\) is the image of
\((\widehat\gamma,\widehat u)\).

\begin{lemma}[The native quotient over the whole punctured cusp]
\label{cusp-cone-native-descent}
There is a well-defined real-analytic fibrewise group epimorphism
\begin{equation}
 \overline\pi:N_0\setminus W
 \longrightarrow D_0^*\times\overline T^2
 \label{cusp-cone-native-family-map}
\end{equation}
whose kernel is the connected two-torus family defined by \(K_\Lambda\).
The target local system is trivial.  On the positive radial fibre \(F_T\),
\(t=e^{-2\pi T}\), the isomorphism
\begin{equation}
 \iota_T:\overline T^2\longrightarrow Q_T,\qquad
 [(a,b)]\longmapsto[Y_T(a,b)^{\mathsf T}]
 \label{cusp-cone-radial-isometry}
\end{equation}
satisfies
\begin{equation}
 \pi_T=\iota_T\circ\overline\pi|_{F_T}.
 \label{cusp-cone-native-restriction}
\end{equation}
If \(\overline T^2\) is given the pullback by \(\iota_T\) of the
Euclidean quotient metric on \(Q_T=\mathbb R^2/Y_T\mathbb Z^2\), then
\(\iota_T\) is an isometry and
\begin{equation}
 \operatorname{sys}(\overline T^2,g_T)
 =\operatorname{sys}(Q_T)\geq T-C.
 \label{cusp-cone-native-systole}
\end{equation}
\end{lemma}

\begin{proof}
On the universal cover of \(D_0^*\), use the real-analytic flat marking
\[
 (z,x+\Lambda)\longmapsto[z,\Pi(z)x].
\]
The deck generator acts by
\[
 (z,x+\Lambda)\longmapsto(g_0z,M_0x+\Lambda).
\]
The displayed matrix \(M_0\) gives
\begin{equation}
 (M_0-I)(a,b,c,d)^{\mathsf T}=(0,0,b,-a)^{\mathsf T},\qquad
 \operatorname{im}(M_0-I)=K_\Lambda.
 \label{cusp-cone-monodromy-image}
\end{equation}
Consequently \(M_0\) preserves \(K_\Lambda\) and induces the identity on
\(\overline\Lambda\).  The formula
\[
 [z,\Pi(z)x]\longmapsto
 \bigl([z],[x\bmod K_\Lambda]\bigr)
\]
is invariant under the deck action and proves
\cref{cusp-cone-native-family-map}.  The lattice \(K_\Lambda\) is primitive,
so the kernel is connected.  On the radial fibre, projection to the
\((e_2,e_4)\)-coordinates sends
\(a\widehat\gamma+b\widehat u\) to \(Y_T(a,b)^{\mathsf T}\) and kills
\(K_\Lambda\).  This proves
\cref{cusp-cone-radial-isometry,cusp-cone-native-restriction}.  The metric
and systolic assertion is the quotient calculation in
\cref{cusp-spec-native-quotient}.
\end{proof}

Define
\begin{equation}
 \mathcal A=
 \frac{[0,\epsilon]\times S^1\times\overline T^2}
 {\{0\}\times S^1\times\overline T^2\sim v_0,\quad
  \{\epsilon\}\times S^1\times\overline T^2\sim v_1}.
 \label{cusp-cone-A-definition}
\end{equation}
Thus
\begin{equation}
 \mathcal A\cong
 \operatorname{Susp}(S^1\times\overline T^2)
 \cong\operatorname{Susp}(T^3),
 \label{cusp-cone-A-suspension}
\end{equation}
where \(\operatorname{Susp}\) denotes the unreduced suspension.  The extra
circle is the angular coordinate of the punctured base.  If one deliberately
forgets that angle, the evident map
\(\mathcal A\to\operatorname{Susp}(\overline T^2)\) gives a smaller target
with the same restriction to every fixed radial fibre.  We retain
\(\mathcal A\) because it records the entire punctured-cusp quotient.

\begin{theorem}[Exact localized \(S^{12}\) route]
\label{cusp-cone-S12-route}
Choose \(x_0\in W\), and base \(\mathcal A\) at \(v_0\).  There is a
continuous surjective based quotient map
\begin{equation}
 q_2:(X,x_0)\longrightarrow(\mathcal A,v_0)
 \label{cusp-cone-q2}
\end{equation}
with the following exact description.  For \(x\in N_0\setminus W\), write
\[
 t(f(x))=re^{i\phi},\qquad 0<r<\epsilon,
\]
and let \(\overline\pi_f(x)\in\overline T^2\) be the fibre component of
\(\overline\pi(x)\).  Then
\begin{equation}
 q_2(x)=[r,e^{i\phi},\overline\pi_f(x)],\qquad
 q_2(W)=v_0,\qquad
 q_2\bigl(X\setminus\operatorname{int}N_0\bigr)=v_1.
 \label{cusp-cone-q2-formula}
\end{equation}
On \(f^{-1}(|t|=\epsilon)\), the value is \(v_1\).

The smash square is a global based map
\begin{equation}
 q_2^{\wedge2}:X\wedge X\longrightarrow
 \mathcal A\wedge\mathcal A.
 \label{cusp-cone-S12-map}
\end{equation}
Using \(X\cong_{\mathrm{diff}}S^6\), a based choice of the diffeomorphism
identifies its source with
\begin{equation}
 X\wedge X\cong S^6\wedge S^6\cong S^{12}.
 \label{cusp-cone-S12-source}
\end{equation}
For \(0<e^{-2\pi T}<\epsilon\), the canonical product-to-smash quotient
maps restrict to embeddings on \(F_T^2\) and on the corresponding interior
target-product cross-section.  Under those embeddings and \(\iota_T^2\),
the induced smash map restricts exactly to
\begin{equation}
 \pi_T\times\pi_T:F_T\times F_T\longrightarrow Q_T\times Q_T.
 \label{cusp-cone-S12-fibre-map}
\end{equation}
With the Euclidean \(\ell^2\)-product flat quotient metric,
\begin{equation}
 \operatorname{sys}(Q_T\times Q_T)\geq T-C,\qquad
 (Q_T\times Q_T,(0,0))\longrightarrow(\mathbb R^4,0)
 \label{cusp-cone-S12-R4-limit}
\end{equation}
in pointed local-isometric, pointed smooth, and pointed
Gromov--Hausdorff topology.  Once \(T-C>4R\), the radius-\(R\) ball is
exactly Euclidean.

The integral homology of the targets is
\begin{equation}
 \widetilde H_k(\mathcal A;\mathbb Z)\cong
 \begin{cases}
  \mathbb Z^3,&k=2,3,\\
  \mathbb Z,&k=4,\\
  0,&\text{otherwise},
 \end{cases}
 \label{cusp-cone-A-homology}
\end{equation}
and
\begin{equation}
 \widetilde H_k(\mathcal A\wedge\mathcal A;\mathbb Z)\cong
 \begin{cases}
  \mathbb Z^9,&k=4,\\
  \mathbb Z^{18},&k=5,\\
  \mathbb Z^{15},&k=6,\\
  \mathbb Z^6,&k=7,\\
  \mathbb Z,&k=8,\\
  0,&\text{otherwise}.
 \end{cases}
 \label{cusp-cone-A-smash-homology}
\end{equation}
In particular, the target is not \(S^4\), \(S^{12}\), or a compact
four-manifold.
\end{theorem}

\begin{proof}
The equivalence relation in \cref{cusp-cone-A-definition} is closed: it is
the union of the diagonal with the squares of the two closed end
cross-sections.  A compact Hausdorff space modulo a closed equivalence
relation is compact Hausdorff.  Rescaling \(r\) to the unit interval gives
exactly the unreduced suspension in \cref{cusp-cone-A-suspension}.

Well-definedness of the interior formula follows from
\cref{cusp-cone-native-descent}.  A full turn in \(\phi\) changes the
marking by \(M_0\), but \(M_0\) is the identity modulo \(K_\Lambda\).

We check both collapsed strata.  If \(U\) is an open neighborhood of
\(v_0\), its inverse image in the cylinder is open and contains the compact
set \(\{0\}\times S^1\times\overline T^2\).  Hence, for some \(\delta>0\),
it contains
\([0,\delta)\times S^1\times\overline T^2\).  Therefore
\(q_2^{-1}(U)\) contains \(f^{-1}(|t|<\delta)\), a neighborhood of \(W\).
The same compactness argument at the other end shows that every neighborhood
of \(v_1\) contains the image of
\((\epsilon-\delta,\epsilon]\times S^1\times\overline T^2\).  Its preimage
together with \(X\setminus\operatorname{int}N_0\) is a neighborhood of the
outer boundary.  Continuity on the interior is
\cref{cusp-cone-native-family-map}; the closed pasting lemma gives
continuity on \(X\).

Every interior point has a preimage because \(\overline\pi\) is fibrewise
surjective.  The nonempty sets \(W\) and
\(X\setminus\operatorname{int}N_0\) map onto \(v_0\) and \(v_1\),
respectively.  Thus \(q_2\) is surjective.  A continuous surjection from
compact \(X\) to Hausdorff \(\mathcal A\) is a quotient map.

The map is based, so functoriality gives \cref{cusp-cone-S12-map}.  The
standard CW quotient gives \(S^m\wedge S^n\cong S^{m+n}\)
\cite[Chapter~0, Smash Product]{Hatcher2002}, proving
\cref{cusp-cone-S12-source}.

The base point lies in \(W\), so \(F_T^2\) is disjoint from the source fat
wedge.  Its image lies in an interior cross-section of
\(\mathcal A^2\), disjoint from the target fat wedge.  The source and target
fat wedges are closed in the compact Hausdorff products \(X^2\) and
\(\mathcal A^2\), so both smash quotients are Hausdorff.  Their restrictions
to the relevant compact subspaces are continuous and one-to-one, and are
therefore embeddings.
Equations \cref{cusp-cone-native-restriction,cusp-cone-q2-formula} give
\cref{cusp-cone-S12-fibre-map}.

The period lattice of the product target is
\(\operatorname{diag}(Y_T,Y_T)\mathbb Z^4\).  Its systole is the systole of
\(Y_T\mathbb Z^2\), proving the first assertion in
\cref{cusp-cone-S12-R4-limit}.  If \(x,y\in B_R(0)\subset\mathbb R^4\)
and \(0\ne\lambda\) is a product period, then
\[
 \|x-y+\lambda\|\geq\|\lambda\|-\|x-y\|
 \geq(T-C)-2R>2R\geq\|x-y\|
\]
when \(T-C>4R\).  The quotient projection is an exact isometry on the
Euclidean radius-\(R\) ball.  Conversely, any point of quotient distance
less than \(R\) has a shortest lattice representative of norm less than
\(R\), so this embedded ball is the entire quotient radius-\(R\) ball.
This proves all three pointed convergences.

Reduced suspension shifts reduced homology by one.  Since
\[
 H_1(T^3)=\mathbb Z^3,\qquad
 H_2(T^3)=\mathbb Z^3,\qquad H_3(T^3)=\mathbb Z,
\]
\cref{cusp-cone-A-homology} follows.  All groups are free, so the Kunneth
theorem for smash products has no Tor terms.  The ranks in degrees
\(4,5,6,7,8\) are
\[
 3\cdot3,\quad 2(3\cdot3),\quad
 2(3\cdot1)+3\cdot3,\quad 2(3\cdot1),\quad1,
\]
which proves \cref{cusp-cone-A-smash-homology}.
\end{proof}

\paragraph{The localized one-circle construction.}
The full-monodromy-invariant character also gives a smaller cusp-supported
target.  Define
\begin{equation}
 \mathcal B=
 \frac{[0,\epsilon]\times S^1\times S^1}
 {\{0\}\times S^1\times S^1\sim b_0,\quad
  \{\epsilon\}\times S^1\times S^1\sim b_1}
 \cong\operatorname{Susp}(T^2),
 \label{cusp-cone-B-definition}
\end{equation}
where the two circle coordinates in an interior cross-section are the base
angle and the value of \(\theta_\gamma\).

\begin{proposition}[Localized \(S^{24}\) route with target
\(\operatorname{Susp}(T^2)^{\wedge4}\)]
\label{cusp-cone-S24-local-route}
There is a continuous surjective based quotient map
\begin{equation}
 q_{1,\mathrm{cusp}}:(X,x_0)\longrightarrow(\mathcal B,b_0)
 \label{cusp-cone-q1-local}
\end{equation}
given by
\begin{equation}
 q_{1,\mathrm{cusp}}(x)
 =[r,e^{i\phi},\theta_\gamma(x)]
 \quad(0<|t|=r<\epsilon),\qquad
 q_{1,\mathrm{cusp}}(W)=b_0,\qquad
 q_{1,\mathrm{cusp}}
 \bigl(X\setminus\operatorname{int}N_0\bigr)=b_1.
 \label{cusp-cone-q1-local-formula}
\end{equation}
Its fourth smash power has source \(X^{\wedge4}\cong S^{24}\), and on
four synchronized radial fibres it is exactly
\begin{equation}
 \chi_{\gamma,T}^{\times4}:F_T^4\longrightarrow C_T^4.
 \label{cusp-cone-S24-local-fibre-map}
\end{equation}
Consequently its intrinsic target slices satisfy the Euclidean limit
\cref{cusp-cone-S24-R4-limit}.  The target homology is
\begin{equation}
 \widetilde H_k(\mathcal B;\mathbb Z)\cong
 \begin{cases}
  \mathbb Z^2,&k=2,\\
  \mathbb Z,&k=3,\\
  0,&\text{otherwise},
 \end{cases}
 \label{cusp-cone-B-homology}
\end{equation}
and
\begin{equation}
 \widetilde H_k(\mathcal B^{\wedge4};\mathbb Z)\cong
 \begin{cases}
  \mathbb Z^{16},&k=8,\\
  \mathbb Z^{32},&k=9,\\
  \mathbb Z^{24},&k=10,\\
  \mathbb Z^8,&k=11,\\
  \mathbb Z,&k=12,\\
  0,&\text{otherwise}.
 \end{cases}
 \label{cusp-cone-B-smash-homology}
\end{equation}
\end{proposition}

\begin{proof}
The monodromy invariance
\(\varepsilon_\gamma M_0=\varepsilon_\gamma\) makes the interior formula
single-valued after a turn in the base angle.  The two endpoint-continuity
arguments are identical to those in the proof of
\cref{cusp-cone-S12-route}, with compact cross-section \(T^2\) in place of
\(T^3\).  Fibrewise surjectivity of \(\theta_\gamma\), together with the
nonempty inner fibre and exterior, proves surjectivity.  Compact-to-Hausdorff
makes the map a quotient map.

Since \(x_0\in W\), the fourth smash power exists.  The radial regular
fibres and their circle targets avoid the two closed fat wedges.  The
restrictions of the Hausdorff smash quotients to these compact subspaces are
continuous and one-to-one, hence embeddings, and the fibre map is exactly
\cref{cusp-cone-S24-local-fibre-map}.  The intrinsic metric statement is
therefore the four-circle calculation
\cref{cusp-cone-S24-R4-limit}.  Finally, suspension shifts the free homology
of \(T^2\), giving \cref{cusp-cone-B-homology}.  In the fourth smash power,
choosing \(j\) of the four factors in degree \(3\) and the other
\(4-j\) factors in degree \(2\) gives degree \(8+j\) and rank
\(\binom4j2^{4-j}\).  This is \cref{cusp-cone-B-smash-homology}.
\end{proof}

\paragraph{The full-base circle target before coning.}
Choose pairwise disjoint closed coordinate discs \(D_0,D_1,D_2\) about the
three singular values and put
\begin{equation}
 P=B\setminus\bigcup_{j=0}^2\operatorname{int}D_j.
 \label{cusp-cone-pair-of-pants}
\end{equation}
For \(j=1,2\), let \(C_j\) be the solid-torus Seifert target in
\cref{cusp-spec-finite-Seifert-target}, and let \(h_j^{S^1}\) be the
target overlap map in \cref{cusp-spec-finite-overlap-identity}.  Glue
\begin{equation}
 \mathcal M=(P\times S^1)
 \underset{h_1^{S^1}}{\cup}C_1
 \underset{h_2^{S^1}}{\cup}C_2.
 \label{cusp-cone-M-definition}
\end{equation}
The gluing is along boundary tori over \(\partial D_1,\partial D_2\).
The remaining boundary is
\begin{equation}
 \partial\mathcal M=\partial D_0\times S^1\cong T^2.
 \label{cusp-cone-M-boundary}
\end{equation}
Define
\begin{equation}
 \mathcal C=\mathcal M\cup_{\partial\mathcal M}
 \operatorname{Cone}(\partial\mathcal M),
 \label{cusp-cone-C-definition}
\end{equation}
and write \(c_0\) for the cone vertex.
This changes the codomain at the cusp: the whole boundary two-torus,
including the target circle, is collapsed to \(c_0\).  It therefore does
not contradict the proved nonextension of the literal map
\(\chi_\gamma\) to the fixed product target \(B\times S^1\).

\begin{theorem}[Exact cone-compactified \(S^{24}\) route]
\label{cusp-cone-S24-route}
There is a continuous surjective based quotient map
\begin{equation}
 q_1:(X,x_0)\longrightarrow(\mathcal C,c_0),\qquad x_0\in W,
 \label{cusp-cone-q1}
\end{equation}
defined as follows:
\begin{enumerate}
\item over \(P\), it is
\((f,\theta_\gamma):f^{-1}(P)\to P\times S^1\), equivalently the
restriction of \(\chi_\gamma\) to \(f^{-1}(P)\);
\item over \(D_j\), \(j=1,2\), it is the finite-filling map
\(N_j\to C_j\) from \cref{cusp-spec-finite-Seifert-target};
\item on \(N_0\setminus W\), writing \(t=re^{i\phi}\), it is
\begin{equation}
 x\longmapsto[r,(e^{i\phi},\theta_\gamma(x))]
 \in\operatorname{Cone}(\partial D_0\times S^1),
 \label{cusp-cone-q1-cusp-formula}
\end{equation}
and every point of \(W\) maps to \(c_0\).
\end{enumerate}
The fourth smash power is a global based map
\begin{equation}
 q_1^{\wedge4}:X^{\wedge4}\longrightarrow\mathcal C^{\wedge4}.
 \label{cusp-cone-S24-map}
\end{equation}
Using \(X\cong_{\mathrm{diff}}S^6\), its source is
\begin{equation}
 X^{\wedge4}\cong(S^6)^{\wedge4}\cong S^{24}.
 \label{cusp-cone-S24-source}
\end{equation}
On four synchronized radial fibres, the canonical product-to-smash quotient
maps restrict to embeddings on the source fibre product and the interior
target-product cross-section.  Under those embeddings,
\cref{cusp-cone-S24-map} restricts exactly to
\begin{equation}
 \chi_{\gamma,T}^{\times4}:F_T^4\longrightarrow C_T^4,\qquad
 C_T=\mathbb R/\ell_T\mathbb Z.
 \label{cusp-cone-S24-fibre-map}
\end{equation}
With its Euclidean \(\ell^2\)-product flat quotient metric,
\begin{equation}
 \operatorname{sys}(C_T^4)=\ell_T
 =T+H-B+O(T^{-1}),\qquad
 (C_T^4,0)\longrightarrow(\mathbb R^4,0)
 \label{cusp-cone-S24-R4-limit}
\end{equation}
in pointed local-isometric, pointed smooth, and pointed
Gromov--Hausdorff topology.

The target homology is
\begin{equation}
 \widetilde H_k(\mathcal C;\mathbb Z)\cong
 \begin{cases}
  \mathbb Z,&k=2,3,\\
  0,&\text{otherwise},
 \end{cases}
 \label{cusp-cone-C-homology}
\end{equation}
and
\begin{equation}
 \widetilde H_k(\mathcal C^{\wedge4};\mathbb Z)\cong
 \begin{cases}
  \mathbb Z,&k=8,12,\\
  \mathbb Z^4,&k=9,11,\\
  \mathbb Z^6,&k=10,\\
  0,&\text{otherwise}.
 \end{cases}
 \label{cusp-cone-C-smash-homology}
\end{equation}
\end{theorem}

\begin{proof}
The compact spaces \(P\times S^1,C_1,C_2\) are glued along closed boundary
tori by homeomorphisms, producing a compact orientable three-manifold
\(\mathcal M\) with boundary \(\partial D_0\times S^1\).  Attaching the cone
along this closed boundary gives a compact Hausdorff finite CW space
\(\mathcal C\).

The overlap identity
\[
 \chi_\gamma\circ h_j=h_j^{S^1}\circ\overline\chi_j
\]
proves that the first two pieces of \(q_1\) agree.  On \(\partial D_0\), the
cone base is glued to the remaining boundary of \(P\times S^1\), so the
first and third pieces agree.  If \(U\) is a neighborhood of \(c_0\), its
inverse image in the cone contains \([0,\delta)\times T^2\) for some
\(\delta>0\), by compactness of the cone base.  Therefore \(q_1^{-1}(U)\)
contains \(f^{-1}(|t|<\delta)\), a neighborhood of \(W\).  This proves
continuity at the new stratum; continuity elsewhere follows from the input
maps and the closed pasting lemma.

Every regular circle quotient and both finite Seifert targets are
surjective, and \(W\) maps onto the cone vertex.  Thus \(q_1\) is a
continuous surjection from compact \(X\) to Hausdorff \(\mathcal C\), hence
a quotient map.  It is based, so its fourth smash power exists.  Iterating
\(S^m\wedge S^n\cong S^{m+n}\) gives \cref{cusp-cone-S24-source}.

The base point lies in \(W\), while each radial regular fibre is disjoint
from \(W\).  Hence \(F_T^4\) and its product-circle image avoid the closed
source and target fat wedges.  The canonical quotient maps from the source
product and the indicated interior target-product cross-section to their
smash products restrict to embeddings.  Under those embeddings, the induced
smash map restricts to \(\chi_{\gamma,T}^{\times4}\), which itself retains
its displayed nontrivial torus kernel.
The quotient-metric calculation in
\cref{cusp-spec-global-circle-circumference} gives \(\ell_T\).
The product lattice is \(\ell_T\mathbb Z^4\), so the radius-\(R\) argument,
including the shortest-representative proof that the embedded Euclidean ball
is the whole quotient ball, in the proof of \cref{cusp-cone-S12-route} proves
\cref{cusp-cone-S24-R4-limit}.

For homology, let \(\rho_0,\rho_1,\rho_2\) be the boundary loops of the pair
of pants, with \(\rho_0+\rho_1+\rho_2=0\), and let \(\eta\) be the regular
target-circle class.  With the boundary-loop and circle orientations used
in the finite-filling presentation, the meridian relations are
\begin{equation}
 3\rho_1-\eta=0,\qquad4\rho_2+\eta=0.
 \label{cusp-cone-Seifert-meridians}
\end{equation}
Reversing the circle orientation reverses both displayed \(\eta\)-signs
and changes none of the following groups.  Thus
\begin{equation}
 H_1(\mathcal M;\mathbb Z)
 =\frac{\mathbb Z\langle\eta,\rho_1,\rho_2\rangle}
 {\langle3\rho_1-\eta,\;4\rho_2+\eta\rangle}
 \cong\mathbb Z\langle g\rangle,
 \label{cusp-cone-M-H1}
\end{equation}
where
\begin{equation}
 \rho_0=-g,\qquad\rho_1=4g,\qquad
 \rho_2=-3g,\qquad\eta=12g.
 \label{cusp-cone-M-boundary-map}
\end{equation}
In particular,
\(H_1(\partial\mathcal M)=\mathbb Z\langle\rho_0,\eta\rangle
\to H_1(\mathcal M)\) is surjective.  The orientable three-manifold
\(\mathcal M\) has one torus boundary, \(H_3(\mathcal M)=0\), and
\(\chi(\mathcal M)=\frac12\chi(\partial\mathcal M)=0\); hence
\(H_2(\mathcal M;\mathbb Q)=0\).  To obtain the integral statement without
suppressing torsion, Poincare--Lefschetz duality and the universal
coefficient theorem give
\[
 H_2(\mathcal M;\mathbb Z)
 \cong H^1(\mathcal M,\partial\mathcal M;\mathbb Z)
 \cong \operatorname{Hom}
 \bigl(H_1(\mathcal M,\partial\mathcal M;\mathbb Z),\mathbb Z\bigr),
\]
because \(H_0(\mathcal M,\partial\mathcal M;\mathbb Z)=0\).  Thus
\(H_2(\mathcal M;\mathbb Z)\) is free; its rank is zero, so it vanishes.

The cone attachment is the homotopy cofiber of
\(\partial\mathcal M\hookrightarrow\mathcal M\).  Since the boundary
inclusion is a cofibration,
\[
 \widetilde H_k(\mathcal C;\mathbb Z)
 \cong H_k(\mathcal M,\partial\mathcal M;\mathbb Z).
\]
Poincare--Lefschetz duality and
\(H^0(\mathcal M)=H^1(\mathcal M)=\mathbb Z\),
\(H^2(\mathcal M)=H^3(\mathcal M)=0\), give
\cref{cusp-cone-C-homology}.  The homology is free, so the smash Kunneth
theorem says that four factors of degree \(2\) or \(3\) produce ranks
\(\binom4j\) in degree \(8+j\), proving
\cref{cusp-cone-C-smash-homology}.
\end{proof}

\begin{definition}[The three radial target-slice embeddings]
\label{cusp-cone-target-slice-embeddings}
Put \(r_T=e^{-2\pi T}\) and let
\[
 \vartheta_T:\mathbb R/\mathbb Z\longrightarrow
 C_T=\mathbb R/\ell_T\mathbb Z,\qquad [a]\longmapsto[\ell_Ta].
\]
For \(y_1,y_2\in Q_T\), define
\[
 j_T(y_1,y_2)=
 \bigl[[r_T,1,\iota_T^{-1}(y_1)],
       [r_T,1,\iota_T^{-1}(y_2)]\bigr]_{\wedge}
 \in\mathcal A^{\wedge2}.
\]
For \(c_1,\ldots,c_4\in C_T\), define
\[
 k_T^{\mathcal B}(c_1,\ldots,c_4)=
 \bigl[[r_T,1,\vartheta_T^{-1}(c_1)],\ldots,
       [r_T,1,\vartheta_T^{-1}(c_4)]\bigr]_{\wedge}
 \in\mathcal B^{\wedge4},
\]
and, using the cusp-cone piece of \(\mathcal C\),
\[
 k_T^{\mathcal C}(c_1,\ldots,c_4)=
 \bigwedge_{j=1}^4[r_T,(1,\vartheta_T^{-1}(c_j))]
 \in\mathcal C^{\wedge4}.
\]
These are embeddings because every displayed factor lies away from the smash
basepoint, so equality in the smash quotient is equality in every factor;
compact-to-Hausdorff injectivity then gives a topological embedding.
\end{definition}

\begin{proposition}[Intrinsic expansion versus ambient collapse]
\label{cusp-cone-metric-scope}
Let \(D_{12}\), \(D_{24}^{\mathcal B}\), and
\(D_{24}^{\mathcal C}\) be arbitrary fixed metrics compatible with the
topologies of \(\mathcal A^{\wedge2}\), \(\mathcal B^{\wedge4}\), and
\(\mathcal C^{\wedge4}\), respectively.  In their corresponding ambient
metrics, the radial target
slices satisfy
\begin{equation}
 \operatorname{diam}_{\mathrm{amb}}
 \bigl(j_T(Q_T^2)\bigr)\longrightarrow0,\qquad
 \operatorname{diam}_{\mathrm{amb}}
 \bigl(k_T^{\mathcal B}(C_T^4)\bigr)\longrightarrow0,\qquad
 \operatorname{diam}_{\mathrm{amb}}
 \bigl(k_T^{\mathcal C}(C_T^4)\bigr)\longrightarrow0.
 \label{cusp-cone-ambient-collapse}
\end{equation}
By contrast, in the intrinsic quotient metrics,
\begin{equation}
 \operatorname{diam}(Q_T^2)\geq\tfrac12(T-C)\longrightarrow\infty,\qquad
 \operatorname{diam}(C_T^4)\geq\tfrac12\ell_T\longrightarrow\infty.
 \label{cusp-cone-intrinsic-expansion}
\end{equation}
Thus \cref{cusp-cone-S12-R4-limit,cusp-cone-S24-R4-limit} assert
intrinsic, flat pointed limits of varying slices, not Sormani--Wenger
intrinsic-flat convergence and not an ambient pointed \(\mathbb R^4\) limit
inside one fixed compact metric target.
\end{proposition}

\begin{proof}
Every neighborhood of the relevant cone vertex contains every sufficiently
small radial cross-section, uniformly in all compact angular and torus
coordinates.  In a compatible metric, every such cross-section eventually
lies in an arbitrarily small ball about the vertex.  Products followed by
the smash quotient preserve this uniform convergence, proving
\cref{cusp-cone-ambient-collapse}.

For a flat torus with systole \(s\), choose a shortest primitive period
\(\lambda\).  The class of \(\lambda/2\) has distance \(s/2\) from zero:
for every period \(\mu\), the nonzero period \(\lambda-2\mu\) has length at
least \(s\).  Hence the diameter is at least \(s/2\).
Equations \cref{cusp-cone-native-systole,cusp-cone-S24-R4-limit} prove
\cref{cusp-cone-intrinsic-expansion}.
\end{proof}

\begin{proposition}[Exact almost-complex obstruction on \(S^{12}\)]
\label{cusp-cone-S12-almost-complex-obstruction}
Let \(M\) be a smooth manifold homeomorphic to \(S^{12}\).  There is no
bundle endomorphism \(J:TM\to TM\) satisfying \(J^2=-\operatorname{Id}_{TM}\).
This conclusion has two independent proofs which use only the homotopy type
of \(M\): the original Borel--Serre reduced-power calculation at the single
prime \(5\), and Bott integrality.  For the standard embedded sphere there is
also an independent stable-tangent Pontryagin-class contradiction.

These statements concern an unstabilized complex structure on the actual
rank-twelve tangent bundle.  A stable complex structure on
\(TS^{12}\oplus\varepsilon^r_{\mathbb R}\), a complex structure confined to
an open subset or to fibres of a map, or an endomorphism which ceases to be a
bundle endomorphism along a singular set is not an almost-complex structure
on \(TS^{12}\) and therefore lies outside the hypothesis.
\end{proposition}

\begin{proof}
Assume that \(J\) exists and put \(E=(TM,J)\), a complex vector bundle of
rank six.  Orient \(M\) by the interleaved complex frame
\[
 (v_1,Jv_1,\ldots,v_6,Jv_6),
\]
write \([M]_J\) for the resulting fundamental class, and choose
\(u_J\in H^{12}(M;\mathbb Z)\) by
\(\langle u_J,[M]_J\rangle=1\).  Since \(M\) has the integral cohomology of
\(S^{12}\),
\begin{equation}
 c_1(E)=c_2(E)=c_3(E)=c_4(E)=c_5(E)=0.
 \label{cusp-cone-S12-lower-Chern-zero}
\end{equation}
The top-Chern/Euler identity uses exactly the orientation just fixed, so
Poincare--Hopf gives
\begin{equation}
 c_6(E)=e(E_{\mathbb R})=e(TM)=2u_J.
 \label{cusp-cone-S12-top-Chern-two}
\end{equation}
Thus reversing a previously chosen orientation can change the displayed
generator but cannot change the coefficient two.

We first specialize the original calculation of Borel and Serre
\cite[\S12, Proposition~12.4, Lemma~12.5, and Proposition~12.6,
p.~430; \S15, Proposition~15.1, p.~434]{BorelSerre1953}.
They write \(C_{2j}\) for the universal class now denoted \(c_j\).
For \(j=6\), the admissible prime is exactly
\[
 p=5<6,\qquad 5\nmid6,
\]
and their identity in \(H^*(BU(6);\mathbb F_5)\) reads
\begin{equation}
 \overline c_6
 =6^{-1}\mathcal P_5^1(\overline c_2)
   +Q(\overline c_1,\ldots,\overline c_5).
 \label{cusp-cone-S12-Borel-Serre-specialization}
\end{equation}
Here \(6^{-1}=1\) in \(\mathbb F_5\).  To retain the exact coefficient rather
than suppressing the lower-class polynomial, pass to the splitting space,
write \(x_1,\ldots,x_6\) for the degree-two Chern roots, put
\(s_r=\sum_i x_i^r\), and use
\(\mathcal P_5^1(x_i)=x_i^5\) and the Cartan formula.  The resulting universal
identity is
\begin{align}
 \mathcal P_5^1(\overline c_2)
 ={}&\sum_{i<j}(x_i^5x_j+x_ix_j^5)=s_5s_1-s_6 \notag\\
 ={}&6\overline c_6+\overline c_1^4\overline c_2
 -\overline c_1^3\overline c_3
 -4\overline c_1^2\overline c_2^2
 +\overline c_1^2\overline c_4
 +7\overline c_1\overline c_2\overline c_3
 -\overline c_1\overline c_5 \notag\\
 &\quad+2\overline c_2^3-6\overline c_2\overline c_4
 -3\overline c_3^2.
 \label{cusp-cone-S12-P1-exact}
\end{align}
Consequently, for the lower-Chern ideal
\(I=(\overline c_1,\ldots,\overline c_5)\),
\begin{equation}
 \mathcal P_5^1(\overline c_2)
 \equiv6\overline c_6=\overline c_6\pmod I,
 \label{cusp-cone-S12-P1-lower-ideal}
\end{equation}
which verifies the exact top-class coefficient in
\cref{cusp-cone-S12-Borel-Serre-specialization}.  Pulling back along the
classifying map of \(E\), using \cref{cusp-cone-S12-lower-Chern-zero} and
naturality of \(\mathcal P_5^1\), gives
\[
 \overline{c_6(E)}=\mathcal P_5^1(\overline{c_2(E)})
 =\mathcal P_5^1(0)=0.
\]
But \cref{cusp-cone-S12-top-Chern-two} gives
\(\overline{c_6(E)}=2\overline u_J\ne0\) in
\(H^{12}(M;\mathbb F_5)\), the first contradiction.

For the second contradiction, choose an orientation-preserving homotopy
equivalence \(f:S^{12}\to M\).  Bott integrality says that the sixth Chern
class of every complex bundle on \(S^{12}\) is divisible by
\((6-1)!=120\)
\cite[Chapter~20, \S9, Theorem~9.6 and Corollary~9.8,
p.~308]{Husemoller1994}; see also
\cite[\S6]{KonstantisParton2018}.  The coefficient can be checked without
normalization ambiguity: Newton's sixth identity and
\cref{cusp-cone-S12-lower-Chern-zero} give
\[
 \operatorname{ch}_6(f^*E)
 =\frac{s_6(f^*E)}{6!}
 =-\frac{c_6(f^*E)}{5!}.
\]
The integral Bott generator therefore forces
\(c_6(f^*E)=120k f^*u_J\) for some \(k\in\mathbb Z\), whereas
\cref{cusp-cone-S12-top-Chern-two} gives
\(c_6(f^*E)=2f^*u_J\).  Thus \(120k=2\), which has no integral solution.

Finally take \(M=S^{12}\) with its standard smooth embedding in
\(\mathbb R^{13}\).  There is an explicit stable trivialization
\begin{equation}
 TS^{12}\oplus\varepsilon^1_{\mathbb R}
 \xrightarrow{\ \cong\ }\varepsilon^{13}_{\mathbb R},
 \qquad (x;v,a)\longmapsto(x;v+ax),
 \label{cusp-cone-S12-stable-trivialization}
\end{equation}
so \(p_3(TS^{12})=0\).  On the other hand,
\[
 E_{\mathbb R}\otimes_{\mathbb R}\mathbb C
 \cong E\oplus\overline E,
 \qquad
 p_3(E_{\mathbb R})=-c_6(E\oplus\overline E).
\]
Equation \cref{cusp-cone-S12-lower-Chern-zero} and
\(c_6(\overline E)=c_6(E)\) give the exact relation
\begin{equation}
 p_3(TS^{12})=-2c_6(E)=-4u_J,
 \label{cusp-cone-S12-Pontryagin-contradiction}
\end{equation}
contradicting stable triviality in the torsion-free group
\(H^{12}(S^{12};\mathbb Z)\).  This is the Pontryagin--Euler proof recorded
independently in
\cite[Lemma~4.1 and Theorem~4.2]{KonstantisParton2018}.

Every contradiction above requires a globally defined \(J\) on the actual
bundle \(TM\).  Stabilizing that bundle, deleting points, restricting to
regular fibres, or allowing \(J\) to fail at a singular locus removes that
hypothesis rather than producing an exception to it.
\end{proof}

\begin{theorem}[Sharp obstructions to stronger interpretations]
\label{cusp-cone-obstructions}
The sphere routes above have the following exact limitations.
\begin{enumerate}
\item Let \(r\geq1\), and let \(N\subset\Lambda^{\oplus r}\) be a
primitive kernel defining a connected fibrewise torus quotient over the
full regular base, or over any finite cover of that base.  Equip the quotient
with the norm induced by the unrescaled marked Euclidean metric on the
original fibres.  If its quotient systole tends to infinity on the radial
cusp, then
\begin{equation}
 \operatorname{rank}(\Lambda^{\oplus r}/N)\leq r.
 \label{cusp-cone-rank-bound}
\end{equation}
Thus two copies cannot give a full-base four-torus quotient.  The
\(S^{12}\) map in \cref{cusp-cone-S12-route} evades the obstruction only by
localizing to the cusp and collapsing the exterior.

\item No nontrivial finite base cover, including one extending with
branching over the original singular values, makes \(K_\Lambda\) a
sub-local system over the full regular base.  The local cyclic change
\(t=s^M\) used for the auxiliary degree-\(M^2\) cover cannot extend as a
connected cover of \(B\cong\mathbb P^1\) branched only at \(p_0\).

\item There is no map from \(X^2\), \(X\wedge X\), \(X^4\), or
\(X^{\wedge4}\) to a fixed \(T^4\) whose restriction retains all four
target character classes of the displayed product quotient.  There is no
map to a fixed \(S^4\) whose restriction is a degree-one collapse of
\(Q_T^2\), or of the four-circle target, for even one regular fibre.  The
conical targets evade these obstructions because the fibre character and
orientation classes do not extend as global target classes.

\item Cartesian products do not produce the named spheres:
\begin{equation}
 X^2\simeq S^6\times S^6\not\simeq S^{12},\qquad
 X^4\simeq(S^6)^4\not\simeq S^{24}.
 \label{cusp-cone-products-not-spheres}
\end{equation}
The two- and four-fold synchronized fibre products have real dimensions
\(10\) and \(18\), respectively, and generally have singular completed
fibres.  Joins give
\begin{equation}
 X*X\cong S^{13},\qquad X^{*4}\cong S^{27}.
 \label{cusp-cone-join-dimensions}
\end{equation}

\item Since \(X\cong S^6\) is simply connected, every connected unbranched
finite cover of \(X\), \(X^r\), \(X\wedge X\cong S^{12}\), or
\(X^{\wedge4}\cong S^{24}\) is trivial.  The local degree-\(M^2\)
unspooling cover is not the restriction of a nontrivial connected
unbranched cover of any of these compact sources.  Branched or normalized
compactifications are different objects and are not constructed here.

\item By \cref{cusp-cone-S12-almost-complex-obstruction}, no smooth manifold
homeomorphic to \(S^{12}\) admits an almost-complex structure.  Nor does
\(S^{24}\), with any differentiable structure, by the Borel--Serre theorem
\cite[\S15, Proposition~15.1 and the following remark,
p.~434]{BorelSerre1953}.  Thus the smash sources cannot be promoted to
compact complex spheres.
\end{enumerate}
\end{theorem}

\begin{proof}
Every element of \(K_\Lambda^{\oplus r}\) has uniformly bounded physical
length on the radial cusp.  If one survived in the quotient lattice, it
would be a nonzero period of bounded quotient norm.  Hence
\begin{equation}
 K_\Lambda^{\oplus r}\subset N.
 \label{cusp-cone-kernel-contained}
\end{equation}
For the full monodromy group,
\(A_1\widehat w=\widehat u-\widehat w\) and invariance force every copy of
\(\widehat u\) into \(N\).  Therefore
\[
 L_\Lambda^{\oplus r}
 =\mathbb Z\langle\widehat u,\widehat w,\widehat\delta\rangle^{\oplus r}
 \subset N,
\]
which proves \cref{cusp-cone-rank-bound}.

The same conclusion holds after a finite base cover.  Let \(H\) be its
finite-index subgroup and put \(H_*=A_2A_1=I+N_*\), \(N_*^2=0\).
Every finite-index subgroup contains \(H_*^n\) for some \(n>0\), while
\begin{equation}
 H_*^n\widehat w
 =\widehat w+n(\widehat u-\widehat\delta).
 \label{cusp-cone-finite-index-orbit}
\end{equation}
Equations
\cref{cusp-cone-kernel-contained,cusp-cone-finite-index-orbit} give
\(n\widehat u\in N\).  Primitivity gives \(\widehat u\in N\), proving the
same rank bound.  The orbit also gives
\(H_*^nK_\Lambda\ne K_\Lambda\) for every \(n>0\), so no finite cover of
the regular base preserves the native kernel.  A compact cover branched
over the singular values restricts to such a cover.  Finally, a connected
finite cover of \(\mathbb P^1\) branched only at \(p_0\) would restrict to
a connected unbranched finite cover of
\(\mathbb P^1\setminus\{p_0\}\cong\mathbb C\), hence has degree one.

For a fixed \(T^4\) target, \(H^1(X^r;\mathbb Z)=0\), as does the first
cohomology of the smash spheres.  A global map therefore pulls every target
character class back to zero.  On a regular fibre the quotient homomorphisms
inject the target character lattice into the character lattice of the
source torus, giving four independent nonzero pullbacks, a contradiction.

For the fixed \(S^4\) obstruction, let
\(\omega_T\in H^2(Q_T;\mathbb Z)\) be the orientation class.  Since
\(\pi_T^*:H^1(Q_T)\to H^1(F_T)\) injects two independent classes and torus
cohomology is their exterior algebra,
\begin{equation}
 (\pi_T\times\pi_T)^*(\omega_T\times\omega_T)
 \ne0\quad\text{in }H^4(F_T^2;\mathbb Z).
 \label{cusp-cone-fibre-H4-nonzero}
\end{equation}
A degree-one collapse \(Q_T^2\to S^4\) pulls the orientation generator
back to \(\pm(\omega_T\times\omega_T)\).  But
\[
 H^4(X^2;\mathbb Z)=H^4(S^{12};\mathbb Z)=0,
\]
so no global extension has this restriction.  The four-circle case is the
same: its orientation class pulls back to the exterior product of four
independent circle characters, while
\(H^4(X^4)=H^4(S^{24})=0\).

The Kunneth theorem gives \(H_6(X^2)=\mathbb Z^2\), whereas
\(H_6(S^{12})=0\), and gives nonzero groups of ranks \(4,6,4\) in degrees
\(6,12,18\) for \(X^4\), whereas \(S^{24}\) has none.  The real dimension
of an \(r\)-fold synchronized fibre product is
\(6r-2(r-1)=4r+2\), giving \(10\) and \(18\).  The join formula
\(S^m*S^n\cong S^{m+n+1}\) proves
\cref{cusp-cone-join-dimensions}.  The remaining assertions follow from
simple connectivity, \cref{cusp-cone-S12-almost-complex-obstruction}, and
the cited Borel--Serre theorem for \(S^{24}\).
\end{proof}

\paragraph{Canonicity and exact scope.}
The lattice \(K_\Lambda=\ker(M_0-I)=\operatorname{im}(M_0-I)\) and the
punctured-cusp quotient are intrinsic to the marked cusp monodromy.
The compact map \(q_2\) additionally depends on the chosen cusp coordinate,
closed cusp disc, and the decision to collapse its outer boundary and the
exterior.  The identifications
\(X^{\wedge2}\cong S^{12}\) and \(X^{\wedge4}\cong S^{24}\) depend on a
based choice of \(X\cong S^6\).  None of these choices is claimed to be
holomorphic or canonical in the algebraic category.

The map \(q_2^{\wedge2}\) is the direct two-copy construction from the
native \(T^2\times\mathbb R^2\) cusp: it is globally defined on the smash
sphere, and its radial regular-fibre quotients are \(Q_T\times Q_T\), which
converge pointed local-isometrically, smoothly, and Gromov--Hausdorff to
\(\mathbb R^4\).  It is not a torus quotient over the full
regular base.  The map \(q_{1,\mathrm{cusp}}^{\wedge4}\) is the parallel
cusp-supported construction with target
\(\operatorname{Susp}(T^2)^{\wedge4}\).  The stronger map
\(q_1^{\wedge4}\) uses the quotient over every regular and finite-filled
fibre before the cusp cone is attached; four copies are minimal for that
full-base all-expanding group quotient.
Neither construction gives a compact complex \(S^{12}\) or \(S^{24}\), an
ambient Euclidean limit in a fixed compact metric target, or a claim about
the interacting Yang--Mills spectrum.

The argument does not rule out a separately constructed branched or
toroidal compactification of the auxiliary finite covers after additional
branching away from \(p_0\).  Such a construction would have to specify its
global base cover, normalize and resolve every pulled-back singular fibre,
track the lifted metric, and determine the topology of the resulting total
space.  No such smooth compact sphere construction follows from the present
data.


% Canonical continuation of cusp_spectral_geometry.tex.  It uses the notation
% and labels of cusp_globalization_replacement.tex and
% cusp_cone_sphere_routes.tex, which are included immediately before it.

\subsubsection{Optimal sphere exponents and the exact Euclidean boundary}
\label{cusp-cone-optimal-strengthening}

Within the class of native connected fibrewise torus quotients equipped with
the induced, unscaled quotient of the marked Euclidean metric, the two sphere
dimensions in the preceding constructions are forced by two different
integral quotient problems.  Cusp-locally, each copy has an expanding
quotient of rank two, so two copies are necessary and sufficient within that
class.  Over the full regular base, monodromy leaves at most one expanding
quotient direction per copy, so four copies are necessary and sufficient
within the same class.  We now prove this scoped optimality, identify the
compact conical carriers, and construct a noncompact comparison manifold in
which the limiting \(\mathbb R^4\) is an actual boundary fibre.  There is only
a continuous comparison map from this manifold to the compact conical
carrier: it set-theoretically collapses the new boundary, is a homeomorphism
away from that boundary, and is nonproper and nonquotient.

Throughout this subsubsection,
\[
 B_R^Z(z):=\{w\in Z:d_Z(z,w)\leq R\}
\]
denotes the closed radius-\(R\) ball; when the ambient space and center are
clear, superscripts and centers are suppressed.

\begin{theorem}[Optimality within the native connected torus-quotient class]
\label{cusp-cone-optimal-exponents}
Consider only native connected fibrewise torus quotients of copies of the
original marked fibres, and equip every such quotient with the induced,
unscaled Euclidean quotient metric.

\begin{enumerate}
\item For a native connected torus quotient of \(r\) copies over the
punctured cusp, if the quotient systole tends to infinity on the positive
radial ray, then its lattice rank is at most \(2r\).  A rank-four quotient
therefore requires \(r\geq2\).  Equality is attained for \(r=2\) by
\[
 (\Lambda/K_\Lambda)^{\oplus2},
 \qquad K_\Lambda=
 \mathbb Z\langle\widehat w,\widehat\delta\rangle,
\]
whose physical period lattice is
\[
 L_T^{(12)}=\operatorname{diag}(Y_T,Y_T)\mathbb Z^4.
\]
The corresponding smash source is
\(X^{\wedge2}\cong S^{12}\).

\item For a native connected torus quotient of \(r\) copies over the full
regular base, or over a finite cover of that base, radial divergence of the
quotient systole forces its lattice rank to be at most \(r\).  A rank-four
quotient therefore requires \(r\geq4\).  Equality is attained by the four
copies of the full-monodromy-invariant character \(\chi_\gamma\); on the
radial fibre its period lattice is
\[
 L_T^{(24)}=\ell_T\mathbb Z^4,
 \qquad
 \ell_T=T+H-B+O(T^{-1}),
\]
and the corresponding smash source is
\(X^{\wedge4}\cong S^{24}\).
\end{enumerate}

Thus, among precisely these native connected torus quotients with their
induced unscaled marked Euclidean metrics, \(S^{12}\) is the smallest smash
sphere produced by the cusp-local rank-two quotient, and \(S^{24}\) is the
smallest smash sphere produced by a quotient which survives the full
regular-base monodromy.  No optimality is asserted here among arbitrary
continuous quotient maps or after a \(T\)-dependent rescaling of the metric.
\end{theorem}

\begin{proof}
For \(r\) copies, every element of
\[
 K_\Lambda^{\oplus r}
 =\mathbb Z\langle\widehat w,\widehat\delta\rangle^{\oplus r}
\]
has uniformly bounded physical length on the radial cusp.  Let
\(N\subset\Lambda^{\oplus r}\) be a primitive sublattice defining a
connected torus quotient.  If some \(v\in K_\Lambda^{\oplus r}\) had
nonzero class in \(\Lambda^{\oplus r}/N\), that nonzero quotient period
would have quotient norm at most the original norm of \(v\), uniformly in
\(T\).  The quotient systole could not diverge.  Hence
\[
 K_\Lambda^{\oplus r}\subset N,
\]
and
\[
 \operatorname{rank}(\Lambda^{\oplus r}/N)
 \leq 4r-2r=2r.
\]
For \(r=2\), the primitive kernel \(K_\Lambda^{\oplus2}\) attains this
bound.  The cusp monodromy is the identity modulo \(K_\Lambda\), so the
quotient descends over the punctured cusp.  Its radial lattice is exactly
\(\operatorname{diag}(Y_T,Y_T)\mathbb Z^4\), and
\[
 \operatorname{sys}(L_T^{(12)})
 =\operatorname{sys}(Y_T\mathbb Z^2)\geq T-C\longrightarrow\infty.
\]

For a quotient over the full regular base, monodromy invariance gives a
stronger bound.  Since \(\widehat w\in N\) and
\[
 A_1\widehat w=\widehat u-\widehat w,
\]
one has \(\widehat u\in N\).  Together with
\(\widehat w,\widehat\delta\in N\), this places the rank-three lattice
\[
 \mathbb Z\langle\widehat u,\widehat w,
 \widehat\delta\rangle^{\oplus r}
\]
inside \(N\), and hence
\[
 \operatorname{rank}(\Lambda^{\oplus r}/N)\leq r.
\]
The same conclusion holds after a finite base cover.  Indeed, if
\(\Gamma_{\mathrm{mon}}\) is the corresponding finite-index monodromy
subgroup and \(G_*=A_2A_1=I+N_*\), then \(N_*^2=0\), and
\(\Gamma_{\mathrm{mon}}\) contains \(G_*^m\) for
some \(m>0\).  The exact orbit formula
\[
 G_*^m\widehat w
 =\widehat w+m(\widehat u-\widehat\delta)
\]
implies \(m\widehat u\in N\); primitivity of \(N\) implies
\(\widehat u\in N\).

The invariant character \(\chi_\gamma\) has precisely the complementary
rank-one quotient in each copy.  The full-base construction, including the
two finite Seifert fillings and the cone at the unipotent cusp, proves that
four copies attain rank four.  Finally, after choosing the based
diffeomorphism \(X\cong_{\mathrm{diff}}S^6\) already reconstructed in this
workbench,
\(X^{\wedge r}\cong(S^6)^{\wedge r}\cong S^{6r}\), which gives the two
claimed sphere dimensions.
\end{proof}

\begin{theorem}[Exact sphere-to-sphere carriers with moving four-spheres]
\label{cusp-cone-angle-forgotten-sphere-carriers}
There are continuous surjective based quotient maps
\[
 \widehat q_3:(X,x_0)\longrightarrow(S^3,*),
 \qquad
 \widehat q_\gamma:(X,x_0)\longrightarrow(S^2,*),
\]
whose smash powers are continuous surjective quotient maps
\begin{equation}
 \widehat q_3^{\wedge2}:S^{12}\twoheadrightarrow S^6,
 \qquad
 \widehat q_\gamma^{\wedge4}:S^{24}\twoheadrightarrow S^8.
 \label{cusp-cone-sphere-to-sphere-maps}
\end{equation}
Use the marked presentation \(\overline T^2=\mathbb R^2/\mathbb Z^2\) and
put
\[
 K^{(1)}
 =((\mathbb R/\mathbb Z)\times\{0\})
  \cup(\{0\}\times(\mathbb R/\mathbb Z)),
 \qquad
 c:\overline T^2\longrightarrow\overline T^2/K^{(1)}\cong S^2.
\]
Here \(c\) is the quotient collapsing the coordinate one-skeleton.  For
every sufficiently large \(T\), put
\begin{align*}
 r_T&=e^{-2\pi T},\\
 K_T^{(1)}&:=\iota_T(K^{(1)})\subset Q_T,\\
 S_T^2&:=\overline T^2/K^{(1)}\cong Q_T/K_T^{(1)},\\
 c_T&:=c\circ\iota_T^{-1}:Q_T\longrightarrow S_T^2,
\end{align*}
where the identification of the two displayed quotients is induced by
\(\iota_T\).  Within this theorem and its proof, \(c_T\) denotes the
displayed collapse map; the earlier scalar \(c_T=\xi_T-\eta_T\) does not
enter.
At reduced-suspension cross-section height \(r_T\), the first target contains
a based topologically embedded four-sphere, equipped with its intrinsic
quotient length metric,
\[
 \mathbb S_{12,T}^4=S_T^2\wedge S_T^2\hookrightarrow S^6,
\]
and the second target contains, at the same cross-section height, a based
topologically embedded four-sphere, equipped with its intrinsic quotient
length metric,
\[
 \mathbb S_{24,T}^4=C_T^{\wedge4}\hookrightarrow S^8.
\]
On synchronized radial fibres, the exact restrictions factor as
\begin{align}
 F_T^2&\xrightarrow{\ \pi_T\times\pi_T\ }
 Q_T^2\xrightarrow{\ c_T\times c_T\ }(S_T^2)^2
 \xrightarrow{\ \rho_{12,T}\ }S_T^2\wedge S_T^2
 =\mathbb S_{12,T}^4,
 \label{cusp-cone-S12-S4-radial-map}\\
 F_T^4&\xrightarrow{\ \chi_{\gamma,T}^{\times4}\ }
 C_T^4\xrightarrow{\ \rho_{24,T}\ }C_T^{\wedge4}
 =\mathbb S_{24,T}^4,
 \label{cusp-cone-S24-S4-radial-map}
\end{align}
Here \(c_T\) collapses \(K_T^{(1)}\), and \(\rho_{12,T}\) and
\(\rho_{24,T}\) are the smash quotient maps.  Thus the
full collapsed locus in \(Q_T^2\) is exactly
\begin{equation}
 W_{12,T}
 =\bigl(K_T^{(1)}\times Q_T\bigr)
  \cup\bigl(Q_T\times K_T^{(1)}\bigr),
 \label{cusp-cone-S12-exact-collapsed-locus}
\end{equation}
and the full collapsed locus in \(C_T^4\) is exactly
\begin{equation}
 W_{24,T}
 =\{(x_1,x_2,x_3,x_4)\in C_T^4:
       x_j=0\pmod{\ell_T}\ \text{for some }j\}.
 \label{cusp-cone-S24-exact-collapsed-locus}
\end{equation}
The maps
\(\rho_{12,T}\circ(c_T\times c_T):Q_T^2\to
S_T^2\wedge S_T^2\) and
\(\rho_{24,T}:C_T^4\to C_T^{\wedge4}\) send exactly \(W_{12,T}\) and
\(W_{24,T}\), respectively, to the smash basepoint and are injective on
their complements.  The full composites from \(F_T^2\) and \(F_T^4\)
still identify every pair of points identified by
\(\pi_T\times\pi_T\) and \(\chi_{\gamma,T}^{\times4}\), respectively; in
particular, their first arrows retain their nontrivial torus kernels.
Their exact smash-basepoint preimages in the full radial sources are
\[
 \bigl(\pi_T^{-1}(K_T^{(1)})\times F_T\bigr)
 \cup\bigl(F_T\times\pi_T^{-1}(K_T^{(1)})\bigr)
\]
and
\[
 \{(x_1,x_2,x_3,x_4)\in F_T^4:
   \chi_{\gamma,T}(x_j)=0\ \text{in }C_T\ \text{for some }j\},
\]
respectively.  Outside these basepoint preimages, two source points have the
same final image if and only if their images under the corresponding first
arrow in \cref{cusp-cone-S12-S4-radial-map,cusp-cone-S24-S4-radial-map}
are equal.

Choose the Euclidean Riemannian \(\ell^2\)-product length metrics on
\(Q_T^2\) and \(C_T^4\) before taking the smash quotients, and use the
successive quotient length metrics specified in the proof.  With the
specified midpoint basepoints,
\begin{equation}
 (\mathbb S_{12,T}^4,\zeta_{12,T})
 \longrightarrow(\mathbb R^4,0),
 \qquad
 (\mathbb S_{24,T}^4,\zeta_{24,T})
 \longrightarrow(\mathbb R^4,0)
 \label{cusp-cone-moving-S4-R4-limits}
\end{equation}
in pointed local-isometric and pointed Gromov--Hausdorff topology.  Moreover,
the convergence is smooth on the eventually Euclidean basepoint
neighborhoods.  For every \(R>0\), once \(T-C>4R\), respectively once
\(\ell_T>4R\), the corresponding quotient map restricts to an isometry from
the radius-\(R\) ball in \(Q_T^2\), respectively \(C_T^4\), \emph{onto} the
entire radius-\(R\) ball in the indicated smash sphere.  The domain ball, and
hence the quotient ball, is exactly a Euclidean four-ball.
\end{theorem}

\begin{proof}
Write \(\Sigma_r(Z,z_0)\) for the reduced suspension
\[
 \frac{[0,\epsilon]\times Z}
 {(\{0,\epsilon\}\times Z)\cup([0,\epsilon]\times\{z_0\})},
\]
where the displayed union is collapsed to one basepoint.

Use the marked presentation
\(\overline T^2=\mathbb R^2/\mathbb Z^2\), let
\[
 K^{(1)}=
 ((\mathbb R/\mathbb Z)\times\{0\})
 \cup(\{0\}\times(\mathbb R/\mathbb Z)),
\]
and let
\[
 c:\overline T^2\longrightarrow
 \overline T^2/K^{(1)}\cong S^2
\]
be the degree-one quotient obtained by collapsing the one-skeleton.  The
formula
\[
 \widehat p_2:\mathcal A\longrightarrow
 \Sigma_r(\overline T^2/K^{(1)},c(K^{(1)}))\cong S^3,
 \qquad
 [r,e^{i\phi},u]\longmapsto[r,c(u)]
\]
is well defined: the two suspension ends of \(\mathcal A\) and the entire
cylinder over \(K^{(1)}\) map to the reduced-suspension basepoint.  It is a
continuous surjection.  Set
\(\widehat q_3=\widehat p_2\circ q_2\).

For every sufficiently large \(T\), define the transported one-skeleton and
its collapse by
\begin{equation}
 \begin{aligned}
 K_T^{(1)}&:=\iota_T(K^{(1)})\subset Q_T,\\
 S_T^2&:=\overline T^2/K^{(1)}\cong Q_T/K_T^{(1)},\\
 c_T&:=c\circ\iota_T^{-1}:Q_T\longrightarrow S_T^2.
 \end{aligned}
 \label{cusp-cone-transported-skeleton-collapse}
\end{equation}
The displayed target identification is the one induced by \(\iota_T\), so
\(c_T^{-1}(*)=K_T^{(1)}\), where
\(*=c(K^{(1)})=c_T(K_T^{(1)})\).

Give \(Q_T\) its Euclidean quotient length metric and \(S_T^2\) the quotient
length metric under \(c_T\).  Put
\[
 y_*=[(1/2,1/2)]\in\overline T^2,
 \qquad y_T:=\iota_T(y_*)\in Q_T,
 \qquad z_T:=c_T(y_T)\in S_T^2.
\]
Equivalently, the metric on \(\overline T^2\) used below is
\(g_T:=\iota_T^*g_{Q_T}\).  For every \(v\in\mathbb R^2\), the exact
decomposition
\[
 Y_T=T\begin{pmatrix}0&1\\-1&0\end{pmatrix}+E_T,
 \qquad\|E_T\|_{\mathrm{op}}\leq C,
\]
gives
\[
 \|Y_Tv\|\geq(T-C)\|v\|.
\]
The standard flat distance from \(y_*\) to \(K^{(1)}\) is \(1/2\);
therefore
\begin{equation}
 d_T:=\operatorname{dist}_{Q_T}(y_T,K_T^{(1)})
 =\operatorname{dist}_{g_T}(y_*,K^{(1)})
 \geq\frac{T-C}{2}.
 \label{cusp-cone-distance-to-skeleton}
\end{equation}
For a length-space quotient which collapses a closed subset \(K\) to one
point, the distance between two noncollapsed classes represented by \(x,y\)
is
\[
 \min\{d(x,y),d(x,K)+d(y,K)\}.
\]
The two displayed terms are upper bounds, obtained respectively from a
direct path and from paths from \(x\) and \(y\) to \(K\).  Conversely, a
rectifiable quotient path either avoids the collapsed point, in which case
its compact image lifts to \(X\setminus K\) and has length at least
\(d(x,y)\), or meets that point, in which case its two pieces have total
length at least \(d(x,K)+d(y,K)\).  Taking infima proves the formula.
If \(T-C>4R\), then \(d_T>2R\).  The displayed formula and the systolic
bound show that
\[
 c_T:B_R^{Q_T}(y_T)\longrightarrow B_R^{S_T^2}(z_T)
\]
is an isometry onto the entire target ball, and both balls are exactly
Euclidean two-balls.  Indeed, the second term in the minimum cannot realize
a distance at most \(2R\), and any point of the target ball has a unique
noncollapsed representative; the same formula puts that representative in
the source ball.

At reduced-suspension coordinate \(r_T=e^{-2\pi T}\), \(S_T^2\) is a based
embedded subspace of \(S^3\): its point \(*\) is the global
reduced-suspension basepoint.  Hence
\[
 \mathbb S_{12,T}^4=S_T^2\wedge S_T^2
\]
embeds in \(S^3\wedge S^3=S^6\).  Equip \(Q_T^2\) first with the Euclidean
Riemannian \(\ell^2\)-product length metric, equip \((S_T^2)^2\) with the
\(\ell^2\)-product of the two quotient length metrics, and let
\[
 \rho_{12,T}:(S_T^2)^2\longrightarrow S_T^2\wedge S_T^2
\]
carry the quotient length metric obtained by collapsing the fat wedge
\[
 \mathcal W_{12,T}
 =(S_T^2\times\{*\})\cup(\{*\}\times S_T^2).
\]
Since \(c_T^{-1}(*)=K_T^{(1)}\), the inverse image of
\(\mathcal W_{12,T}\) under \(c_T\times c_T\) is exactly \(W_{12,T}\) in
\cref{cusp-cone-S12-exact-collapsed-locus}.  On
\(Q_T^2\setminus W_{12,T}\), the map \(c_T\times c_T\) is injective and its
image avoids \(\mathcal W_{12,T}\); the map \(\rho_{12,T}\) is injective on
the complement of \(\mathcal W_{12,T}\).  Hence their composite identifies
all of \(W_{12,T}\) to the smash basepoint and makes no other
identifications.

The point \(\zeta_{12,T}=[z_T,z_T]\) has distance \(d_T\) from this fat
wedge.  The one-factor onto isometry above gives an onto isometry from the
radius-\(R\) ball about \((y_T,y_T)\) in \(Q_T^2\) to the corresponding
ball about \((z_T,z_T)\) in \((S_T^2)^2\).  Applying the same collapsed-set
distance formula to \(\rho_{12,T}\), and using \(d_T>2R\), shows that
\begin{equation}
 \rho_{12,T}\circ(c_T\times c_T):
 B_R^{Q_T^2}((y_T,y_T))
 \xrightarrow{\ \cong_{\mathrm{isom}}\ }
 B_R^{\mathbb S_{12,T}^4}(\zeta_{12,T})
 \label{cusp-cone-S12-onto-exact-ball}
\end{equation}
is onto as well as isometrically embedded.  The systolic bound
\(\operatorname{sys}(Q_T^2)\geq T-C>4R\) makes its domain exactly a
Euclidean four-ball.  On a radial fibre,
\[
 (\widehat p_2\circ q_2)|_{F_T}
 =[r_T,c(\iota_T^{-1}(\pi_T(\,\cdot\,)))]
 =[r_T,c_T(\pi_T(\,\cdot\,))],
\]
and taking two factors followed by \(\rho_{12,T}\) proves the entire
factorization in \cref{cusp-cone-S12-S4-radial-map}.

For the four-copy construction, choose \(1\in S^1\) as character basepoint
and define
\[
 \widehat p_1:\mathcal B\longrightarrow\Sigma_r(S^1,1)\cong S^2,
 \qquad
 [r,e^{i\phi},\theta]\longmapsto[r,\theta].
\]
The formula forgets the base angle and collapses the two suspension ends
and the entire cylinder with \(\theta=1\) to the reduced-suspension
basepoint.  It is a continuous surjection.  Put
\(\widehat q_\gamma=\widehat p_1\circ q_{1,\mathrm{cusp}}\).

At the interior suspension coordinate \(r_T=e^{-2\pi T}\), identify the
character circle \(C_T=\mathbb R/\ell_T\mathbb Z\) with
\(\mathbb R/\mathbb Z\subset S^2\) by \([x]\mapsto[x/\ell_T]\).  It is a
based embedded circle, so its fourth smash is an embedded
\[
 \mathbb S_{24,T}^4=C_T^{\wedge4}\subset(S^2)^{\wedge4}=S^8.
\]
Equip \(C_T^4\) first with the Euclidean Riemannian \(\ell^2\)-product
length metric, and let
\[
 \rho_{24,T}:C_T^4\longrightarrow C_T^{\wedge4}
\]
carry the quotient length metric obtained by collapsing exactly the fat
wedge \(W_{24,T}\) of
\cref{cusp-cone-S24-exact-collapsed-locus}.  Away from that fat wedge the
smash quotient is injective, so no further points are identified.  The
midpoint
\[
 \zeta_{24,T}
 =[(\ell_T/2,\ell_T/2,\ell_T/2,\ell_T/2)]
\]
has distance \(\ell_T/2\) from that fat wedge.  If \(\ell_T>4R\), neither a
nonzero period nor a path through the collapsed fat wedge can shorten a
distance between points in the radius-\(R\) ball of \(C_T^4\).  Conversely,
the collapsed-set distance formula shows that every point of the
radius-\(R\) quotient ball has its unique noncollapsed representative in
that product ball.  Consequently
\begin{equation}
 \rho_{24,T}:
 B_R^{C_T^4}((\ell_T/2,\ell_T/2,\ell_T/2,\ell_T/2))
 \xrightarrow{\ \cong_{\mathrm{isom}}\ }
 B_R^{\mathbb S_{24,T}^4}(\zeta_{24,T})
 \label{cusp-cone-S24-onto-exact-ball}
\end{equation}
is onto as well as isometrically embedded, and its domain is exactly a
Euclidean four-ball because \(\operatorname{sys}(C_T^4)=\ell_T>4R\).  On
the radial source the formula discards only the base angle and then performs
this smash quotient, giving
\cref{cusp-cone-S24-S4-radial-map}.

Both one-copy maps are continuous surjections from compact \(X\) to
Hausdorff spheres, hence quotient maps, and they are based.  Their smash
powers are surjective; compact-to-Hausdorff again makes them quotient maps.
Using the same chosen based diffeomorphism \(X\cong_{\mathrm{diff}}S^6\),
the homeomorphisms
\[
 X^{\wedge2}\cong S^{12},\quad(S^3)^{\wedge2}\cong S^6,
 \qquad
 X^{\wedge4}\cong S^{24},\quad(S^2)^{\wedge4}\cong S^8
\]
prove \cref{cusp-cone-sphere-to-sphere-maps}.  The exact Euclidean-ball
statements prove all three convergences in
\cref{cusp-cone-moving-S4-R4-limits}.
\end{proof}

\begin{corollary}[The two-copy rank-two sphere carrier is null-homotopic]
\label{cusp-cone-S12-sphere-map-null}
After the based identifications in
\cref{cusp-cone-sphere-to-sphere-maps}, the continuous surjective quotient
\[
 \widehat q_3^{\wedge2}:S^{12}\twoheadrightarrow S^6
\]
is based null-homotopic.  In fact, \(g\wedge g:S^{12}\to S^6\) is based
null-homotopic for every based map \(g:S^6\to S^3\).
\end{corollary}

\begin{proof}
Apply the full external-smash calculation in
\cref{cusp-s12-full-carrier-null} to
\(g=\widehat q_3\).  That theorem proves the stronger assertion that
\(x\wedge y:S^{12}\to S^6\) is null-homotopic for every pair of based
maps \(x,y:S^6\to S^3\).
\end{proof}

\begin{corollary}[Every repeated smash of the circle-quotient carrier is null]
\label{cusp-cone-S24-sphere-map-null}
Put \(f=\widehat q_\gamma:S^6\to S^2\) after the chosen based
identification \(X\cong S^6\).  For every \(r\geq2\),
\begin{equation}
 \bigl[f^{\wedge r}\bigr]=0
 \quad\text{in}\quad
 \pi_{6r}(S^{2r}).
 \label{cusp-cone-S24-null-maps}
\end{equation}
In particular,
\begin{equation}
 \bigl[\widehat q_\gamma^{\wedge2}\bigr]
 =0\in\pi_{12}(S^4),
 \qquad
 \bigl[\widehat q_\gamma^{\wedge4}\bigr]
 =0\in\pi_{24}(S^8).
 \label{cusp-cone-S24-null-special-cases}
\end{equation}
Thus the displayed surjective quotient
\(\widehat q_\gamma^{\wedge4}:S^{24}\twoheadrightarrow S^8\) is globally
null-homotopic.  This does not alter either its point-set surjectivity or the
exact intrinsic \(\mathbb R^4\) limit on the moving four-spheres
\(\mathbb S_{24,T}^4\): those are properties of the specified representative
and its radial restrictions, whereas null-homotopy classifies that
representative only up to based homotopy.
\end{corollary}

\begin{proof}
Associativity of the smash product gives the
exact factorization
\begin{equation}
 \begin{tikzcd}[column sep=large]
 S^6\wedge S^6
 \arrow[r,"f\wedge\operatorname{id}_{S^6}"]&
 S^2\wedge S^6
 \arrow[r,"\operatorname{id}_{S^2}\wedge f"]&
 S^2\wedge S^2 .
 \end{tikzcd}
 \label{cusp-cone-S24-smash-square-factorization}
\end{equation}
Under the standard based sphere identifications, the first arrow is the
sixfold suspension of \(f\), up to the harmless twist homeomorphisms dictated
by the left-versus-right suspension convention:
\begin{equation}
 f\wedge\operatorname{id}_{S^6}
 =\Sigma^6f:S^{12}\longrightarrow S^8.
 \label{cusp-cone-S24-sixfold-suspension}
\end{equation}
Freudenthal's suspension theorem is already in its isomorphism range here,
because \(12\leq 2\cdot8-2\), and remains in that range after every further
suspension.  It therefore identifies \(\pi_{12}(S^8)\) with the fourth
stable stem.  The low-stem calculation gives
\begin{equation}
 \pi_{12}(S^8)\cong\pi_4^{\mathrm S}=0
 \quad
 \text{\cite[Corollary~4.24, p.~360]{Hatcher2002};
 \cite[Table~1, p.~113]{IsaksenWangXu2023}}.
 \label{cusp-cone-fourth-stable-stem}
\end{equation}
Hence \(\Sigma^6f\) is based null-homotopic, so the composite in
\cref{cusp-cone-S24-smash-square-factorization}, namely
\(f\wedge f\), is based null-homotopic.  For \(r\geq2\), the canonical
associativity homeomorphisms identify
\begin{equation}
 f^{\wedge r}=(f\wedge f)\wedge f^{\wedge(r-2)},
\end{equation}
with the convention \(f^{\wedge0}=\operatorname{id}_{S^0}\).  Smashing a
based null-homotopy with any based map is a based null-homotopy.  This proves
\cref{cusp-cone-S24-null-maps,cusp-cone-S24-null-special-cases}.
\end{proof}

\begin{corollary}[Literal \(\mathbb R^4\) open cells in the sphere targets]
\label{cusp-cone-literal-R4-open-cells}
Let \(*\) denote the smash basepoint and set
\[
 U_{12,T}:=\mathbb S_{12,T}^4\setminus\{*\},
 \qquad
 U_{24,T}:=\mathbb S_{24,T}^4\setminus\{*\}.
\]
The embedded moving four-spheres restrict to topological embeddings
\[
 U_{12,T}\cong\mathbb R^4\hookrightarrow S^6,
 \qquad
 U_{24,T}\cong\mathbb R^4\hookrightarrow S^8.
\]
More exactly, the quotient-stage maps restrict to homeomorphisms
\begin{align}
 (Q_T\setminus K_T^{(1)})^2
 &\xrightarrow{\ \cong\ }U_{12,T},
 \label{cusp-cone-U12-homeomorphism}\\
 (C_T\setminus\{0\})^4
 &\xrightarrow{\ \cong\ }U_{24,T}.
 \label{cusp-cone-U24-homeomorphism}
\end{align}
These copies of \(\mathbb R^4\) are nonclosed topological open cells: they
are open in their moving four-spheres, not open in the ambient \(S^6\) or
\(S^8\), and their closures in those ambient spheres are
\(\mathbb S_{12,T}^4\) and \(\mathbb S_{24,T}^4\), respectively.  They are
not globally isometric to Euclidean \(\mathbb R^4\) for either the assigned
compact quotient metrics or the restriction of any fixed compatible metric
on the compact sphere targets.  The exact midpoint balls in
\cref{cusp-cone-S12-onto-exact-ball,cusp-cone-S24-onto-exact-ball} give the
precise metric statement: as \(T\to\infty\), every fixed Euclidean
radius-\(R\) ball occurs isometrically and onto the corresponding pointed
ball.
\end{corollary}

\begin{proof}
For a quotient which collapses a closed subset \(K\) to a point, restriction
to the complement gives a homeomorphism
\[
 X\setminus K\xrightarrow{\ \cong\ }(X/K)\setminus\{*\}.
\]
Indeed, the complement of the quotient point is open, its inverse image is
the saturated open set \(X\setminus K\), and restriction of a quotient map
over an open subset of its target is again a quotient map.  The restricted
map is bijective, and a bijective quotient map is a homeomorphism.
Consequently
\[
 Q_T\setminus K_T^{(1)}
 \cong S_T^2\setminus\{*\}.
\]
The marked coordinate one-skeleton \(K^{(1)}\) is the \(1\)-skeleton of the
standard one-vertex CW structure on \(\overline T^2\), with one \(0\)-cell,
two \(1\)-cells, and one \(2\)-cell.  Its complement is therefore the open
\(2\)-cell \((0,1)^2\cong\mathbb R^2\).  Transport by
\(\iota_T\) gives
\[
 Q_T\setminus K_T^{(1)}\cong\mathbb R^2.
\]
Likewise \(C_T\setminus\{0\}\) is an open interval and hence is homeomorphic
to \(\mathbb R\).

The complement of the smash basepoint is the product of the complements of
the factor basepoints.  Therefore
\[
 (S_T^2\wedge S_T^2)\setminus\{*\}
 \cong(S_T^2\setminus\{*\})^2\cong\mathbb R^4,
\]
and
\[
 C_T^{\wedge4}\setminus\{*\}
 \cong(C_T\setminus\{0\})^4\cong\mathbb R^4.
\]
The full moving four-spheres are topologically embedded in their compact
sphere targets, so these restrictions are topological embeddings.  In each
case the omitted smash basepoint lies in the closure of the open cell, which
proves the closure assertions.  Finally, every compatible metric on a
compact target is bounded, as is each compact quotient metric assigned to a
moving four-sphere, whereas Euclidean \(\mathbb R^4\) is unbounded.  Thus no
global isometry is possible; the exact metric content is precisely the
onto-ball calculation already proved.
\end{proof}

\begin{remark}[The moving four-spheres are null-homotopic in the targets]
\label{cusp-cone-moving-S4-null-homotopy}
For a fixed interior reduced-suspension coordinate, the inclusion
\(S_T^2\hookrightarrow S^3\) contracts through based maps by sliding that
coordinate to a suspension endpoint; throughout the homotopy the collapsed
meridian remains at the global basepoint.  Hence
\(S_T^2\wedge S_T^2\hookrightarrow S^6\) is null-homotopic.  The identical
sliding argument contracts \(C_T\hookrightarrow S^2\), and therefore also
contracts \(C_T^{\wedge4}\hookrightarrow S^8\).  Thus the moving metric
four-spheres do not contradict the fixed-target orientation obstruction:
their fundamental classes die in the ambient targets.
\end{remark}

\begin{remark}[Homotopical status of the two sphere-to-sphere maps]
\label{cusp-cone-sphere-map-homotopy-status}
The maps \(S^{12}\to S^6\) and \(S^{24}\to S^8\) induce zero on all
positive-degree integral homology and cohomology groups and have zero
ordinary Hurewicz image.  They have no classical degree because source and
target dimensions differ.  The full four-copy map
\(\widehat q_\gamma^{\wedge4}:S^{24}\to S^8\) is based null-homotopic by
\cref{cusp-cone-S24-sphere-map-null}.  The distinct two-copy map
\(\widehat q_3^{\wedge2}:S^{12}\to S^6\) is also based null-homotopic, by
\cref{cusp-cone-S12-sphere-map-null}.  In both cases the vanishing is a
theorem about the full specified sphere-valued map, not merely its Hurewicz
image or the inclusion of a moving four-sphere. The one-copy class
\([\widehat q_3]\in\pi_6(S^3)\) is now zero by
Theorem~\ref{cusp-null:thm}; the one-copy class
\([\widehat q_\gamma]\in\pi_6(S^2)\) remains open.
\end{remark}

\begin{proposition}[The compact cone-on-\(T^4\) carriers]
\label{cusp-cone-explicit-carriers}
Fix \(T_0\) sufficiently large and put
\(r_0=e^{-2\pi T_0}<\epsilon\).  Define
\[
 \mathcal D_{12}
 =\frac{[0,r_0]\times(\overline T^2)^2}
 {\{0\}\times(\overline T^2)^2\sim d_{12}},
 \qquad
 \mathcal D_{24}
 =\frac{[0,r_0]\times(S^1)^4}
 {\{0\}\times(S^1)^4\sim d_{24}}.
\]
Both spaces are homeomorphic to \(\operatorname{Cone}(T^4)\).  There are
topological embeddings
\begin{align}
 \delta_{12}:\mathcal D_{12}&\hookrightarrow
 \mathcal A\wedge\mathcal A,
 &
 [r,u_1,u_2]&\longmapsto
 [\,[r,1,u_1],[r,1,u_2],],
 \label{cusp-cone-D12-embedding}\\
 \delta_{24}:\mathcal D_{24}&\hookrightarrow
 \mathcal C^{\wedge4},
 &
 [r,\theta_1,\ldots,\theta_4]&\longmapsto
 \bigwedge_{j=1}^4[r,(1,\theta_j)].
 \label{cusp-cone-D24-embedding}
\end{align}
Here \(1\in S^1\) fixes the positive radial base angle.  For every
\(T\geq T_0\), put \(r_T=e^{-2\pi T}\).  The restrictions of the sphere
quotient maps to the synchronized radial fibres are exactly the
cross-sections at height \(r_T\):
\begin{align}
 F_T^2&\xrightarrow{\ \pi_T\times\pi_T\ }
 Q_T^2\cong\{r_T\}\times T^4
 \subset\mathcal D_{12},
 \label{cusp-cone-D12-slice}\\
 F_T^4&\xrightarrow{\ \chi_{\gamma,T}^{\times4}\ }
 C_T^4\cong\{r_T\}\times T^4
 \subset\mathcal D_{24}.
 \label{cusp-cone-D24-slice}
\end{align}
Every one of these embedded four-torus slices is null-homotopic in its
carrier, and hence in the full smash target.

At the cone vertex \(d_\nu\), \(\nu\in\{12,24\}\), the local integral
homology is
\begin{equation}
 H_k(\mathcal D_\nu,
       \mathcal D_\nu\setminus\{d_\nu\};\mathbb Z)
 \cong
 \begin{cases}
  \mathbb Z^4,&k=2,\\
  \mathbb Z^6,&k=3,\\
  \mathbb Z^4,&k=4,\\
  \mathbb Z,&k=5,\\
  0,&\text{otherwise}.
 \end{cases}
 \label{cusp-cone-carrier-local-homology}
\end{equation}
In particular, \(d_\nu\) is not an integral homology-manifold point, with or
without boundary, of any dimension.  Consequently it is not a topological
manifold point, with or without boundary, of any dimension.
\end{proposition}

\begin{proof}
The two displayed domains are cones on
\((\overline T^2)^2\cong T^4\) and \((S^1)^4\cong T^4\).  At \(r=0\),
every factor on the right of
\cref{cusp-cone-D12-embedding,cusp-cone-D24-embedding} is the relevant
smash basepoint, so the formulae descend from the cylinders to the cones.
For \(0<r\leq r_0\), none of the factors is the smash basepoint.  Equality
in the smash quotient is then equality in every factor, and the coordinate
descriptions of the suspension and cusp cone give equality of \(r\) and of
all torus coordinates.  Thus each \(\delta_\nu\) is injective.  Its domain
is compact and its target is Hausdorff, so it is an embedding.

The already proved one-factor restriction formulae, with
\(r_T=e^{-2\pi T}\), are
\[
 q_2(x)=[r_T,1,\iota_T^{-1}(\pi_T(x))]
 \quad(x\in F_T)
\]
and, in the cusp-cone piece of \(\mathcal C\),
\[
 q_1(x)
 =[r_T,(1,\vartheta_T^{-1}(\chi_{\gamma,T}(x)))]
 \quad(x\in F_T).
\]
Applying these identities in two and four factors gives
\cref{cusp-cone-D12-slice,cusp-cone-D24-slice}.  Sliding the cone coordinate
from \(r_T\) to zero contracts each slice to the cone vertex.  Finally,
excision and the long exact sequence of the pair give, for \(L=T^4\),
\[
 H_k(\operatorname{Cone}(L),
     \operatorname{Cone}(L)\setminus\{v\};\mathbb Z)
 \cong\widetilde H_{k-1}(L;\mathbb Z).
\]
The ranks \(4,6,4,1\) of \(H_1,H_2,H_3,H_4\) of \(T^4\) prove
\cref{cusp-cone-carrier-local-homology}.  At an interior point of an
integral homology \(n\)-manifold the local integral homology is \(\mathbb Z\)
in degree \(n\) and zero in every other degree, whereas at a boundary point
of an integral homology manifold with boundary all these local groups
vanish.  The simultaneous nonzero groups in degrees \(2,3,4,5\), including
ranks greater than one, agree with neither pattern for any \(n\).  Every
topological manifold is locally an integral homology manifold, so the final
assertion follows as well.
\end{proof}

\begin{theorem}[Point-set quotient maps and exact homotopy types]
\label{cusp-cone-target-homotopy-types}
The smash maps are continuous surjective quotient maps
\[
 q_2^{\wedge2}:S^{12}\twoheadrightarrow\mathcal A^{\wedge2},
 \qquad
 q_1^{\wedge4}:S^{24}\twoheadrightarrow\mathcal C^{\wedge4}.
\]
Their targets have homotopy types
\begin{align}
 \mathcal A^{\wedge2}
 &\simeq
 \bigvee^9S^4\vee\bigvee^{18}S^5\vee
 \bigvee^{15}S^6\vee\bigvee^6S^7\vee S^8,
 \label{cusp-cone-A-smash-wedge}\\
 \mathcal C^{\wedge4}
 &\simeq
 S^8\vee\bigvee^4S^9\vee\bigvee^6S^{10}
 \vee\bigvee^4S^{11}\vee S^{12}.
 \label{cusp-cone-C-smash-wedge}
\end{align}
Both maps induce zero on every positive-degree integral homology and
cohomology group and have zero ordinary Hurewicz image.  Neither map has a
continuous homotopy section.  These conclusions do not determine whether
the unstable homotopy class of either map is zero.
\end{theorem}

\begin{proof}
Let \(f:(K,k_0)\to(L,\ell_0)\) be a surjective based map between compact
Hausdorff spaces.  Every non-basepoint class of \(L^{\wedge r}\) has a
representative \((y_1,\ldots,y_r)\) with each \(y_i\ne\ell_0\).  Choosing
\(x_i\in f^{-1}(y_i)\) proves surjectivity of \(f^{\wedge r}\); the smash
basepoint is also in its image.  The source smash is compact.  The target
fat wedge is closed, hence the target smash is Hausdorff.  A continuous
surjection from a compact space to a Hausdorff space is a quotient map.
This proves the first assertion for \(q_2\) and \(q_1\).

The suspension splitting
\[
 \Sigma(X_1\times\cdots\times X_n)
 \simeq
 \bigvee_{\varnothing\ne I\subseteq\{1,\ldots,n\}}
 \Sigma\!\left(\bigwedge_{i\in I}X_i\right)
\]
gives
\[
 \mathcal A=\Sigma T^3
 \simeq\bigvee^3S^2\vee\bigvee^3S^3\vee S^4.
\]
Finite smash products distribute over finite wedges.  The ordered degree
sums give multiplicities \(9,18,15,6,1\) in dimensions
\(4,5,6,7,8\), proving \cref{cusp-cone-A-smash-wedge}.

For \(\mathcal C\), write \(x=\rho_1\), \(y=\rho_2\), and
\(h=\eta\).  The pair-of-pants product and the two exact Seifert fillings
give
\[
 \pi_1(\mathcal M)=
 \langle x,y,h\mid[x,h]=[y,h]=1,
 x^3h^{-1}=1,\ y^4h=1\rangle.
\]
The remaining boundary torus is generated by \(h\) and
\(\rho_0=(xy)^{-1}\).  Coning this torus kills their normal closure, and
therefore
\[
 \pi_1(\mathcal C)
 \cong\langle x,y\mid x^3=1,\ y^4=1,\ xy=1\rangle=1,
\]
because \(y=x^{-1}\) and \(\gcd(3,4)=1\).  The previously computed
relative homology gives
\[
 \widetilde H_2(\mathcal C;\mathbb Z)=\mathbb Z,
 \qquad
 \widetilde H_3(\mathcal C;\mathbb Z)=\mathbb Z,
\]
and no other reduced homology.  Since \(\mathcal C\) is a simply connected
finite three-complex, Hurewicz supplies a sphere representing the
\(H_2\)-generator, and the Whitehead exact sequence makes
\(\pi_3(\mathcal C)\to H_3(\mathcal C)\) surjective.  A wedge of
representing maps
\(S^2\vee S^3\to\mathcal C\) is a homology equivalence between simply
connected CW complexes, hence a homotopy equivalence.  Thus
\[
 \mathcal C\simeq S^2\vee S^3.
\]
Taking the fourth smash power yields \cref{cusp-cone-C-smash-wedge} with
the binomial multiplicities \(1,4,6,4,1\).

The first target has no reduced homology above degree eight and the second
has none above degree twelve.  The source spheres have no positive homology
below degrees twelve and twenty-four.  Hence all positive-degree homology
and cohomology homomorphisms induced by the two maps vanish.  Their
Hurewicz images lie in \(H_{12}(\mathcal A^{\wedge2})=0\) and
\(H_{24}(\mathcal C^{\wedge4})=0\), respectively.  A section would split a
nonzero lower homology group of the target into the zero lower homology of
the source sphere, which is impossible.  Higher homotopy groups of the
displayed wedges need not vanish, so none of these arguments decides the
unstable null-homotopy question.
\end{proof}

\begin{theorem}[Exact decompactification spaces with boundary
\(\mathbb R^4\)]
\label{cusp-cone-exact-decompactification-space}
For \(\nu\in\{12,24\}\), set
\[
 A_T^{(12)}=\operatorname{diag}(Y_T,Y_T),
 \qquad
 A_T^{(24)}=\ell_T I_4,
 \qquad s_0=T_0^{-1}.
\]
On \([0,s_0]\times\mathbb R^4\), define an equivalence relation
\(\sim_\nu\) fibrewise over \(s\) by
\begin{align*}
 (0,x)\sim_\nu(0,y)&\iff x=y,\\
 (s,x)\sim_\nu(s,y)&\iff
 x-y\in A_{1/s}^{(\nu)}\mathbb Z^4
 \quad(0<s\leq s_0).
\end{align*}
Let
\[
 \mathfrak E_\nu
 =([0,s_0]\times\mathbb R^4)/\!\sim_\nu,
 \qquad p_\nu:\mathfrak E_\nu\longrightarrow[0,s_0].
\]
Then \(\mathfrak E_\nu\) is a Hausdorff, second-countable topological
five-manifold with boundary, and
\begin{align}
 p_{12}^{-1}(s)&\cong Q_{1/s}^2 &(s>0),&
 &p_{12}^{-1}(0)&\cong\mathbb R^4,
 \label{cusp-cone-E12-fibres}\\
 p_{24}^{-1}(s)&\cong C_{1/s}^4 &(s>0),&
 &p_{24}^{-1}(0)&\cong\mathbb R^4.
 \label{cusp-cone-E24-fibres}
\end{align}
The other boundary component is the compact four-torus over \(s=s_0\).

Give every positive fibre its Euclidean quotient metric and the zero fibre
its Euclidean metric.  For every \(R>0\), there is \(s_R>0\) such that,
for \(0\leq s<s_R\), writing
\(q_{\nu,s}(x)=[s,x]\), the quotient map restricts to an exact isometry
\[
 q_{\nu,s}:B_R^{\mathbb R^4}(0)
 \xrightarrow{\ \cong_{\mathrm{isom}}\ }
 B_R^{p_\nu^{-1}(s)}([s,0])
\]
onto the entire quotient radius-\(R\) ball, not merely onto an embedded
subset of that ball.
Consequently the positive fibres converge to the zero fibre in pointed
local-isometric, pointed smooth, and pointed Gromov--Hausdorff topology.

There are continuous surjective comparison maps which set-theoretically
collapse the zero boundary fibres,
\[
 \begin{tikzcd}[column sep=large]
 \mathbb R^4=p_{12}^{-1}(0)
 \arrow[r,hook]&
 \mathfrak E_{12}
 \arrow[r,"\beta_{12}"]&
 \mathcal D_{12}
 \arrow[r,hook,"\delta_{12}"]&
 \mathcal A^{\wedge2},
 \end{tikzcd}
\]
and
\[
 \begin{tikzcd}[column sep=large]
 \mathbb R^4=p_{24}^{-1}(0)
 \arrow[r,hook]&
 \mathfrak E_{24}
 \arrow[r,"\beta_{24}"]&
 \mathcal D_{24}
 \arrow[r,hook,"\delta_{24}"]&
 \mathcal C^{\wedge4}.
 \end{tikzcd}
\]
For \(s>0\), \(\beta_\nu\) is the specified marked homeomorphism from
\(p_\nu^{-1}(s)\) to the radial torus slice at \(r=e^{-2\pi/s}\).  Its
restriction
\[
 \beta_\nu:\mathfrak E_\nu\setminus p_\nu^{-1}(0)
 \xrightarrow{\ \cong\ }\mathcal D_\nu\setminus\{d_\nu\}
\]
is a homeomorphism, while
\[
 \beta_\nu^{-1}(d_\nu)=p_\nu^{-1}(0)\cong\mathbb R^4.
\]
Neither \(\beta_\nu\) is proper or a quotient map.  It factors through a
continuous bijection
\[
 \overline\beta_\nu:
 \mathfrak E_\nu/p_\nu^{-1}(0)\longrightarrow\mathcal D_\nu,
\]
but this bijection is not a homeomorphism; equivalently, the topology of
\(\mathcal D_\nu\) is not the quotient topology induced by \(\beta_\nu\).
Thus the limiting \(\mathbb R^4\) is an actual boundary fibre of
\(\mathfrak E_\nu\), whereas the continuous comparison map sends that entire
fibre to the cone vertex.
\end{theorem}

\begin{proof}
We give a common proof.  Write
\[
 L_s=A_{1/s}^{(\nu)}\mathbb Z^4\quad(s>0),
 \qquad L_0=\{0\}.
\]
The matrices vary continuously for \(s>0\), and
\[
 \operatorname{sys}(L_s)\longrightarrow\infty
 \qquad(s\downarrow0)
\]
by \(\operatorname{sys}(Y_T\mathbb Z^2)\geq T-C\) in the
\(S^{12}\) case and \(\operatorname{sys}(\ell_T\mathbb Z^4)=\ell_T\)
in the \(S^{24}\) case.

The equivalence relation is closed.  Indeed, suppose
\((s_j,x_j)\sim_\nu(s_j,y_j)\) and
\((s_j,x_j,y_j)\to(s,x,y)\).  If \(s>0\), write
\(y_j-x_j=A_{1/s_j}^{(\nu)}n_j\).  On a compact neighborhood of \(s\),
the inverse matrices are uniformly bounded, so boundedness of
\(y_j-x_j\) forces \(n_j\) to lie in a finite set.  Passing to a
subsequence makes \(n_j\) constant and proves \(x\sim_\nu y\).  If
\(s=0\), any nonzero \(n_j\) would give
\[
 \|y_j-x_j\|\geq\operatorname{sys}(L_{s_j})\longrightarrow\infty,
\]
contrary to convergence.  Hence \(x=y\).

The quotient projection \(q_\nu:[0,s_0]\times\mathbb R^4\to
\mathfrak E_\nu\) is open.  If \(U\) is open, its saturation is the union
of \(U\) with the images, for \(n\in\mathbb Z^4\), of
\[
 U\cap((0,s_0]\times\mathbb R^4)
 \quad\text{under}\quad
 (s,x)\longmapsto(s,x+A_{1/s}^{(\nu)}n),
\]
and every one of these images is open.  A closed equivalence relation with
an open quotient map has Hausdorff quotient: two inequivalent points admit
a product neighborhood disjoint from the relation, and the saturations of
the two factor neighborhoods are disjoint open sets.  Since the source is
second countable and the quotient map is open, images of a countable base
form a countable base of the quotient.

Fix \(x_0\in\mathbb R^4\).  Choose \(a>0\) and \(s_1>0\) such that
\(\operatorname{sys}(L_s)>2a\) for \(0<s<s_1\).  The quotient projection
is one-to-one on
\([0,s_1)\times B_a^\circ(x_0)\), where
\(B_a^\circ(x_0)=\{x:\|x-x_0\|<a\}\), and its restriction is an open continuous
bijection onto an open subset of \(\mathfrak E_\nu\); it is therefore a
homeomorphism.  These are half-space charts at \(s=0\).  The usual evenly
covered small balls give product charts for \(s>0\), and the same argument
at \(s=s_0\) gives the second boundary component.  This proves the
manifold assertion and the fibre identifications.

For the metric assertion, choose \(s_R\) so that
\(\operatorname{sys}(L_s)>4R\) for \(s<s_R\).  If
\(x,y\in B_R(0)\) and \(0\ne\lambda\in L_s\), then
\[
 \|x-y+\lambda\|
 \geq\|\lambda\|-\|x-y\|
 >4R-2R\geq\|x-y\|.
\]
Thus no translated representative shortens the Euclidean distance on that
ball, so \(q_{\nu,s}\) is an exact isometric embedding there.  It is also
onto the full quotient ball.  Indeed, if \(\xi=[s,x]\) satisfies
\(d([s,0],\xi)\leq R\), discreteness of \(L_s\) makes the infimum in
\[
 d([s,0],[s,x])=\min_{\lambda\in L_s}\|x+\lambda\|
\]
attained.  For a minimizing \(\lambda_0\), the representative
\(x+\lambda_0\) lies in \(B_R^{\mathbb R^4}(0)\) and maps to \(\xi\).  At
\(s=0\), \(q_{\nu,0}\) is the identity.  This proves the asserted onto
isometry, and since it holds for every fixed \(R\), it proves all three
pointed convergences.

It remains to define these comparison maps, which set-theoretically collapse
the zero fibres.  For \(s>0\), put
\(T=1/s\) and \(r=e^{-2\pi T}\).  In the \(S^{12}\) case, write
\(x=(x_1,x_2)\in\mathbb R^2\oplus\mathbb R^2\) and set
\[
 \beta_{12}([s,x])
 =\bigl[r,
 \iota_T^{-1}([x_1]_{Y_T\mathbb Z^2}),
 \iota_T^{-1}([x_2]_{Y_T\mathbb Z^2})\bigr].
\]
In the \(S^{24}\) case, write
\(x=(x_1,\ldots,x_4)\) and set
\[
 \beta_{24}([s,x])
 =\bigl[r,
 [x_1/\ell_T],\ldots,[x_4/\ell_T]\bigr].
\]
These formulae are well defined on the positive fibres and are continuous
there.  Define \(\beta_\nu([0,x])=d_\nu\) for every \(x\in\mathbb R^4\).
Every neighborhood of \(d_\nu\) in the cone contains all cross-sections
with sufficiently small \(r\).  Since \(e^{-2\pi/s}\to0\) as
\(s\downarrow0\), continuity at the zero fibre is uniform in \(x\).
Surjectivity follows because each positive fibre maps onto its complete
cross-section and the zero fibre maps onto the vertex.  The restriction away
from \(p_\nu^{-1}(0)\) is a homeomorphism because
\(s\mapsto e^{-2\pi/s}\) is a homeomorphism from \((0,s_0]\) to
\((0,r_0]\), and the displayed marked-torus identifications are fibrewise
homeomorphisms varying continuously with \(s\).  The vertex fibre is exactly
\(p_\nu^{-1}(0)\), so \(\beta_\nu\) is not proper.

It is also not a quotient map.  Choose
\(a\in\mathbb R^4\setminus\mathbb Z^4\), for example
\(a=(1/2,1/2,1/2,1/2)\), choose \(s_j\downarrow0\), and put
\[
 z_j=[s_j,A_{1/s_j}^{(\nu)}a]\in\mathfrak E_\nu.
\]
For \(\nu=12\),
\[
 \inf_{n\in\mathbb Z^4}
 \|A_{1/s_j}^{(12)}(a+n)\|
 \geq(1/s_j-C)\operatorname{dist}(a,\mathbb Z^4)
 \longrightarrow\infty,
\]
and for \(\nu=24\) the same infimum equals
\[
 \ell_{1/s_j}\operatorname{dist}(a,\mathbb Z^4)
 \longrightarrow\infty.
\]
More explicitly, for every fixed \(x_0\in\mathbb R^4\),
\[
 \inf_{n\in\mathbb Z^4}
 \|A_{1/s_j}^{(\nu)}(a+n)-x_0\|
 \geq
 \inf_{n\in\mathbb Z^4}\|A_{1/s_j}^{(\nu)}(a+n)\|-\|x_0\|
 \longrightarrow\infty.
\]
Hence \(\{z_j:j\geq1\}\) is closed in \(\mathfrak E_\nu\): it has no
accumulation in a positive fibre and the last estimate shows that it
eventually misses every half-space chart about every point of the zero
fibre.  But
\[
 \beta_\nu(z_j)=[e^{-2\pi/s_j},a]\longrightarrow d_\nu,
\]
so \(\{\beta_\nu(z_j):j\geq1\}\) is not closed in \(\mathcal D_\nu\).
Since \(\beta_\nu\) is injective off the zero fibre and the latter maps to
\(d_\nu\), the inverse image of this nonclosed set is the closed set
\(\{z_j:j\geq1\}\).  Thus \(\beta_\nu\) is not a quotient map.  This proves
the two diagrams and their exact topological scope.  Since \(\beta_\nu\) is
constant exactly on \(p_\nu^{-1}(0)\) and otherwise injective, it induces the
continuous bijection \(\overline\beta_\nu\) displayed in the theorem.  Were
that bijection a homeomorphism, \(\beta_\nu\) would be the composite of the
quotient projection with a homeomorphism and hence a quotient map, contrary
to the counterexample.  Therefore \(\overline\beta_\nu\) is not a
homeomorphism.
\end{proof}

\begin{corollary}[Sharp global meaning of the Euclidean limit]
\label{cusp-cone-sharp-global-R4-scope}
The native routes constructed above exhibit the underlying space
\(\mathbb R^4\) in the following three precise senses:
\begin{enumerate}
\item as each of the nonclosed topologically embedded open cells
\(U_{12,T}\subset S^6\) and \(U_{24,T}\subset S^8\);
\item as the pointed intrinsic Euclidean limit of the four-torus slices in the
global quotient carriers of \(S^{12}\) and \(S^{24}\); and
\item as the Euclidean decompactification boundary fibre of
\(\mathfrak E_{12}\) and \(\mathfrak E_{24}\).
\end{enumerate}
It does not produce \(\mathbb R^4\) as the whole compact sphere, as the
whole compact conical target, or as a continuous quotient target of either
compact space.  Nor does the unrescaled restriction of any fixed compatible
ambient target metric have \(\mathbb R^4\) as its slice limit, and the open
cells in item \(1\) are not globally Euclidean-isometric in such a metric.
\end{corollary}

\begin{proof}
The three positive assertions are
\cref{cusp-cone-literal-R4-open-cells,cusp-cone-explicit-carriers,%
cusp-cone-exact-decompactification-space}.
For the negative assertions, a continuous image of a compact sphere or
compact conical target is compact, so it cannot be all of \(\mathbb R^4\).
The local homology calculation
\cref{cusp-cone-carrier-local-homology} shows that the carrier vertex is not
an integral homology-manifold point, and hence not a manifold point, of any
dimension, with or without boundary.  Finally, in every fixed compatible
metric on a compact cone target, all sufficiently deep radial slices lie
uniformly in an arbitrarily small neighborhood of the vertex, so their ambient
diameters tend to zero.  Their intrinsic diameters tend to infinity because
their intrinsic systoles do.  Hence without a \(T\)-dependent rescaling no
fixed ambient metric realizes the intrinsic Euclidean limit on those slices.
The fixed cusp-adapted blow-ups are treated separately in
\cref{sphere-metric-fixed-realization}.  The comparison diagrams record the
unrescaled distinction exactly: they collapse the new \(\mathbb R^4\)
boundary rather than identify it with the compact target.
\end{proof}

\paragraph{Scope.}
The constructions above use only the original, unscaled fibre metrics and
the quotient maps derived from the native \(S^6\) cusp.  They do not use the
auxiliary finite-cover unspooling.  The conical carriers depend on the
chosen cusp coordinate, positive radial ray, and based identification
\(X\cong S^6\).  The spaces \(\mathfrak E_\nu\) additionally remember the
marked Euclidean quotient metrics.  None of these data defines a compact
complex structure on \(S^{12}\) or \(S^{24}\), a four-dimensional physical
spacetime, or an interacting Yang--Mills Hamiltonian.


\subsection{The full \texorpdfstring{\(S^{12}\to S^6\)}{S12-to-S6}
carrier is null-homotopic}
\label{cusp-s12-carrier-homotopy}

This subsection determines the based homotopy class of the full two-copy
sphere carrier.  It is important that the calculation concerns the full
quotient map, not merely the inclusion of one of its moving four-spheres.

Recall the exact maps from
\cref{cusp-cone-S12-route,cusp-cone-angle-forgotten-sphere-carriers}.  The
cusp quotient is
\[
 q_2:(X,x_0)\longrightarrow
 \mathcal A=\operatorname{Susp}(S^1\times\overline T^2),
\]
and, for the coordinate one-skeleton
\[
 K^{(1)}=
 ((\mathbb R/\mathbb Z)\times\{0\})
 \cup(\{0\}\times(\mathbb R/\mathbb Z))\subset\overline T^2,
\]
the degree-one quotient
\(c:\overline T^2\to\overline T^2/K^{(1)}\cong S^2\) defines
\[
 \widehat p_2([r,e^{i\phi},u])=[r,c(u)]
 \in\Sigma_r S^2\cong S^3.
\]
Thus
\[
 g:=\widehat q_3=\widehat p_2\circ q_2:(X,x_0)\longrightarrow(S^3,*)
\]
is the specified continuous surjective based quotient.  After the fixed
based identification \(X\cong_{\mathrm{diff}}S^6\), its smash square is the
specified continuous surjective based quotient
\[
 g^{\wedge2}:S^{12}=S^6\wedge S^6
 \longrightarrow S^3\wedge S^3=S^6.
\]

\begin{theorem}[The full external smash pairing is zero]
\label{cusp-s12-full-carrier-null}
For every pair of based maps \(x,y:S^6\to S^3\),
\[
 [x\wedge y]=0\qquad\text{in}\qquad\pi_{12}(S^6).
\]
In particular,
\(\widehat q_3^{\wedge2}:S^{12}\twoheadrightarrow S^6\) is simultaneously
surjective as a point-set quotient and based null-homotopic.
\end{theorem}

\begin{proof}
Write \(E\) for reduced suspension, using the ordered smash coordinates on
spheres.  The external smash product has the exact factorization
\begin{equation}
 x\wedge y
 =
 (x\wedge\operatorname{id}_{S^3})
 \circ
 (\operatorname{id}_{S^6}\wedge y):
 S^6\wedge S^6\longrightarrow S^6\wedge S^3
 \longrightarrow S^3\wedge S^3.
 \label{cusp-s12-smash-factorization}
\end{equation}
Under the displayed ordered sphere identifications, the two factors, in
order of application, represent \(E^6y\in\pi_{12}(S^9)\) and
\(E^3x\in\pi_9(S^6)\).  Equivalently, Toda's
Barratt--Hilton reduced-join formula gives, with
\(p=q=k=h=3\),
\begin{equation}
 [x\wedge y]=(E^3[x])\circ(E^6[y]).
 \label{cusp-s12-toda-smash-formula}
\end{equation}
Here Toda's sign is \((-1)^{(p+k)h}=(-1)^{18}=+1\).  Reversing the
two displayed factors introduces the separate twist sign
\((-1)^{ph}=-1\); no such reversal is used here.

The three exact groups needed below are
\begin{equation}
 \pi_6(S^3)\cong\mathbb Z/12,
 \qquad
 \pi_{12}(S^9)\cong\mathbb Z/24,
 \qquad
 \pi_{12}(S^6)\cong\mathbb Z/2.
 \label{cusp-s12-three-stem-groups}
\end{equation}
The first and third are the integral low unstable groups displayed in
Hatcher's table extracted from Toda; the middle group is in the stable
three-stem
\cite[opening table, p.~339]{Hatcher2002};
\cite[Table~1, p.~113]{IsaksenWangXu2023}.

Put \(\beta=E^6[y]\in\pi_{12}(S^9)\).  Since
\(12[y]=0\), functoriality of suspension gives \(12\beta=0\).  In the
cyclic group \(\mathbb Z/24\), the elements annihilated by \(12\) are
exactly the even elements.  Hence there is
\(\gamma\in\pi_{12}(S^9)\) such that \(\beta=2\gamma\).
Postcomposition with \(E^3[x]:S^9\to S^6\) induces a group homomorphism
\[
 (E^3[x])_*:\pi_{12}(S^9)\longrightarrow\pi_{12}(S^6).
\]
The target has exponent two, so \cref{cusp-s12-toda-smash-formula} gives
\[
 [x\wedge y]
 =(E^3[x])_*(\beta)
 =2(E^3[x])_*(\gamma)
 =0.
\]
This proves the vanishing of the entire external smash pairing and hence
the theorem.
\end{proof}

\begin{corollary}[All higher smash powers]
\label{cusp-s12-all-higher-smashes-null}
For every based \(g:S^6\to S^3\) and every \(r\geq2\),
\[
 g^{\wedge r}:S^{6r}\longrightarrow S^{3r}
\]
is based null-homotopic.
\end{corollary}

\begin{proof}
With the fixed associativity homeomorphisms for smash products,
\[
 g^{\wedge r}=g^{\wedge2}\wedge g^{\wedge(r-2)}.
\]
Smashing a based null-homotopy with a based map gives a based
null-homotopy.
\end{proof}

\begin{remark}[The one-copy invariant and its computed value]
\label{cusp-s12-one-copy-invariant}
The proof does not need to determine the one-copy class
\([g]\in\pi_6(S^3)\cong\mathbb Z/12\).  Toda's suspension calculation
identifies \(\pi_6(S^3)\) with the even subgroup of the stable three-stem;
Hatcher gives an integral Serre-spectral-sequence derivation of the same
injectivity and identifies twice the stable Hopf--\(\nu\) generator with a
generator of the suspended \(\mathbb Z/12\)-summand
\cite[Theorem~5.41(b), p.~573, and pp.~577--578]{HatcherSpectralSequences}.
Thus, after choosing a stable generator \(\nu_9\), there is a unique
\(a\in\mathbb Z/12\) such that
\begin{equation}
 [E^6g]=2a\nu_9\in\pi_{12}(S^9)\cong\mathbb Z/24.
 \label{cusp-s12-sixfold-suspension-class}
\end{equation}
Because \(a\mapsto2a\) embeds \(\mathbb Z/12\) as the even subgroup of
\(\mathbb Z/24\), \(E^6g\) is null-homotopic if and only if \(g\) is
null-homotopic.  Thus the one-copy class and its sixfold suspension are the
same integer modulo \(12\). It is now zero by Theorem~\ref{cusp-null:thm}.

The exact geometric invariant which computes \(a\) is also identifiable.
Choose \(0<r_0<\epsilon\) and
\(u_0\in\overline T^2\setminus K^{(1)}\), and let
\(y_0=[r_0,c(u_0)]\) lie in the open three-cell of
\(\Sigma_rS^2\).  From the literal formula for \(q_2\),
\begin{equation}
 g^{-1}(y_0)
 =\{x\in f^{-1}(|t|=r_0):\overline\pi_f(x)=u_0\}.
 \label{cusp-s12-regular-inverse-image}
\end{equation}
The fibre of \(\overline\pi\) is the two-torus defined by
\(K_\Lambda=\mathbb Z\langle\widehat w,\widehat\delta\rangle\), and the
cusp monodromy fixes both \(\widehat w\) and \(\widehat\delta\).  Therefore
\eqref{cusp-s12-regular-inverse-image} is the mapping torus of an affine translation of the kernel torus,
which is diffeomorphic to \(S^1\times T^2\cong T^3\): if the translation is
\(k\mapsto k+v\), choose a real lift \(\widetilde v\); the map
\([\theta,k]\mapsto(e^{2\pi i\theta},k+\theta\widetilde v)\) respects
\((1,k)\sim(0,k+v)\) and has the inverse obtained by subtracting
\(\theta\widetilde v\). After choosing a smooth chart in the open
target cell, the differential of \(g\) supplies its normal three-framing.
Pontryagin--Thom identifies the stabilized class of \(g\) with the framed
bordism class of this specifically framed \(T^3\).  Since stabilization on
\(\pi_6(S^3)\) is injective, computing that framing class computes \(a\)
exactly.  The unframed diffeomorphism type \(T^3\), point-set surjectivity,
and the zero ordinary homology map do not compute \(a\).
\end{remark}

\begin{remark}[Scope]
\label{cusp-s12-nullity-scope}
The theorem changes no point-set or metric statement.  A surjective quotient
map may be null-homotopic, and the embedded moving four-spheres and their
intrinsic pointed \(\mathbb R^4\)-limits remain exactly as constructed.  The
theorem concerns the sphere-valued postcomposition
\(\widehat q_3^{\wedge2}\); it does not imply null-homotopy of the earlier
conical-target map
\(q_2^{\wedge2}:S^{12}\to\mathcal A^{\wedge2}\).
\end{remark}

\subsection{The original cusp coefficient: an explicit framed nullbordism}


Let \(X\) be the smooth compact source of the original construction,
\(f:X\to B\) its map, and \(t\) the specified cusp coordinate.
Put \(W=f^{-1}(0)\), choose the original \(x_0\in W\), and retain the
ordered quotient lattice basis induced by
\((\widehat\gamma,\widehat u)\), with kernel generated by
\((\widehat w,\widehat\delta)\). Thus the quotient torus is
\(\overline T^2=\mathbb R^2/\mathbb Z^2\). The original periods are
\begin{equation}
 Z(s)=sB_0+C(t),\qquad e^{2\pi i s}=t,\qquad
 B_0=\begin{pmatrix}0&1\\-1&0\end{pmatrix},\qquad
 C(t)=\begin{pmatrix}6\mu(t)&h(t)\\b(t)-h(t)&\mu(t)\end{pmatrix}.
 \label{cusp-null:eq:periods}
\end{equation}
Every entry of \(C\) is holomorphic at \(t=0\).
The dense torus of the original toric cusp has coordinates
\((x_1,x_2,t)\), with
\(x=\exp(2\pi i\zeta)\). The residual lattice action is exactly
\begin{equation}
 \Psi_\lambda(x,t)=
 \bigl(t^{B_0\lambda}\exp(2\pi iC(t)\lambda)x,t\bigr),
 \qquad\lambda\in\mathbb Z^2.
 \label{cusp-null:eq:action}
\end{equation}
It extends freely and properly discontinuously across the original
unimodular toric filling. In particular the quotient projection is a local
biholomorphism at every toric fixed point. These are the reconstructed
data and proved local properties of the original source
(the reconstructed original manuscript, \S4, source pp.~24--29; its complete cusp data are displayed here).

Let
\[
 K^{(1)}=(\mathbb R/\mathbb Z\times\{0\})\cup(\{0\}\times\mathbb R/\mathbb Z),\qquad
 c:\overline T^2\longrightarrow \overline T^2/K^{(1)}\cong S^2
\]
be the specified degree-one cell quotient. Its restriction to the open
two-cell is taken with the smooth coordinates used in the original
regular-fibre description. Both radial endpoints and the complete
basepoint cylinder are collapsed in
\[
 \Sigma_r(S^2,*)=
 \frac{[0,\epsilon]\times S^2}
 {(\{0,\epsilon\}\times S^2)\cup([0,\epsilon]\times\{*\})}
 \cong S^3.
\]
The original based map is
\begin{equation}
 g(x)=[|t|,c(\overline\pi_f(x))]\quad(0<|t|<\epsilon),\qquad
 g(W)=g(X\setminus\operatorname{int}N_0)=*.
 \label{cusp-null:eq:g}
\end{equation}

\begin{theorem}[The original one-copy cusp map is null-homotopic]
\label{cusp-null:thm}
The based map \(g:(X,x_0)\to(S^3,*)\) in \eqref{cusp-null:eq:g} is based
null-homotopic. More precisely, a regular inverse image with its actual
ordered normal framing induced by \((r,c_1,c_2)\) is framed nullbordant
inside one original cusp chart times \([0,1]\), disjoint from
\(\{x_0\}\times[0,1]\).
\end{theorem}

\begin{proof}
The proof constructs the complete inverse image, compares its actual
framing with an explicitly extendible framing, and gives the extension.

\paragraph{1. The exact complete inverse image.}
Retain the lift
\[
 u^0=\begin{pmatrix}-1/4\\1/4\end{pmatrix},\qquad
 \overline u^0=(3/4,1/4)\pmod{\mathbb Z^2},\qquad
 B_0u^0=\begin{pmatrix}1/4\\1/4\end{pmatrix}.
\]
Thus \(\overline u^0\notin K^{(1)}\), and \(B_0u^0\) lies in the
interior of the original triangle \(\operatorname{conv}(0,e_1,e_2)\).
Use its original unimodular chart
\begin{equation}
 z_0=t/(x_1x_2),\qquad z_1=x_1,\qquad z_2=x_2,
 \qquad t=z_0z_1z_2.
 \label{cusp-null:eq:chart}
\end{equation}
Set
\[
 \ell=\log|t|<0,\qquad R(t)=-2\pi\operatorname{Im}C(t),\qquad
 A(t)=\ell B_0+R(t).
\]
For a representative \(\zeta=Z(s)u+k\), where \(u,k\in\mathbb R^2\),
\(\operatorname{Im}s=-\ell/(2\pi)\) gives the exact identity
\begin{equation}
 \begin{pmatrix}\log|x_1|\\\log|x_2|\end{pmatrix}
 =A(t)u.
 \label{cusp-null:eq:moduli}
\end{equation}
The quotient value is \(u\bmod\mathbb Z^2\).
Since \(B_0\) is orthogonal, for all \(v\in\mathbb R^2\),
\begin{equation}
 \|A(t)v\|\geq (|\ell|-\|R(t)\|)\|v\|.
 \label{cusp-null:eq:invertibility}
\end{equation}
The boundedness of \(R\) therefore makes \(A(t)\) invertible on a
sufficiently small punctured disc.

For any point of the fibre quotient over \(\overline u^0\), subtracting
the appropriate original period, equivalently applying
\(\Psi_{-\lambda}\), gives a representative with the exact lift
\(u=u^0\). This representative is unique: another residual translate
has logarithmic moduli differing by \(A(t)\lambda\), which vanishes
only for \(\lambda=0\). Conversely every dense-torus point satisfying
\(\log|x|=A(t)u^0\) has this quotient value, since the remaining real
parts of \(\zeta\) are exactly the kernel coordinates \(k\).
For each fixed \(t\), varying \(k\bmod\mathbb Z^2\) gives all independent
phases of \(x_1,x_2\).

The following change of coordinates uses the full original period units:
\begin{align}
 h_j^0(t)&=\exp\bigl(-2\pi i(C(t)u^0)_j\bigr),\qquad j=1,2,\\
 \mathcal F(z_0,z_1,z_2)
 &=(w_0,w_1,w_2)
 =\left(\frac{z_0}{h_1^0(t)h_2^0(t)},\,
 h_1^0(t)z_1,\,h_2^0(t)z_2\right).
 \label{cusp-null:eq:straightening}
\end{align}
The functions \(h_j^0\) are nowhere-zero holomorphic units, including
at \(t=0\). The product \(w_0w_1w_2\) is exactly \(t\).
The inverse of \(\mathcal F\) is
\[
 (w_0,w_1,w_2)\longmapsto
 \left(w_0h_1^0(t)h_2^0(t),\,
 \frac{w_1}{h_1^0(t)},\,\frac{w_2}{h_2^0(t)}\right),
 \qquad t=w_0w_1w_2.
\]
Thus this is a biholomorphism of the original chart restricted to the
cusp disc, not a replacement of \(C(t)\).

For \(r_0\in(0,\epsilon)\), the entire inverse image of
\(y_0=[r_0,c(\overline u^0)]\) becomes exactly
\begin{equation}
 T_{r_0}=
 \{|w_0|=\rho_0,\ |w_1|=\rho_1,\ |w_2|=\rho_2\},\qquad
 (\rho_0,\rho_1,\rho_2)=
 (r_0^{1/2},r_0^{1/4},r_0^{1/4}).
 \label{cusp-null:eq:torus}
\end{equation}
Indeed, \(\log|h^0(t)|=2\pi\operatorname{Im}C(t)u^0=-R(t)u^0\),
so \eqref{cusp-null:eq:moduli} gives \(\log|w_j|=\ell(B_0u^0)_j\) for
\(j=1,2\); the product identity gives the exponent \(1/2\) for
\(w_0\). Conversely these moduli recover \(u=u^0\) by invertibility
of \(A(t)\). The angular map is explicitly
\[
 (\arg t,\arg w_1,\arg w_2)\longleftrightarrow
 (\arg w_0,\arg w_1,\arg w_2),\qquad
 \arg w_0=\arg t-\arg w_1-\arg w_2\pmod{2\pi},
\]
with inverse \(\arg t=\arg w_0+\arg w_1+\arg w_2\).
Its integral matrix has determinant one. This includes the complete
angular mapping torus and all kernel phases, with their original signs.

Let \(o\) be the origin of \eqref{cusp-null:eq:chart}. The original free properly
discontinuous action supplies a neighborhood \(V\) of \(o\) on which
the quotient projection is injective. To see the local injectivity
directly, first choose a relatively compact neighborhood; proper
discontinuity leaves only finitely many translates that can intersect it.
None of those nonidentity translates fixes \(o\). Shrinking finitely
many times separates the neighborhood from all of them. The quotient
map is a local biholomorphism, so its restriction to \(V\) is a chart
of the original source.
Choose a polydisc about \(0\) whose closure lies in \(\mathcal F(V)\).
For sufficiently small \(r_0\), both \(T_{r_0}\) and
\begin{equation}
 \{|w_0|\leq\rho_0,\ |w_1|=\rho_1,\ |w_2|=\rho_2\}
 \label{cusp-null:eq:spatial-bound}
\end{equation}
lie in this polydisc. This proves that the whole inverse image embeds
in one chart and supplies space for the bounding manifold below.

\paragraph{2. The actual normal framing, including its angle dependence.}
Put \(L_j=\log|w_j|\), and use the unique real lift of the quotient
coordinate near \(u^0\). From the exact change of coordinates,
\begin{equation}
 \begin{pmatrix}L_1\\L_2\end{pmatrix}
 =\ell B_0u^0+A(t)(u-u^0),\qquad
 L_0=\ell-L_1-L_2.
 \label{cusp-null:eq:nearby-moduli}
\end{equation}
Differentiating on \(u=u^0\) annihilates the term
\((dA(t))(u-u^0)\). Hence no derivative of \(C(t)\) has been omitted:
the exact differential there is
\begin{equation}
 \begin{pmatrix}dL_0\\dL_1\\dL_2\end{pmatrix}
 =J(t)\begin{pmatrix}d\ell\\du_1\\du_2\end{pmatrix},\quad
 J(t)=
 \begin{pmatrix}
 1/2&-A_{11}-A_{21}&-A_{12}-A_{22}\\
 1/4&A_{11}&A_{12}\\
 1/4&A_{21}&A_{22}
 \end{pmatrix}.
 \label{cusp-null:eq:J}
\end{equation}
Adding the second and third rows to the first gives \((1,0,0)\),
and expansion in that row proves \(\det J(t)=\det A(t)\).

Let \(D=Dc_{\overline u^0}\) be the actual two-by-two derivative in
the ordered open-cell target coordinates \((c_1,c_2)\). It is a fixed
invertible matrix; all its entries and signs are retained. Let
\[
 H=\begin{pmatrix}r_0&0&0\\0&D_{11}&D_{12}\\0&D_{21}&D_{22}\end{pmatrix}.
\]
Let \(n_j\) be the real radial vector \(w_j\) in the \(j\)-th complex
coordinate, with other components zero. Then \(dL_i(n_j)=\delta_{ij}\).
Writing \(N=(n_0,n_1,n_2)\) as the matrix with these vector columns, the actual
\(dg\)-dual normal framing is, in the normal quotient bundle,
\begin{equation}
 \mathfrak n(t)=N J(t)H^{-1}.
 \label{cusp-null:eq:actual-frame}
\end{equation}
Indeed, \(dg(N)=HJ(t)^{-1}\), so multiplication by \(J(t)H^{-1}\)
gives the identity. In the Euclidean metric of this chart the radial
vectors are perpendicular to the product torus; a different source
metric changes their representatives by the canonical projection to its
normal bundle and preserves this normal-quotient framing.

Put \(\ell_0=\log r_0\). Decrease \(r_0\) further, if necessary, so
that \(|\ell_0|>2\sup_{|t|\leq r_0}\|R(t)\|\). For
\(0\leq\sigma\leq1\), define only a homotopy of frames by
\[
 A_\sigma(t)=\ell_0B_0+\sigma R(t),
\]
and let \(J_\sigma(t)\) be \eqref{cusp-null:eq:J} with \(A_\sigma\) in place
of \(A\). For every \(v\in\mathbb R^2\),
\[
 \|A_\sigma(t)v\|\geq
 (|\ell_0|-\|R(t)\|)\|v\|>0\quad(v\ne0).
\]
Thus all matrices are invertible. At \(\sigma=0\),
\begin{equation}
 J_0=
 \begin{pmatrix}
 1/2&\ell_0&-\ell_0\\
 1/4&0&\ell_0\\
 1/4&-\ell_0&0
 \end{pmatrix},\qquad \det J_0=\ell_0^2>0.
 \label{cusp-null:eq:J0}
\end{equation}
Continuity also proves \(\det J_\sigma(t)>0\) for every \(\sigma,t\).
The homotopy \(N J_\sigma(t)H^{-1}\), run from \(\sigma=1\) to
\(\sigma=0\), compares the actual frame with the exact constant-matrix
frame \(N P\), where
\begin{equation}
 P=J_0H^{-1}.
 \label{cusp-null:eq:P}
\end{equation}
The original embedded torus and original functions \(C(t)\) have not
changed. A framing homotopy is a framed cylinder: choose a smooth
parameter function constant near its endpoints and use the displayed
frames on \(T_{r_0}\times[0,1]\).

\paragraph{3. A collared framed nullbordism with the exact constant matrix.}
Use a new time coordinate \(v\in[0,1]\). Define
\begin{equation}
 q(v)=
 \begin{cases}
 0,&0\leq v\leq1/4,\\
 \rho_0^2\exp\bigl(4-1/(v-1/4)\bigr),&1/4<v\leq1.
 \end{cases}
 \label{cusp-null:eq:q}
\end{equation}
This function is smooth: all right derivatives of
\(\exp(-1/(v-1/4))\) at \(1/4\) vanish. It is strictly increasing
for \(v>1/4\), satisfies \(q(1/2)=\rho_0^2\), and
\(q'(1/2)=16\rho_0^2>0\).
In the chart times \([0,1]\), put
\begin{equation}
 \mathcal W=
 \{|w_0|^2+q(v)=\rho_0^2,\quad
       |w_1|^2=\rho_1^2,\quad|w_2|^2=\rho_2^2\}.
 \label{cusp-null:eq:bordism}
\end{equation}
No point has \(v>1/2\). The three defining differentials are independent:
the latter two have nonzero gradients in their separate complex factors;
the first has nonzero \(w_0\)-gradient whenever \(w_0\ne0\), and at
\(w_0=0\) it has the nonzero \(v\)-derivative \(q'(1/2)\).
Thus \(\mathcal W\) is a smooth compact four-manifold. Its only
boundary is \(T_{r_0}\times\{0\}\), and on \(0\leq v\leq1/4\)
it is the exact product collar \(T_{r_0}\times[0,1/4]\).
The north locus \(w_0=0,v=1/2\) is an interior cap: the implicit
function theorem expresses \(v\) smoothly in terms of
\((\operatorname{Re}w_0,\operatorname{Im}w_0)\) there. The first
factor is a disk obtained by rotating its radial profile with this smooth
cap and one boundary circle. Therefore
\(\mathcal W\cong D^2\times T^2\).
Its spatial coordinates lie in \eqref{cusp-null:eq:spatial-bound}, so the complete
manifold lies in the injective original chart.

In the Euclidean chart product metric define
\[
 V_0=(w_0,0,0,q'(v)/2),\qquad
 V_1=(0,w_1,0,0),\qquad V_2=(0,0,w_2,0).
\]
These are one half of the gradients of the three defining equations
in \eqref{cusp-null:eq:bordism}. They are mutually orthogonal and never zero.
For example, if \(F_0=|w_0|^2+q(v)-\rho_0^2\), then
\[
 dF_0(V_0)=2|w_0|^2+\frac12q'(v)^2>0
\]
at every point of \(\mathcal W\). On the entire collar \(q'=0\), so
\((V_0,V_1,V_2)=(N,0)\) exactly. Consequently
\begin{equation}
 (V_0,V_1,V_2)P
 \label{cusp-null:eq:extended-frame}
\end{equation}
is a framing of the normal three-plane bundle of \(\mathcal W\),
and its restriction to the collar is exactly the framing \(NP\) from
\eqref{cusp-null:eq:P}, with every original target-coordinate sign retained.
Concatenate the framing cylinder of the preceding paragraph with this
nullbordism, rescaling the two interval lengths. Their framings agree
exactly on the common product collar, so the concatenation is a smooth
framed nullbordism of the actual inverse image.

\paragraph{4. The based map represented by this framed inverse image.}
Choose the injective chart to exclude \(x_0\) if its center is not
\(x_0\). If its center is \(x_0\), the constructed nullbordism still
excludes \(x_0\), since \(|w_1|=\rho_1>0\) and
\(|w_2|=\rho_2>0\) everywhere on it. Compactness therefore gives
tubular neighborhoods for the framed cylinder and cap disjoint from
\(\{x_0\}\times[0,1]\).

The original continuous quotient map is smooth and a submersion near
its entire inverse image of \(y_0\). First make it constant in a
neighborhood of \(x_0\), by a based homotopy that changes nothing near
that inverse image. Explicitly, take a target chart
\(\kappa:O\to\mathbb R^3\) with \(\kappa(*)=0\), whose image is a ball,
and \(y_0\notin O\). By continuity choose a neighborhood \(U\) of
\(x_0\) with \(g(U)\subset O\), and a smooth bump function
\(\beta\) supported in \(U\) and equal to one near \(x_0\).
For \(x\in U\), replace \(g(x)\) during this homotopy by
\(\kappa^{-1}((1-a\beta(x))\kappa(g(x)))\), \(0\leq a\leq1\),
and leave \(g\) fixed off \(U\). Convexity of the target chart image
makes the formula well-defined, the vanishing of \(\beta\) near the
boundary of its support proves continuity of the pasted formula, and
\(\kappa(g(x_0))=0\) makes the homotopy based. It never introduces
an inverse image of \(y_0\). Its endpoint is smooth near \(x_0\) and
agrees with \(g\) near the entire old inverse image.

Denote the endpoint of the preceding homotopy by \(g_b\). Choose one
closed tubular neighborhood \(U_1\) of the inverse image, contained in
the smooth submersion region and with the entire inverse image contained
in \(\operatorname{int}U_1\). The compact set
\(X\setminus\operatorname{int}U_1\) has image under \(g_b\) missing
\(y_0\). In a fixed target metric put
\[
 d_0=\min_{x\in X\setminus\operatorname{int}U_1}
       d(g_b(x),y_0)>0.
\]
Relative smooth approximation, fixed on the whole closed set \(U_1\)
and on a neighborhood of \(x_0\), gives a smooth based representative
\(g_{\mathrm{sm}}\) whose error on
\(X\setminus\operatorname{int}U_1\) is less than \(d_0/2\).
It therefore has the same inverse image and differential on \(U_1\),
and \(d(g_{\mathrm{sm}}(x),y_0)>d_0/2\) on its complement. Thus there
is no uncontrolled intervening region and no additional inverse image.
Make the approximation also closer than a fixed target injectivity
radius; its short geodesics give the based homotopy.

The relative Pontryagin construction now applies to the explicit
framed nullbordism. It produces a based homotopy to a map missing
\(y_0\); stereographic contraction of \(S^3\setminus\{y_0\}\cong\mathbb R^3\)
to the image of \(*\) is based and contracts that last map.
For completeness, the dependence on the framed inverse image in this
argument is precisely the local product-neighborhood construction and
the interpolation lemma of \cite[\S7, Theorem B and Lemma 4,
pp.~43, 48--50]{CuspNullMilnor1965}. Their relative versions fix the source
basepoint because all tubular modifications are away from it and
stereographic interpolation between equal values keeps that value
fixed. A framed cylinder realizes a framing homotopy; the empty
Pontryagin manifold realizes a map missing the chosen regular value.
This proves the claimed based null-homotopy of the exact map \(g\).
\end{proof}

\begin{corollary}
Under the original fixed based identification \(X\cong S^6\), the
one-copy coefficient is
\[
 [g]=0\in\pi_6(S^3)\cong\mathbb Z/12,\qquad a=0\pmod{12},\qquad
 [E^6g]=0\in\pi_{12}(S^9)\cong\mathbb Z/24.
\]
The previously defined homomorphism
\(\Theta_g:\mathbb Z/12\to\pi_6(S^3)\), \(\Theta_g([n])=n[g]\), and
its composite with the marked peripheral character are zero.
\end{corollary}
\begin{proof}
The theorem states that the actual based map is null-homotopic.
Precomposition with the fixed based identification preserves a based
null-homotopy. Suspension preserves it as well. Thus the unique
coefficient of \([g]\) in any generator of \(\mathbb Z/12\) is zero.
For every \([n]\), \(\Theta_g([n])=n\cdot0=0\), and every
composite with this homomorphism is zero.
\end{proof}

\paragraph{Exact scope.}
The proof computes this sphere-valued postcomposition. It does not
compute the homotopy class of the earlier conical-target map
\(q_2:X\to\operatorname{Susp}(S^1\times\overline T^2)\), and it
does not assert a null-homotopy for that different target. Point-set
surjectivity of \(g\) and all of its stated metric restrictions remain
properties of the original map: a homotopy-class computation does not
alter the defining formula \eqref{cusp-null:eq:g}.
The retained-angle companion is proved separately in
Theorem~\ref{h24:ret:main}: $p_3q_2:X\to S^4$ represents
$\eta_4\eta_5\ne0$, so $q_2$ is nonnull. Its Hopf smash
represents $12\nu_6\ne0$ in $\pi_9(S^6)$. These positive
computations use the actual four-framed inverse torus of that different
postcomposition; they do not reverse or contradict the preceding
null-homotopy through $\widehat p_2$.


\begin{thebibliography}{9}
\bibitem{CuspNullMilnor1965}
J.~W. Milnor, \emph{Topology from the Differentiable Viewpoint},
University Press of Virginia, 1965, \S7, Theorem B and Lemma 4,
pp.~42--50. \url{https://www.maths.ed.ac.uk/~aar/papers/milnortop.pdf}.
\end{thebibliography}


\subsection{Fixed-target metric realization and its sharp dimensional scope}
\label{sphere-metric-limit-note}

This note concerns the codomains of the three based maps
\[
 q_2^{\wedge 2}:X^{\wedge 2}\cong S^{12}
   \longrightarrow Z_{12}:=\mathcal A^{\wedge 2},
\]
\[
 q_{1,\mathrm{cusp}}^{\wedge 4}:X^{\wedge 4}\cong S^{24}
   \longrightarrow Z_{24}^{\mathcal B}:=\mathcal B^{\wedge 4},
 \qquad
 q_1^{\wedge 4}:X^{\wedge 4}\cong S^{24}
   \longrightarrow Z_{24}^{\mathcal C}:=\mathcal C^{\wedge 4}.
\]
In particular, \(S^{12}\) and \(S^{24}\) are the sources, not the spaces
in which the metric limits occur.  We construct one fixed compatible
length metric on each codomain which realizes the native
four-dimensional flat limits on the target slices.  The realization is
sharp: after the same rescaling, the full ambient limits have dimensions
\(8\) and \(12\), and the collapsed smash base point goes to infinity.

Put
\[
 r_T=e^{-2\pi T},\qquad
 J=\begin{pmatrix}0&1\\-1&0\end{pmatrix},\qquad
 Y_T=TJ+E_T,
\]
where \(E_T\) is the bounded matrix in the radial cusp expansion.  Choose
and fix \(T_0\) and \(C>0\) so that
\(\|E_T\|_{\mathrm{op}}\leq C\) for every \(T\geq T_0\), in the fixed
Euclidean coordinates used below.  This notation avoids confusing the
error matrix with the quotient circle \(C_T\) used below.  On the
fixed torus
\[
 \overline T^2=(\overline\Lambda\otimes\mathbb R)/\overline\Lambda
              \cong\mathbb R^2/\mathbb Z^2
\]
let
\[
 h_T:=\iota_T^*g_{Q_T}=dx^{\mathsf T}Y_T^{\mathsf T}Y_Tdx.
\]
Thus
\begin{equation}
 T^{-2}h_T\longrightarrow h_\infty:=dx^{\mathsf T}dx
 \label{sphere-metric-normalized-Q-limit}
\end{equation}
in every \(C^k\) norm in the fixed flat coordinates.  Indeed, the cusp
coefficients are holomorphic in \(t=e^{-2\pi T}\), so their \(T\)-derivatives
are exponentially small, while \(T^{-1}Y_T=J+T^{-1}E_T\).

Use normalized angular coordinates
\[
 u,\theta\in\mathbb R/\mathbb Z;
\]
\(u\) is the argument of the base coordinate and \(\theta\) is the value of
the invariant character.  Increase this same \(T_0\) once more, if
necessary, so that all radial cusp formulas used below hold for
\(T\geq T_0\) and \(e^{-2\pi T_0}<\epsilon\), and put
\[
 r_0=e^{-2\pi T_0}.
\]
For \(0<r\leq r_0\), write
\[
 T(r):=-\frac{1}{2\pi}\log r.
\]

\paragraph{Target-slice notation.}
We use the embeddings
\(j_T,k_T^{\mathcal B},k_T^{\mathcal C}\) defined in
\cref{cusp-cone-target-slice-embeddings}; the normalized angle \(u=0\) used
below is the same positive radial base angle denoted by \(1\in S^1\) there.

\begin{theorem}[Topology-forced collapse for every fixed compatible metric]
\label{sphere-metric-uniform-collapse}
Let \(D_{12},D_{24}^{\mathcal B},D_{24}^{\mathcal C}\) be arbitrary fixed
metrics compatible with the respective compact target topologies, and let
\(*\) denote the corresponding smash base point.  Then
\begin{align}
 d_H^{D_{12}}\bigl(j_T(Q_T^2),\{*\}\bigr)&\longrightarrow0,\\
 d_H^{D_{24}^{\mathcal B}}
  \bigl(k_T^{\mathcal B}(C_T^4),\{*\}\bigr)&\longrightarrow0,\\
 d_H^{D_{24}^{\mathcal C}}
  \bigl(k_T^{\mathcal C}(C_T^4),\{*\}\bigr)&\longrightarrow0.
 \label{sphere-metric-Hausdorff-collapse}
\end{align}
Consequently, with the metrics restricted from any one fixed ambient
target, all three target slices converge in Gromov--Hausdorff topology to
a point, not to \(\mathbb R^4\).
\end{theorem}

\begin{proof}
We prove the first assertion; the other two are identical.  Let
\[
 G:[0,r_0]^2\times(T^3)^2\longrightarrow Z_{12}
\]
be the product of the two inner-cone quotient maps followed by the smash
quotient.  It is continuous and
\(G(0,0,z_1,z_2)=*\) for every \((z_1,z_2)\in(T^3)^2\).  Since the domain
is compact, uniform continuity with respect to \(D_{12}\) gives
\[
 \sup_{(z_1,z_2)\in(T^3)^2}
 D_{12}\bigl(G(r,r,z_1,z_2),*\bigr)\longrightarrow0.
\]
The set \(j_T(Q_T^2)\) is contained in such a cross-section, which proves
the first Hausdorff convergence.  Replacing \(T^3\) by \(T^2\) and taking
four factors proves the other two.  Hausdorff convergence to a singleton
implies Gromov--Hausdorff convergence to that singleton.
\end{proof}

The preceding theorem is unrescaled.  It does not prevent a blow-up of one
fixed metric.  The following construction gives the strongest natural
positive statement.

\begin{theorem}[One fixed cusp-adapted metric realizes the vertical
\(\mathbb R^4\) limits]
\label{sphere-metric-fixed-realization}
There are fixed compatible length metrics
\[
 D_{12}\text{ on }Z_{12},\qquad
 D_{24}^{\mathcal B}\text{ on }Z_{24}^{\mathcal B},\qquad
 D_{24}^{\mathcal C}\text{ on }Z_{24}^{\mathcal C}
\]
with the following properties.

On the inner cone of \(\mathcal A\), the metric is induced by
\begin{equation}
 g_{\mathcal A}
 =dr^2+r^2du^2+\frac{r^2}{T(r)^2}h_{T(r)}.
 \label{sphere-metric-A-fixed-metric}
\end{equation}
On the cusp cone of either \(\mathcal B\) or \(\mathcal C\), it is induced
by
\begin{equation}
 g_{\mathcal D}=dr^2+r^2(du^2+d\theta^2),
 \qquad\mathcal D\in\{\mathcal B,\mathcal C\}.
 \label{sphere-metric-BC-fixed-metric}
\end{equation}
The metrics on the smash powers are the metric quotients of the product
length metrics after the closed fat wedge is collapsed.  Write
\(g_{Z_{12}}\) and \(g_{Z_{24}^{\mathcal D}}\) for the resulting smooth
Riemannian tensors on the open regular strata.

Let
\[
 p_T^{12}:=j_T(0,0),\qquad
 p_T^{24,\mathcal D}:=k_T^{\mathcal D}(0,0,0,0).
\]
Put
\[
 \Sigma_T^{12}:=j_T(Q_T^2),\qquad
 \Sigma_T^{24,\mathcal D}:=k_T^{\mathcal D}(C_T^4).
\]
Let \(d_{T,12}^{\mathrm{slice}}\) and
\(d_{T,24,\mathcal D}^{\mathrm{slice}}\) be the intrinsic Riemannian path
metrics of the first fundamental forms on these embedded slices.  These
must be distinguished from the restricted ambient distance functions
\[
 d_{T,12}^{\mathrm{amb}}
 :=D_{12}|_{\Sigma_T^{12}\times\Sigma_T^{12}},\qquad
 d_{T,24,\mathcal D}^{\mathrm{amb}}
 :=D_{24}^{\mathcal D}|_{\Sigma_T^{24,\mathcal D}
                         \times\Sigma_T^{24,\mathcal D}}.
\]
In general the latter can be smaller, because an ambient path is allowed
to leave the slice.

Then the induced first fundamental forms on the target slices satisfy the
exact identities
\begin{equation}
 \left(\frac{T}{r_T}\right)^2
  j_T^*g_{Z_{12}}
 =g_{Q_T}\oplus g_{Q_T},
 \label{sphere-metric-A-slice-exact}
\end{equation}
and
\begin{equation}
 \left(\frac{\ell_T}{r_T}\right)^2
  (k_T^{\mathcal D})^*g_{Z_{24}^{\mathcal D}}
 =g_{C_T}^{\oplus4}.
 \label{sphere-metric-BC-slice-exact}
\end{equation}
Consequently,
\begin{align}
 \left(\Sigma_T^{12},\frac{T}{r_T}d_{T,12}^{\mathrm{slice}},
       p_T^{12}\right)
 &\longrightarrow(\mathbb R^4,0),\\
 \left(\Sigma_T^{24,\mathcal D},
       \frac{\ell_T}{r_T}d_{T,24,\mathcal D}^{\mathrm{slice}},
       p_T^{24,\mathcal D}\right)
 &\longrightarrow(\mathbb R^4,0)
 \label{sphere-metric-slice-fixed-blowups}
\end{align}
in pointed local-isometric and pointed Gromov--Hausdorff topology, and
smoothly on the exhausting Euclidean slice balls.

The full ambient pointed blow-ups are instead
\begin{align}
 \left(Z_{12},\frac{T}{r_T}D_{12},p_T^{12}\right)
 &\longrightarrow(\mathbb R^8,0),\\
 \left(Z_{24}^{\mathcal D},
       \frac{\ell_T}{r_T}D_{24}^{\mathcal D},
       p_T^{24,\mathcal D}\right)
 &\longrightarrow(\mathbb R^{12},0).
 \label{sphere-metric-ambient-fixed-blowups}
\end{align}
as pointed metric spaces in pointed Gromov--Hausdorff topology.  More
precisely, for every fixed radius the corresponding rescaled target ball is
eventually contained in an open regular-stratum coordinate neighborhood,
and the pulled-back Riemannian tensors converge there in every \(C^k\) norm.
Thus smooth convergence is asserted on an exhaustion by regular-stratum
neighborhoods, not globally on the singular compact targets.  Under these
limits, the slice \(\mathbb R^4\) is respectively the coordinate subspace
obtained by fixing the two radial and two base-angle coordinates, or by
fixing the four radial and four base-angle coordinates.
\end{theorem}

\begin{proof}
Equation \eqref{sphere-metric-normalized-Q-limit} shows that the metrics on
the compact cross-sections of \eqref{sphere-metric-A-fixed-metric} remain
uniformly bounded and converge to a positive definite flat metric.
Therefore the cross-section diameter is \(O(r)\).  Moreover the radial
coordinate is \(1\)-Lipschitz and a radial segment from level \(r\) to the
vertex has length \(r\).  The length completion consequently adds exactly
one point at \(r=0\), and its topology is the prescribed cone topology.
The same assertions are immediate for
\eqref{sphere-metric-BC-fixed-metric}.  Interpolation on a closed annulus
away from the vertex extends these metrics to fixed compatible length
metrics on the compact spaces \(\mathcal A,\mathcal B,\mathcal C\).

If a compact metric space \((E,d)\) and a closed subset \(W\subset E\) are
given, the metric quotient which collapses \(W\) to one point is
\[
 d_{E/W}([x],[y])
 =\min\{d(x,y),d(x,W)+d(y,W)\}.
\]
Applying this to the fat wedges in the product length spaces gives the
claimed fixed compatible metrics on the smash targets.

On the \(r_T\)-slice of \(\mathcal A\), with \(u=0\), the first fundamental
form is
\[
 \frac{r_T^2}{T^2}h_T.
\]
The definition \(h_T=\iota_T^*g_{Q_T}\) proves
\eqref{sphere-metric-A-slice-exact}.  On the \(r_T\)-slice of
\(\mathcal B\) or \(\mathcal C\), with \(u=0\), the form is
\(r_T^2d\theta^2\).  Under
\(\vartheta_T:[\theta]\mapsto[\ell_T\theta]\), the native quotient-circle
metric is \(\ell_T^2d\theta^2\).  This proves
\eqref{sphere-metric-BC-slice-exact}.  The native systole estimates
\[
 \operatorname{sys}(Q_T^2)\geq T-C,
 \qquad
 \operatorname{sys}(C_T^4)=\ell_T
\]
then prove \eqref{sphere-metric-slice-fixed-blowups}: every prescribed
radius ball is eventually exactly a Euclidean four-ball.

It remains to compute the ambient limits.  In one \(\mathcal A\)-factor,
put
\[
 \lambda_T=\frac{T}{r_T},\qquad
 s=\lambda_T(r-r_T),\qquad
 \eta=Tu,\qquad z=Y_Tx
\]
on the universal cover of the torus directions.  On every fixed bounded
set in \((s,\eta,z)\), one has
\[
 \frac r{r_T}=1+\frac sT,
 \qquad
 T(r)=T+O(T^{-1}),
 \qquad
 Y_{T(r)}Y_T^{-1}=I+O(T^{-2}),
\]
with the same estimates after any fixed number of derivatives.  Thus
\(\lambda_T^2g_{\mathcal A}\) converges in these coordinates to
\[
 ds^2+d\eta^2+dz_1^2+dz_2^2.
\]
The inward radial boundary is at distance \(T\), the \(u\)-period is \(T\),
and the \(z\)-period lattice is \(Y_T\mathbb Z^2\), whose systole is at
least \(T-C\).  The outward edge of the fixed cusp chart is also at scaled
distance
\[
 \frac{T}{r_T}(r_0-r_T)
 =T\left(\frac{r_0}{r_T}-1\right)\longrightarrow\infty.
\]
Thus the inward vertex, all periods, and the outward edge recede to
infinity.  A single factor converges on exhausting regular-stratum
coordinate neighborhoods to \(\mathbb R^4\), and two factors converge to
\(\mathbb R^8\).

For one \(\mathcal B\)- or \(\mathcal C\)-factor, put
\[
 \lambda_T=\frac{\ell_T}{r_T},\qquad
 s=\lambda_T(r-r_T),\qquad
 \eta=\ell_Tu,\qquad \zeta=\ell_T\theta.
\]
Then
\[
 \lambda_T^2g_{\mathcal D}
 =ds^2+\left(\frac r{r_T}\right)^2
       (d\eta^2+d\zeta^2),
\]
which converges smoothly on bounded coordinate sets to the Euclidean metric
on \(\mathbb R^3\).  The two circle periods and the inward radial boundary
are all at scaled distance \(\ell_T\to\infty\), while the outward edge is
at scaled distance
\[
 \frac{\ell_T}{r_T}(r_0-r_T)
 =\ell_T\left(\frac{r_0}{r_T}-1\right)\longrightarrow\infty.
\]
Four factors therefore converge on exhausting regular-stratum coordinate
neighborhoods to \(\mathbb R^{12}\).

Finally, a level-\(r_T\) factor point has distance exactly \(r_T\) from its
cone vertex.  Indeed, the radial segment has that length.  A competing
curve contained in the inner chart has length at least \(r_T\), while one
which leaves has inner-chart portions of total length at least
\((r_0-r_T)+r_0>r_T\).  The Riemannian-\(\ell^2\)-product distance from a
synchronized point to the fat wedge is the minimum of the factor
distances, and hence is also \(r_T\).

For the chosen scale put
\[
 b_T=\lambda_Tr_T,\qquad a_T=\lambda_T(r_0-r_T),
\]
and let \(\sigma_T\) be the least scaled period in a factor.  Thus
\(\sigma_T\geq T-C\) in the \(\mathcal A\) calculation and
\(\sigma_T=\ell_T\) in the \(\mathcal B,\mathcal C\) calculation.  The
preceding computations give
\[
 a_T\longrightarrow\infty,\qquad
 b_T\longrightarrow\infty,\qquad
 \sigma_T\longrightarrow\infty.
\]
Fix \(R\).  A curve of rescaled length at most \(R\) from \(p_T\) meets
neither the fat wedge nor \(r=r_0\) for all large \(T\), so the
radius-\(R\) target ball lies in the open regular stratum.  If \(x,y\) lie
in that ball, a competitor through the fat wedge has length at least
\[
 (b_T-R)+(b_T-R)>2R,
\]
and a competitor which exits the inner chart and returns has length at
least
\[
 2(a_T-R)>2R.
\]
Every remaining competitor of length at most \(2R\) stays in
\(r/r_T=1+o_R(1)\).  After lifting it to the angular and torus covers, a
nonzero deck translation has length at least
\((1-o_R(1))(\sigma_T-2R)>2R\) for all large \(T\).  Since the path through
\(p_T\) has length at most \(2R\), none of these three kinds of global
shortcut computes the distance between \(x\) and \(y\).  Thus the metric
on every fixed target ball is computed entirely in one lifted
regular-stratum chart.  The \(C^\infty\) convergence of the pulled-back
tensors on these exhausting charts implies the pointed
Gromov--Hausdorff limits in
\eqref{sphere-metric-ambient-fixed-blowups}.
\end{proof}

\begin{theorem}[Complete intrinsic scalar-rescaling trichotomy]
\label{sphere-metric-scaling-trichotomy}
For the fixed metrics of
\cref{sphere-metric-fixed-realization}, let \(\lambda_T>0\) be arbitrary
and put
\[
 \beta_T:=\lambda_Tr_T.
\]
Choose one of the following triples, retaining the indicated ambient
metric:
\[
 \begin{split}
  (\Sigma_T,d_T^{\mathrm{slice}},d_T^{\mathrm{amb}},p_T;D)
  &=(\Sigma_T^{12},d_{T,12}^{\mathrm{slice}},
      d_{T,12}^{\mathrm{amb}},p_T^{12};D_{12}),\\
  (\Sigma_T,d_T^{\mathrm{slice}},d_T^{\mathrm{amb}},p_T;D)
  &=(\Sigma_T^{24,\mathcal D},
      d_{T,24,\mathcal D}^{\mathrm{slice}},
      d_{T,24,\mathcal D}^{\mathrm{amb}},
      p_T^{24,\mathcal D};D_{24}^{\mathcal D}),
      \qquad \mathcal D\in\{\mathcal B,\mathcal C\}.
 \end{split}
\]
After passing to a subsequence so that \(\beta_T\) has a limit in
\([0,\infty]\), the intrinsic pointed slice limit is
\begin{equation}
 (\Sigma_T,\lambda_Td_T^{\mathrm{slice}},p_T)
 \longrightarrow
 \begin{cases}
  \{\mathrm{pt}\},&\beta_T\longrightarrow0,\\
  ((\mathbb R/L\mathbb Z)^4,0),
      &\beta_T\longrightarrow L\in(0,\infty),\\
  (\mathbb R^4,0),&\beta_T\longrightarrow\infty.
 \end{cases}
 \label{sphere-metric-scaling-trichotomy-result}
\end{equation}
The convergence in
\eqref{sphere-metric-scaling-trichotomy-result} is pointed
Gromov--Hausdorff in all three cases.  In the finite positive case the flat
metrics converge smoothly under the fixed marked-torus identifications; in
the unrolling case the convergence is pointed local-isometric and smooth on
the exhausting Euclidean balls.
In the middle line the lattice is orthogonally equivalent to the square
lattice \(L\mathbb Z^4\).  At the same time,
\begin{equation}
 \lambda_TD(p_T,*)=\beta_T
 \label{sphere-metric-vertex-distance}
\end{equation}
for every one of the displayed choices.

Only in the unrolling regime \(\beta_T\to\infty\) do we compare the
intrinsic metric with the restricted ambient distance.  In that regime, for
every \(R\in(0,\infty)\),
\begin{equation}
 \sup_{x,y\in
  \overline B_{\lambda_Td_T^{\mathrm{slice}}}(p_T,R)}
 \left|
  \lambda_Td_T^{\mathrm{amb}}(x,y)
  -\lambda_Td_T^{\mathrm{slice}}(x,y)
 \right|\longrightarrow0.
 \label{sphere-metric-intrinsic-ambient-agreement}
\end{equation}
Moreover, there are \(\varepsilon_{T,R}\to0\) such that
\begin{equation}
 \overline B_{\lambda_Td_T^{\mathrm{slice}}}(p_T,R)
 \subset
 \overline B_{\lambda_Td_T^{\mathrm{amb}}}(p_T,R)
 \subset
 \overline B_{\lambda_Td_T^{\mathrm{slice}}}
             (p_T,(1+\varepsilon_{T,R})R).
 \label{sphere-metric-intrinsic-ambient-ball-comparison}
\end{equation}
Thus the restricted ambient metric has the same pointed
\(\mathbb R^4\) limit on these fixed slice balls.  No agreement is asserted
in the regimes \(\beta_T\to0\) or \(\beta_T\to L\in(0,\infty)\), where
radial or fat-wedge shortcuts need not disappear.  Every rescaling which
unrolls the intrinsic slice sends the cone/fat-wedge base point to infinity,
so no such limit is a tangent cone based at that point.

In the same regime \(\beta_T\to\infty\), the corresponding full ambient
bubbles are
\begin{align}
 (Z_{12},\lambda_TD_{12},p_T^{12})
 &\longrightarrow(\mathbb R^8,0),\\
 (Z_{24}^{\mathcal D},\lambda_TD_{24}^{\mathcal D},
  p_T^{24,\mathcal D})
 &\longrightarrow(\mathbb R^{12},0),
 \qquad\mathcal D\in\{\mathcal B,\mathcal C\}.
 \label{sphere-metric-arbitrary-unrolling-ambient-bubbles}
\end{align}
These are pointed Gromov--Hausdorff limits of the singular metric targets.
The convergence of tensors is \(C^\infty\) only under coordinate
identifications on exhausting open regular-stratum neighborhoods; no
global smooth convergence of the singular targets is claimed.
\end{theorem}

\begin{proof}
For the \(Q_T^2\) route, the rescaled period matrix in each two-dimensional
factor is
\[
 \frac{\lambda_Tr_T}{T}Y_T
 =\beta_T\left(J+\frac{E_T}{T}\right).
\]
If \(\beta_T\to L\in(0,\infty)\), this lattice converges to
\(LJ\mathbb Z^2\); two factors give the square four-torus in
\eqref{sphere-metric-scaling-trichotomy-result}.  If
\(\beta_T\to0\), the lengths of a lattice basis tend to zero, so the
diameter tends to zero.  If \(\beta_T\to\infty\), the reverse triangle
inequality gives systole
\[
 \frac{\beta_T}{T}(T-C)
 =\beta_T\left(1-\frac CT\right)\longrightarrow\infty,
\]
so every fixed-radius ball is eventually Euclidean.

For a \(C_T^4\) slice, the four periods after rescaling have the common
length
\[
 \lambda_Tr_T=\beta_T,
\]
which gives the same three intrinsic cases directly.

Consider first one factor.  A curve from a level-\(r_T\) point to the cone
vertex which remains in the inner chart has length at least \(r_T\), because
the \(dr^2\)-term makes \(r\) \(1\)-Lipschitz there.  If it leaves the
chart, it has a first crossing of \(r=r_0\) and, before ending at the
vertex, a last re-entry through \(r=r_0\); its two inner-chart portions
have total length at least
\((r_0-r_T)+r_0>r_T\).  The radial segment has length \(r_T\), so the
factor distance to its based vertex is exactly \(r_T\).  The
Riemannian-\(\ell^2\)-product distance from every synchronized slice point
to the fat wedge is the minimum of these factor distances and hence is
also \(r_T\).  The metric-quotient formula in the preceding proof therefore
gives \eqref{sphere-metric-vertex-distance}.

It remains to prove the restricted-ambient comparison, so assume
\(\beta_T\to\infty\) and fix \(R\).  Two points in
\(\overline B_{\lambda_Td_T^{\mathrm{slice}}}(p_T,R)\) have intrinsic
distance at most \(2R/\lambda_T\).  The slice path is an ambient competitor,
so
\[
 d_T^{\mathrm{amb}}(x,y)\leq d_T^{\mathrm{slice}}(x,y).
\]
The quotient shortcut through the fat wedge has rescaled length
\(2\lambda_Tr_T=2\beta_T\), and hence is longer than \(2R\) for all large
\(T\).  A path which leaves the fixed cusp chart must reach \(r=r_0\).
Its outward cost is controlled by
\begin{equation}
 \lambda_T(r_0-r_T)
 =\beta_T\left(\frac{r_0}{r_T}-1\right)\longrightarrow\infty,
 \label{sphere-metric-outward-shortcut-exclusion}
\end{equation}
so no path of rescaled length at most \(2R\) can use the interpolating
annulus or the global part of a factor.

Consider now an ambient competitor of length at most
\(2R/\lambda_T+o(\lambda_T^{-1})\) which stays in the cusp chart.  In every
factor its radial coordinate satisfies
\[
 |r-r_T|\leq \frac{2R+o(1)}{\lambda_T},
 \qquad
 \frac r{r_T}=1+O_R(\beta_T^{-1}).
\]
On this tube, the horizontal tensors in the directions tangent to the slice
are uniformly \((1+\varepsilon_{T,R})\)-bilipschitz to their values at
\(r=r_T\), where \(\varepsilon_{T,R}\to0\).  For
\(\mathcal B\) and \(\mathcal C\) this follows directly from
\((r/r_T)^2=1+O_R(\beta_T^{-1})\).  For \(\mathcal A\), the same conclusion
follows, without changing the fixed coordinates or the matrix \(Y_T\), from
\[
 \begin{aligned}
  T(r)-T&=O_R(\beta_T^{-1}),\\
  \frac{T(r)}T&=1+O_R((\beta_TT)^{-1}),\\
  Y_{T(r)}Y_T^{-1}&=I+O_R((\beta_TT)^{-1}).
 \end{aligned}
\]
Projecting the competitor to the torus or character directions, and
discarding its radial and base-angle components, produces a path between
the same slice endpoints.  The bilipschitz estimate therefore gives
\[
 d_T^{\mathrm{slice}}(x,y)
 \leq(1+\varepsilon_{T,R})d_T^{\mathrm{amb}}(x,y).
\]
Together with the opposite inequality and
\(\lambda_Td_T^{\mathrm{slice}}(x,y)\leq2R\), this proves
\eqref{sphere-metric-intrinsic-ambient-agreement}.  The first inclusion in
\eqref{sphere-metric-intrinsic-ambient-ball-comparison} follows from
\(d_T^{\mathrm{amb}}\leq d_T^{\mathrm{slice}}\).  For a point \(x\) in
the middle ball, apply the same fat-wedge and outward exclusions to a
near-minimizing ambient path from \(p_T\) to \(x\), whose rescaled length is
at most \(R+o(1)\), and project that path to the slice directions.  This
gives
\[
 \lambda_Td_T^{\mathrm{slice}}(p_T,x)
 \leq(1+\varepsilon_{T,R})R
\]
after enlarging \(\varepsilon_{T,R}\) by a quantity tending to zero, and
proves the second inclusion.

For completeness, the full ambient bubbles at an arbitrary unrolling scale
are obtained without changing any coordinate or orientation.  In one
\(\mathcal A\)-factor use
\[
 s=\lambda_T(r-r_T),\qquad
 \eta=\beta_Tu,\qquad
 z=\frac{\beta_T}{T}Y_Tx.
\]
Writing \(A_T(r)=Y_{T(r)}Y_T^{-1}\), one obtains the exact expression
\[
 \lambda_T^2g_{\mathcal A}
 =ds^2+\left(\frac r{r_T}\right)^2d\eta^2
  +\left(\frac r{r_T}\right)^2
   \left(\frac{T}{T(r)}\right)^2
   dz^{\mathsf T}A_T(r)^{\mathsf T}A_T(r)dz.
\]
On every fixed bounded coordinate set,
\[
 \frac r{r_T}=1+O_R(\beta_T^{-1}),\qquad
 T(r)-T=O_R(\beta_T^{-1}),\qquad
 A_T(r)=I+O_R((\beta_TT)^{-1}),
\]
with the same estimates after every fixed number of derivatives.  Hence
the tensor converges smoothly to the Euclidean metric on \(\mathbb R^4\).
The \(\eta\)-period is \(\beta_T\), and the \(z\)-period lattice is
\[
 \frac{\beta_T}{T}Y_T\mathbb Z^2,
 \qquad
 \operatorname{sys}\left(
   \frac{\beta_T}{T}Y_T\mathbb Z^2\right)
 \geq\beta_T\left(1-\frac CT\right)\longrightarrow\infty.
\]
In one \(\mathcal B\)- or \(\mathcal C\)-factor use
\[
 s=\lambda_T(r-r_T),\qquad
 \eta=\beta_Tu,\qquad \zeta=\beta_T\theta.
\]
Then
\[
 \lambda_T^2g_{\mathcal D}
 =ds^2+\left(\frac r{r_T}\right)^2
       (d\eta^2+d\zeta^2),
\]
and both angular periods are \(\beta_T\to\infty\).

In all three cases the inward boundary and the fat wedge are at scaled
distance \(\beta_T\), while a lift starting at radius \(r_T\) incurs the
one-way radial cost in
\eqref{sphere-metric-outward-shortcut-exclusion} before its first outward
crossing of \(r=r_0\).  Fix \(R\).  For large \(T\), all these quantities
and all nonzero period lengths exceed \(4R\).
If \(x,y\) lie in the ambient rescaled radius-\(R\) ball, their distances
from the fat wedge are at least \(\beta_T-R\); hence a quotient path through
the smash base point has length at least \(2(\beta_T-R)>2R\).  Every
remaining competitor lifts to the complement of the fat wedge in the
product.  If such a lift exits the inner product chart, then some factor
has a first crossing of \(r=r_0\); the radial variation before that crossing
is at least \(\lambda_T(r_0-r_T)-R\), and returning to the lift of \(y\)
incurs the same lower bound.  Its length is therefore greater than \(2R\).
Finally, a path of length at most \(2R\) remains in the tube
\(r/r_T=1+O_R(\beta_T^{-1})\); there every nonzero deck translation has
length tending to infinity, so such a path cannot wrap around an angular
or torus period.  The displayed coordinate neighborhoods therefore exhaust
the pointed balls without quotient or global shortcuts.
Their smooth tensor convergence proves the pointed Gromov--Hausdorff
limits in
\eqref{sphere-metric-arbitrary-unrolling-ambient-bubbles}, with smooth
convergence confined to this regular-stratum exhaustion.
\end{proof}

\begin{corollary}[Unit-volume normalization does not give
\(\mathbb R^4\)]
\label{sphere-metric-unit-volume}
The four-volume of \(Q_T^2\) is \((\det Y_T)^2\sim T^4\).  Its unit-volume
length factor is therefore \((\det Y_T)^{-1/2}\sim T^{-1}\), and its
unit-volume limit is a square flat four-torus, not \(\mathbb R^4\).
Likewise, the four-volume of \(C_T^4\) is \(\ell_T^4\), and its unit-volume
normalization is exactly the square unit four-torus.
\end{corollary}

\begin{proposition}[What vertex tangents and compact asymptotic cones can
mean]
\label{sphere-metric-tangent-scope}
The compact target topology alone determines no metric tangent cone.
For an exact conic length metric built from fixed piecewise-Riemannian
metrics on the links, however, the vertex links have topological and
Hausdorff dimensions \(7\) for \(Z_{12}\) and \(11\) for both
\(Z_{24}\) targets.  The resulting vertex cones have topological and
Hausdorff dimensions \(8\) and \(12\), and are not \(\mathbb R^4\).

More precisely, writing \(L_3=T^3\) and \(L_2=T^2\), one has the based
topological models
\[
 Z_{12}\cong\operatorname{Susp}(L_3*L_3),
 \qquad
 Z_{24}^{\mathcal B}
 \cong\operatorname{Susp}(L_2*L_2*L_2*L_2).
\]
At the suspension base point their conic links are respectively
\(L_3*L_3\) and \(L_2^{*4}\).  Also
\[
 (\mathcal C,c_0)
 \cong
 (\mathcal M/\partial\mathcal M,[\partial\mathcal M]),
 \qquad
 (Z_{24}^{\mathcal C},*)
 \cong
 (\mathcal M^4/\partial(\mathcal M^4),
  [\partial(\mathcal M^4)]),
\]
so an exact conic metric at the collapsed boundary has link
\(\partial(\mathcal M^4)\).  At the collapsed point of
\(Z_{24}^{\mathcal C}\), the complete integral local-homology calculation is
\begin{equation}
 H_i\bigl(Z_{24}^{\mathcal C},
          Z_{24}^{\mathcal C}\setminus\{*\};\mathbb Z\bigr)
 \cong
 \begin{cases}
  \mathbb Z^{\binom4{i-1}},&2\leq i\leq5,\\
  0,&i=6,7,\\
  \mathbb Z^{\binom4{i-8}},&8\leq i\leq12,\\
  0,&\text{otherwise}.
 \end{cases}
 \label{sphere-metric-C-smash-local-homology}
\end{equation}
In particular, its local homology is not concentrated in top degree.  For
the first two targets the top local homology occurs in degree \(8\) and
\(12\), respectively; for the third target
\eqref{sphere-metric-C-smash-local-homology} includes a top
\(\mathbb Z\) in degree \(12\).  None has the local homology of
\(\mathbb R^4\).

Finally, the standard asymptotic cone obtained by blowing down any fixed
compact metric target is a point.  Indeed, if \(R_n\to\infty\), then
\[
 \operatorname{diam}(Z,R_n^{-1}D)
 =R_n^{-1}\operatorname{diam}(Z,D)\longrightarrow0.
\]
Thus the limits in \cref{sphere-metric-fixed-realization} are
moving-base-point blow-ups, often called bubbles.  They are neither
vertex-based tangent cones nor asymptotic cones of the compact targets.
\end{proposition}

\begin{proof}
For compact spaces \(V_i\) with closed subspaces \(A_i\), there is a
canonical homeomorphism
\[
 (V_1/A_1)\wedge\cdots\wedge(V_m/A_m)
 \cong
 \frac{V_1\times\cdots\times V_m}
      {\bigcup_iV_1\times\cdots\times A_i\times\cdots\times V_m}.
\]
Take \(V_i=\operatorname{Cone}(L_i)\) and \(A_i=L_i\), its outer boundary.
The product of \(m\) cones is homeomorphic to the cone on the iterated join
\(L_1*\cdots*L_m\), and its displayed union of faces is the outer link.
Collapsing that link gives its suspension.  This proves the first two
models and the link dimensions
\[
 \dim(T^3*T^3)=3+3+1=7,
 \qquad
 \dim((T^2)^{*4})=4\cdot2+3=11.
\]

We next prove the based homeomorphism for \(\mathcal C\), rather than
identifying the two presentations without a map.  Put
\(L=\partial\mathcal M\), write the attached cone as
\[
 \operatorname{Cone}(L)
 =L\times[0,1]/(L\times\{0\}),
\]
with base \(L\times\{1\}\), and choose a collar
\(\kappa:L\times[0,1]\hookrightarrow\mathcal M\) satisfying
\(\kappa(x,0)=x\).  There is a homeomorphism
\[
 \Phi:\mathcal C\setminus\{c_0\}
       \longrightarrow\operatorname{int}\mathcal M
\]
given on the punctured cone, the inner collar, and the remaining core by
\[
 \Phi([x,t])=\kappa(x,t/e)\quad(0<t\leq1),\qquad
 \Phi(\kappa(x,s))=\kappa(x,e^{s-1})\quad(0\leq s\leq1),
\]
\[
 \Phi(q)=q
 \quad\text{for}\quad
 q\in\mathcal M\setminus\kappa(L\times[0,1)).
\]
The first two formulas agree at the glued points
\([x,1]=\kappa(x,0)\), and the second and third agree at
\(\kappa(x,1)\).  Their collar-coordinate images are respectively
\(L\times(0,e^{-1}]\) and \(L\times[e^{-1},1]\); hence the displayed
piecewise map and its piecewise inverse are continuous.

The compact Hausdorff spaces \(\mathcal C\) and
\(\mathcal M/L\) are both one-point compactifications of
\(\operatorname{int}\mathcal M\): for \(\mathcal C\) use \(\Phi\), and for
\(\mathcal M/L\) use the quotient map on
\(\mathcal M\setminus L=\operatorname{int}\mathcal M\).  Uniqueness of the
one-point compactification therefore extends \(\Phi\) to a based
homeomorphism
\begin{equation}
 (\mathcal C,c_0)
 \xrightarrow{\ \cong\ }
 (\mathcal M/\partial\mathcal M,[\partial\mathcal M]).
 \label{sphere-metric-C-boundary-quotient-homeomorphism}
\end{equation}
Only after establishing
\eqref{sphere-metric-C-boundary-quotient-homeomorphism} do we apply the
quotient identity at the start of the proof.  It gives
\[
 \mathcal C^{\wedge4}
 \cong
 \frac{\mathcal M^4}
 {\displaystyle\bigcup_{a=1}^4
  \mathcal M^{a-1}\times\partial\mathcal M
  \times\mathcal M^{4-a}}
 =\mathcal M^4/\partial(\mathcal M^4),
\]
and carries the smash base point to the collapsed boundary.

Put \(N=\mathcal M^4\).  Rounding its product corners preserves the
homeomorphism type of the pair \((N,\partial N)\) and makes
\(\partial N\) an eleven-dimensional closed manifold.  A boundary collar
descends in \(N/\partial N\) to a cone neighborhood of the collapsed
point, so the local link is exactly \(\partial N\) and
\begin{equation}
 H_i(N/\partial N,(N/\partial N)\setminus\{*\};\mathbb Z)
 \cong\widetilde H_{i-1}(\partial N;\mathbb Z).
 \label{sphere-metric-boundary-link-local-shift}
\end{equation}

The integral homology calculation in the proof of
\cref{cusp-cone-S24-route} for the specific compact orientable
three-manifold \(\mathcal M\) gives
\[
 H_j(\mathcal M;\mathbb Z)\cong
 \begin{cases}
  \mathbb Z,&j=0,1,\\
  0,&\text{otherwise}.
 \end{cases}
\]
All groups are free, so the integral Kunneth theorem gives
\[
 H_j(N;\mathbb Z)\cong
 \begin{cases}
  \mathbb Z^{\binom4j},&0\leq j\leq4,\\
  0,&\text{otherwise}.
 \end{cases}
\]
Poincare--Lefschetz duality, with the product orientation retained, gives
\[
 H_j(N,\partial N;\mathbb Z)
 \cong H^{12-j}(N;\mathbb Z)
 \cong
 \begin{cases}
  \mathbb Z^{\binom4{j-8}},&8\leq j\leq12,\\
  0,&\text{otherwise}.
 \end{cases}
\]
The long exact sequence of \((N,\partial N)\) now has disjoint absolute
and relative degree ranges.  It gives
\[
 \widetilde H_j(\partial N;\mathbb Z)\cong
 \begin{cases}
  \mathbb Z^{\binom4j},&1\leq j\leq4,\\
  0,&j=5,6,\\
  \mathbb Z^{\binom4{j-7}},&7\leq j\leq11,\\
  0,&\text{otherwise}.
 \end{cases}
\]
Substitution in
\eqref{sphere-metric-boundary-link-local-shift} proves
\eqref{sphere-metric-C-smash-local-homology}.

An exact metric cone is invariant under every radial blow-up at its vertex,
and its Hausdorff dimension is one plus that of its fixed
piecewise-Riemannian link.  For the first target,
\[
 H_8\bigl(\operatorname{Cone}(T^3*T^3),T^3*T^3;\mathbb Z\bigr)
 \cong\mathbb Z;
\]
the analogous top local group for the \(\mathcal B\) target is
\(\mathbb Z\) in degree \(12\), while the full \(\mathcal C\) calculation
is \eqref{sphere-metric-C-smash-local-homology}.  Euclidean
\(\mathbb R^4\) has local homology \(\mathbb Z\) only in degree \(4\), so
none of these conic vertex models is \(\mathbb R^4\).  The diameter
calculation proves the final assertion.
\end{proof}

\paragraph{Consequences for the three sphere routes.}
The exact positive statement is that one fixed metric on each compact
codomain contains the native \(\mathbb R^4\) decompactification as a
vertical four-plane in a moving-point Euclidean bubble.  The whole bubbles
are \(\mathbb R^8\) for the \(S^{12}\) route and \(\mathbb R^{12}\) for the
\(S^{24}\) routes.  The factor \(\beta_T=\lambda_Tr_T\) proves that this
dimension discrepancy cannot be removed by a different scalar rescaling
of the same fixed cusp-adapted metric: the slice unrolls only when the
smash base point has already gone to infinite scaled distance.

No metric on a source sphere is supplied by these quotient constructions.
If an independent smooth Riemannian metric is placed on \(S^{12}\) or
\(S^{24}\), its ordinary tangent spaces are \(\mathbb R^{12}\) or
\(\mathbb R^{24}\).  Pulling a codomain metric back through
\(q_2^{\wedge2}\) or \(q_1^{\wedge4}\), on the regular strata where the
maps are differentiable, gives a degenerate quadratic form along the
positive-dimensional quotient kernels, not a Riemannian source metric.

The moving four-sphere carriers
\[
 \mathbb S_{12,T}^4=S_T^2\wedge S_T^2\subset S^6,\qquad
 \mathbb S_{24,T}^4=C_T^{\wedge4}\subset S^8
\]
of \cref{cusp-cone-angle-forgotten-sphere-carriers} carry separately
assigned metrics.  Namely, \(S_T^2\) first receives the quotient length
metric obtained from
\((\overline T^2,h_T)\cong(Q_T,g_{Q_T})\) by collapsing
\(K^{(1)}\subset\overline T^2\), equivalently by collapsing
\(\iota_T(K^{(1)})\subset Q_T\).  Then
\(\mathbb S_{12,T}^4\) receives the quotient length metric of the
Riemannian \(\ell^2\)-product \((S_T^2)^2\) after its fat wedge is
collapsed.  Likewise, \(\mathbb S_{24,T}^4\) receives the quotient length
metric of the Riemannian \(\ell^2\)-product \(C_T^4\) after its fat wedge is
collapsed.  These metrics are assigned directly to the moving carriers;
they are unrescaled intrinsic quotient metrics, not restrictions of
\(D_{12}\) or \(D_{24}^{\mathcal B}\), whose domains are the conical
targets, and no fixed ambient metric on \(S^6\) or \(S^8\) is used to
define them.  Their Euclidean balls are
based at the midpoint classes
\[
 \begin{aligned}
  y_*&=[(1/2,1/2)]\in\overline T^2,\qquad z_T=c(y_*),\\
  \zeta_{12,T}&=[z_T,z_T],\\
  \zeta_{24,T}
  &=[(\ell_T/2,\ell_T/2,\ell_T/2,\ell_T/2)].
 \end{aligned}
\]
not at the images of \(p_T^{12}\) or \(p_T^{24,\mathcal B}\), which belong
to the subsequently collapsed loci.  In particular, the trichotomy in
\cref{sphere-metric-scaling-trichotomy} is the trichotomy of the
uncollapsed intrinsic flat-torus slices; it is not a finite-\(\beta_T\)
description of the singular quotient-metric moving four-spheres.  Indeed,
the restriction of any one compatible metric on the compact sphere
\(S^6\) or \(S^8\) has diameter bounded by the diameter of that sphere.  By
contrast, in the separately assigned moving metrics the midpoint distances
to the collapsed loci are at least \((T-C)/2\) and exactly
\(\ell_T/2\), respectively, and hence tend to infinity.  Thus the assigned
unrescaled metrics cannot themselves be restrictions of one fixed ambient
metric.  This does not preclude a separately proved realization after a
scalar blow-up of another fixed quotient metric.

Finally, the rank obstruction for a full-base torus quotient must retain
its actual metric hypothesis: \emph{quotient systole tends to infinity}
means the quotient norm induced by the unrescaled Euclidean fibre metric.
Without that qualification the assertion is false, since an arbitrary
\(T\)-dependent scalar multiplication can make the systole of any fixed
nonzero quotient tend to infinity.


\subsection{The exceptional first-coordinate elliptic quotient}

The native real quotient in \cref{cusp-spec-native-exact-sequence} exists
throughout the radial cusp.  A holomorphic quotient obtained from the first
complex coordinate is a different and much more restrictive construction.
We determine its locus without replacing it by an ambient coefficient-space
path.

For \(z\in\mathfrak h_z\), write
\begin{equation}
 \begin{aligned}
  v_{\gamma,z}&=(6\mu(z),\beta(z))^{\mathsf t},&
  v_{u,z}&=(\tau(z),\mu(z))^{\mathsf t},\\
  v_{w,z}&=(1,0)^{\mathsf t},&
  v_{\delta,z}&=(0,1)^{\mathsf t},
 \end{aligned}
 \label{cusp-spec-elliptic-four-periods}
\end{equation}
and put
\begin{equation}
 \mathcal A_z:=\mathbb C^2/
 \mathbb Z\langle v_{\gamma,z},v_{u,z},v_{w,z},v_{\delta,z}\rangle,
 \qquad
 E_{\tau(z)}:=\mathbb C/(\mathbb Z\tau(z)+\mathbb Z).
 \label{cusp-spec-elliptic-source-target}
\end{equation}
These are the marked fibres of \cref{prop:period-input-audit} and the
elliptic curves selected by the same period \(\tau\).

\begin{proposition}[Exact first-coordinate descent and connected kernel]
\label{cusp-spec-first-coordinate-descent}
The linear map
\(\operatorname{pr}_1:\mathbb C^2\to\mathbb C\) induces a surjective
holomorphic homomorphism
\begin{equation}
 \overline{\operatorname{pr}}_{1,z}:\mathcal A_z\longrightarrow
 E_{\tau(z)}
 \label{cusp-spec-first-coordinate-homomorphism}
\end{equation}
if and only if there is a necessarily unique pair
\((a,b)\in\mathbb Z^2\) such that
\begin{equation}
 6\mu(z)=a\tau(z)+b.
 \label{cusp-spec-first-coordinate-condition}
\end{equation}
On this locus define
\begin{equation}
 \tau_f(z):=a\mu(z)-\beta(z).
 \label{cusp-spec-kernel-period}
\end{equation}
Then
\begin{equation}
 \operatorname{Im}\tau_f(z)
 =-D_{\mathrm{can}}(z)>0,
 \qquad
 D_{\mathrm{can}}(z):=\operatorname{Im}\beta(z)
 -\frac{6(\operatorname{Im}\mu(z))^2}{\operatorname{Im}\tau(z)},
 \label{cusp-spec-kernel-period-imaginary-part}
\end{equation}
and the kernel in \cref{cusp-spec-first-coordinate-homomorphism} is connected
and isomorphic to
\begin{equation}
 E_{\tau_f(z)}=\mathbb C/(\mathbb Z+\mathbb Z\tau_f(z)).
 \label{cusp-spec-connected-kernel-elliptic-curve}
\end{equation}
\end{proposition}

\begin{proof}
The four first coordinates in \cref{cusp-spec-elliptic-four-periods} are
\(6\mu(z),\tau(z),1,0\).  Thus the linear projection descends precisely when
\(6\mu(z)\in\mathbb Z\tau(z)+\mathbb Z\), which is
\cref{cusp-spec-first-coordinate-condition}.  The pair is unique because
\(\operatorname{Im}\tau(z)>0\).  The periods \(v_{u,z}\) and \(v_{w,z}\)
map onto the two generators of the target lattice, so the induced map is
surjective.

Replace the marked lattice basis by the unimodular basis whose first vector
is
\begin{equation}
 v_{\gamma,z}-av_{u,z}-bv_{w,z}
 =(0,\beta(z)-a\mu(z))^{\mathsf t}
 =(0,-\tau_f(z))^{\mathsf t}.
 \label{cusp-spec-vertical-unimodular-period}
\end{equation}
Together with \(v_{\delta,z}=(0,1)^{\mathsf t}\), these are exactly the
vertical periods.  Indeed, after subtracting integral multiples of
\(v_{u,z}\) and \(v_{w,z}\), every point in the kernel has first coordinate
zero; and an integral combination of \(v_{u,z},v_{w,z}\) has first
coordinate zero only when both coefficients vanish.  The kernel is therefore
the connected one-dimensional complex torus in
\cref{cusp-spec-connected-kernel-elliptic-curve}.

Taking imaginary parts in
\cref{cusp-spec-first-coordinate-condition} gives
\(6\operatorname{Im}\mu=a\operatorname{Im}\tau\).  Hence
\[
 \operatorname{Im}\tau_f
 =a\operatorname{Im}\mu-\operatorname{Im}\beta
 =\frac{6(\operatorname{Im}\mu)^2}{\operatorname{Im}\tau}
  -\operatorname{Im}\beta
 =-D_{\mathrm{can}}>0
\]
by \cref{prop:period-input-audit}.
\end{proof}

For \((a,b)\in\mathbb Z^2\), define the marked analytic set
\begin{equation}
 \mathcal D_{a,b}:=
 \{z\in\mathfrak h_z:6\mu(z)-a\tau(z)-b=0\},
 \qquad
 \mathcal D:=\bigcup_{(a,b)\in\mathbb Z^2}\mathcal D_{a,b}.
 \label{cusp-spec-elliptic-descent-locus}
\end{equation}

\begin{theorem}[Monodromy, cusp exclusion, and finiteness of the descent locus]
\label{cusp-spec-elliptic-descent-finiteness}
The transformation laws of \cref{prop:period-input-audit} act on the integer
labels in \cref{cusp-spec-elliptic-descent-locus} by
\begin{equation}
 g_0:(a,b)\longmapsto(a,a+b),
 \qquad
 g_2:(a,b)\longmapsto(-b,6+a).
 \label{cusp-spec-descent-label-monodromy}
\end{equation}
Consequently \(\mathcal D\) is invariant under the triangle group and has a
well-defined image in the completed orbifold base \(B\simeq\mathbb P^1\).
That image is finite and is disjoint from some punctured neighborhood of the
distinguished cusp \(p_0\).
\end{theorem}

\begin{proof}
Suppose \(z\in\mathcal D_{a,b}\).  From
\(\tau\circ g_0=\tau-1\) and \(\mu\circ g_0=\mu\),
\[
 6\mu(g_0z)=a(\tau(g_0z)+1)+b
 =a\tau(g_0z)+(a+b),
\]
so \(g_0z\in\mathcal D_{a,a+b}\).  Similarly,
\(\tau\circ g_2=-1/\tau\) and
\(\mu\circ g_2=1+\mu/\tau\) give
\[
 6\mu(g_2z)=6+a+\frac b{\tau(z)}
 =-b\tau(g_2z)+(6+a),
\]
so \(g_2z\in\mathcal D_{-b,6+a}\).  Both label maps are bijections; and
\(g_0,g_2\) generate the triangle group because
\(g_0=(g_1g_2)^{-1}\).  This proves invariance.

We next exclude the cusp with the integer labels allowed to vary.  If no
punctured cusp neighborhood were disjoint from \(\mathcal D\), there would be
points \(z_j\in\mathcal D_{a_j,b_j}\) in the distinguished cusp component
with \(t_c(z_j)\to0\).  Put
\(T_j=\operatorname{Im}\tau(z_j)\) and
\(m_j=\operatorname{Im}\mu(z_j)\).  The cusp construction in
\cref{prop:period-input-audit} and the formulas preceding
\cref{cusp-spec-real-period-matrix} give
\(T_j\to\infty\) while \(m_j\) stays bounded.  Taking imaginary parts yields
\begin{equation}
 a_j=\frac{6m_j}{T_j}\longrightarrow0.
 \label{cusp-spec-cusp-label-a-limit}
\end{equation}
Because \(a_j\) is integral, \(a_j=0\) for all sufficiently large \(j\).
Then \(6\mu(z_j)=b_j\).  The cusp extension of \(\mu\) is bounded, so the
integers \(b_j\) are bounded; after taking a subsequence, they equal one
integer \(b\).  The holomorphic function \(6\mu-b\) extends across
\(t_c=0\) and has zeros accumulating there.  It therefore vanishes on the
cusp disc and, by the identity theorem, on the connected surface
\(\mathfrak h_z\).

More generally, no identity
\(6\mu=a\tau+b\) with fixed integers \(a,b\) can hold on
\(\mathfrak h_z\).  Applying \(g_0\) to such an identity and subtracting it
from the original one gives \(a=0\).  Then \(\mu=b/6\) is constant, whereas
the \(g_2\)-law would give
\begin{equation}
 \frac b6=1+\frac{b}{6\tau}.
 \label{cusp-spec-constant-mu-contradiction}
\end{equation}
This is impossible for the nonconstant map \(\tau\).  The preceding
conclusion \(6\mu-b\equiv0\) is therefore a contradiction, proving cusp
exclusion.

Choose a closed fundamental polygon for the triangle group and remove the
excluded cusp neighborhood.  Its closure \(K\) is compact in
\(\mathfrak h_z\), including the two elliptic fixed points.  The functions
\(\tau,\mu\) are holomorphic there and
\(\inf_K\operatorname{Im}\tau>0\).  At a point of
\(\mathcal D_{a,b}\cap K\),
\begin{equation}
 a=\frac{6\operatorname{Im}\mu}{\operatorname{Im}\tau},
 \qquad
 b=6\mu-a\tau.
 \label{cusp-spec-compact-label-bounds}
\end{equation}
Thus only finitely many integer pairs \((a,b)\) occur on \(K\).  For each of
them, \(6\mu-a\tau-b\) is a nonzero holomorphic function on a neighborhood of
the compact set \(K\), so its zero set meets \(K\) in finitely many points.
The \((3,4,\infty)\) triangle group has one cusp orbit.  Hence cusp exclusion
holds at every lift of the chosen cusp neighborhood.  The set
\(\mathcal D\) is triangle-group invariant by the already proved label
permutation, and it is disjoint from all those lifted cusp neighborhoods.
Consequently every triangle-group orbit contained in \(\mathcal D\) meets
\(K\).  Hence the image of \(\mathcal D\) in the completed orbifold base is
finite.
\end{proof}

\begin{remark}[Exact scope of the elliptic classification]
\label{cusp-spec-elliptic-descent-scope}
The result classifies only the quotient obtained from the specified first
coordinate and the specified target lattice
\(\mathbb Z\tau+\mathbb Z\).  Its cusp exclusion does not prove that
\(\mathcal A_z\) has no other elliptic quotient or elliptic subtorus, that its
N\'eron--Severi group vanishes, or that it carries no smooth integral
magnetic line bundle.  Those are different typed questions.  In particular,
the first-coordinate elliptic construction cannot supply the gap-closing
cusp sequence, whereas the native scalar and classical gauge soft modes
proved above do not use this exceptional locus.
\end{remark}

\begin{theorem}[Base-changed cover monodromy and exact momentum sectors]
\label{cusp-spec-cover-momentum-theorem}
In the ordered integral basis
\((\widehat\gamma,\widehat u,\widehat w,\widehat\delta)\), put
\begin{equation}
 \mathsf S_M:=\operatorname{diag}(1,1,M,M),
 \qquad
 \Lambda^{(M)}:=\mathsf S_M\mathbb Z^4.
 \label{cusp-spec-cover-index-lattice}
\end{equation}
Then \(P_T\Lambda^{(M)}=\Lambda_{T,M}\), and the cover in
\cref{cusp-spec-degree-cover} is induced by the lattice inclusion
\[
 P_T\mathsf S_M\mathbb Z^4\subset P_T\mathbb Z^4.
\]
Writing \(M_0=I+R\), one has
\begin{equation}
 R^2=0,\qquad
 \ker R=\operatorname{im}R
 =\mathbb Z\langle\widehat w,\widehat\delta\rangle,
 \label{cusp-spec-monodromy-kernel-image}
\end{equation}
and
\begin{equation}
 M_0^M=I+MR,\qquad
 \mathsf S_M^{-1}M_0^M\mathsf S_M=M_0.
 \label{cusp-spec-cover-monodromy-conjugacy}
\end{equation}
Thus, after the cyclic base change \(t_c=v^M\), the degree-\(M^2\)
fibre covers form a regular relative cover over \(v\ne0\).  This particular
statement proves only punctured descent.  The distinct completed local
extension, with the required modified source action and the separate
degree-\(M^3\) and degree-\(M^2\) maps, is proved in
\cref{cusp-cover-extension-theorem}.

Equip \(\widetilde F_{T,M}\) with the pullback Euclidean metric.  Its
Riemannian volume is \(M^2D_T\), so normalized Haar measure is exactly
\[
 \mathrm d\widetilde\nu_{T,M}
 :=(M^2D_T)^{-1}\mathrm d\operatorname{vol}_{\widetilde F_{T,M}}.
\]
Let \(\widetilde\Delta_{T,M}\) be the nonnegative self-adjoint scalar
Laplacian on \(L^2(\widetilde F_{T,M},\mathrm d\widetilde\nu_{T,M})\)
obtained from the closed Dirichlet form
\[
 \widetilde{\mathfrak q}_{T,M}[f]
 :=\int_{\widetilde F_{T,M}}|\mathrm df|^2\,
        \mathrm d\widetilde\nu_{T,M},
 \qquad
 \operatorname{Dom}(\widetilde{\mathfrak q}_{T,M})
 =H^1(\widetilde F_{T,M}),
\]
with operator domain
\(\operatorname{Dom}(\widetilde\Delta_{T,M})
=H^2(\widetilde F_{T,M})\).  Its complete orthonormal character basis is
\begin{equation}
 \widetilde\chi_{T,M,k}([x])
 =
 \exp\!\left(
  2\pi i
  \left\langle P_T^{-\mathsf T}\mathsf S_M^{-1}k,x\right\rangle
 \right),
 \qquad k\in\mathbb Z^4,
 \label{cusp-spec-cover-characters}
\end{equation}
and
\begin{equation}
 \operatorname{Spec}(\widetilde\Delta_{T,M})
 =
 \left\{
  4\pi^2
  \left\lvert P_T^{-\mathsf T}\mathsf S_M^{-1}k\right\rvert^2:
  k\in\mathbb Z^4
 \right\}
 \label{cusp-spec-cover-scalar-spectrum}
\end{equation}
with multiplicity.  For this covering map one has
\begin{equation}
 c_{T,M}^*\chi_{T,\ell}
 =\widetilde\chi_{T,M,\mathsf S_M\ell}.
 \label{cusp-spec-cover-character-pullback}
\end{equation}
Consequently the pullback, or equivalently deck-invariant, sector consists
exactly of indices
\begin{equation}
 k\in\mathsf S_M\mathbb Z^4
 \quad\Longleftrightarrow\quad
 k_3\equiv k_4\equiv0\pmod M.
 \label{cusp-spec-cover-invariant-sector}
\end{equation}

There is a constant \(C_{\mathrm{mom}}>0\), independent of sufficiently
large \(T\) and of \(M\ge1\), such that the Euclidean covering radius of
the cover momentum lattice satisfies
\begin{equation}
 \sup_{p\in\mathbb R^4}
 \operatorname{dist}\!\left(
  p,P_T^{-\mathsf T}\mathsf S_M^{-1}\mathbb Z^4
 \right)
 \le C_{\mathrm{mom}}\left(\frac1T+\frac1M\right).
 \label{cusp-spec-cover-momentum-mesh}
\end{equation}
In particular, if \(T_j,M_j\to\infty\), the cover momentum lattices become
locally dense in all of \(\mathbb R^4\), and every \(E\ge0\) is a limit of
scalar eigenvalues in
\cref{cusp-spec-cover-scalar-spectrum}.
\end{theorem}

\begin{proof}
The equality \(P_T\Lambda^{(M)}=\Lambda_{T,M}\) follows by multiplying the
last two period columns \(e_1,e_3\) by \(M\).  The matrix in
\cref{prop:monodromy-certificate} gives
\[
 R(a,b,c,d)^{\mathsf t}=(0,0,b,-a)^{\mathsf t}.
\]
This proves \cref{cusp-spec-monodromy-kernel-image} and \(R^2=0\).
The binomial theorem then gives \(M_0^M=I+MR\).  For every
\((a,b,c,d)^{\mathsf t}\),
\[
 \mathsf S_M^{-1}M_0^M\mathsf S_M
 (a,b,c,d)^{\mathsf t}
 =(a,b,c+b,d-a)^{\mathsf t}
 =M_0(a,b,c,d)^{\mathsf t},
\]
which proves \cref{cusp-spec-cover-monodromy-conjugacy}.  In particular
\(M_0^M\Lambda^{(M)}=\Lambda^{(M)}\), which is precisely the descent
condition after \(t_c=v^M\).

The dual of \(P_T\mathsf S_M\mathbb Z^4\) is
\[
 (P_T\mathsf S_M)^{-\mathsf T}\mathbb Z^4
 =P_T^{-\mathsf T}\mathsf S_M^{-1}\mathbb Z^4.
\]
The standard Fourier theorem for a flat torus proves
\cref{cusp-spec-cover-characters,cusp-spec-cover-scalar-spectrum}.
Substitution of \(k=\mathsf S_M\ell\) proves
\cref{cusp-spec-cover-character-pullback}.  A character on the cover
descends through \(c_{T,M}\) exactly when it is invariant under the deck
group, and the displayed equality shows that this occurs exactly for
\(k\in\mathsf S_M\mathbb Z^4\), proving
\cref{cusp-spec-cover-invariant-sector}.

It remains to prove the quantitative density claim.  Let
\(b_{j,T,M}\) be the four columns of
\(P_T^{-\mathsf T}\mathsf S_M^{-1}\).  The exact matrix
\cref{cusp-spec-full-inverse-transpose}, the boundedness of
\(\alpha_T,\eta_T,c_T,m_T\), and
\[
 L_T,q_T=O(T),\qquad D_T\asymp T^2
\]
give, after increasing a uniform constant \(C_0\),
\[
 \lvert b_{1,T,M}\rvert+\lvert b_{2,T,M}\rvert
 \le\frac{C_0}{T},
 \qquad
 \lvert b_{3,T,M}\rvert+\lvert b_{4,T,M}\rvert
 \le\frac{C_0}{M}.
\]
Given \(p\in\mathbb R^4\), express it uniquely as
\(p=\sum_jt_jb_{j,T,M}\), and choose \(n_j\in\mathbb Z\) with
\(\lvert t_j-n_j\rvert\le\tfrac12\).  Then
\[
 \left\lvert p-\sum_jn_jb_{j,T,M}\right\rvert
 \le\frac12\sum_j\lvert b_{j,T,M}\rvert
 \le\frac{C_0}{2}\left(\frac1T+\frac1M\right).
\]
This proves \cref{cusp-spec-cover-momentum-mesh}.  For \(E>0\), apply it
to any \(p\) with \(4\pi^2|p|^2=E\); continuity gives eigenvalues tending
to \(E\).  For \(E=0\), use \(k=(1,0,0,0)\), whose momentum norm is
\(O(T^{-1})\).
\end{proof}

\paragraph{Scope and residual data.}
The conclusions above use the explicitly declared Euclidean quotient metric.
The canonical period theorem proves the lattice
family but does not itself specify a Riemannian metric, so no assertion about
a different metric follows without an explicit comparison.

The constants \(H\) and \(B\) determine the first relative
\(T^{-1}\)-correction.  When \(B=0\), they do not by themselves determine
which of the two unit pairs is strictly shortest at the next order.  The
canonical functions do determine it through the exact expression
\(B_T(2L_T-B_T)-35m_T^2\) in
\cref{cusp-spec-unit-mode-comparison}; if both \(B=0\) and \(m=0\), its sign
can depend on higher Taylor coefficients of \(b\) and \(\mu\).  This is an
indeterminacy of the two requested zeroth-order constants, not a missing
proof obligation.
The proved construction permits an additive constant in \(\beta\) and chooses
its imaginary part only by the qualitative requirement that the canonical
determinant be negative.  It records no numerical additive normalization.
Accordingly, the formulas retain the resulting value
\(B=\operatorname{Im}b(0)\) and do not silently impose a sign on it.

Finally, the degree-\(M^2\) covers in
\cref{cusp-spec-degree-cover} are local in the marked radial family, exactly
as stated.  The sublattice
\(\mathbb Z\langle\widehat\gamma,\widehat u,
M\widehat w,M\widehat\delta\rangle\) is not invariant under the full cusp
monodromy when \(M>1\): the proved monodromy sends
\(\widehat\gamma\) to \(\widehat\gamma-\widehat\delta\) and
\(\widehat u\) to \(\widehat u+\widehat w\).  Therefore these particular
degree-\(M^2\) covers do not, without additional construction, descend
around the entire punctured cusp.  Its \(M\)-th power does preserve the
sublattice, because it adds \(-M\widehat\delta\) to
\(\widehat\gamma\) and \(M\widehat w\) to \(\widehat u\).  Equivalently,
\(M_0^M=I+M(M_0-I)\).  Thus, after the degree-\(M\) cyclic base change
\(t_c=v^M\), the cover forms a regular relative degree-\(M^2\) cover over
\(v\ne0\).  The pointed-limit proof itself uses no central-fibre extension.
Nevertheless, \cref{cusp-cover-extension-theorem} proves that a corrected
source action extends this cover through \(v=0\) as a finite morphism of
completed \emph{local} cusp families.  It remains true that no cover over the
entire original base, and no cover of the compact threefold \(X\), is used or
constructed here.
